% Application de la biotechnologie dans la production de biens et services — SVT 1ère TI — Cours signé OnBuch+
% Programme officiel MINESEC. Compiler avec : tectonic NCT13-application-biotechnologie-production-c-ti.tex
\newcommand{\DOCMATIERE}{SVT}
\newcommand{\DOCNIVEAU}{1ère TI}
\newcommand{\DOCMODULE}{Module 4 : La Biotechnologie}
\newcommand{\DOCLECON}{NCT13}
\newcommand{\DOCTITRE}{Application de la biotechnologie dans la production de biens et services}
% OnBuch+ — préambule « cours signés » ENRICHI (compilé avec tectonic / XeLaTeX)
% Système visuel « Pop » d'OnBuch+ : fond crème (jamais blanc), bordures encre
% épaisses, ombres « dures » décalées, titres Archivo Black en capitales.
% Version MATHÉMATIQUES : graphes (pgfplots/tikz), boîtes pédagogiques
% (définition, propriété, méthode, attention, le savais-tu, exercice résolu,
% exercice, TP, bilan), page de garde et signature OnBuch+ ancrée en bas.
%
% Macros attendues dans main.tex AVANT \input{preamble} :
%   \DOCMATIERE  \DOCNIVEAU  \DOCMODULE  \DOCLECON (numéro)  \DOCTITRE
\documentclass[11pt]{article}

\usepackage{fontspec}
\usepackage[french]{babel}
\usepackage[a4paper,margin=2cm,top=3.4cm,bottom=2.8cm,headheight=1.4cm,footskip=1.3cm]{geometry}
\usepackage{amsmath,amssymb,amsfonts}
\usepackage{mathtools} % \xmapsto, \coloneqq... souvent utilisés par les modèles
\usepackage{cancel} % simplifications barrées (bilans de forces)
% \abs{z} (module, valeur absolue) et \norm{u} (norme), utilisés par les modèles
\providecommand{\abs}{}\let\abs\relax\DeclarePairedDelimiter{\abs}{\lvert}{\rvert}
\providecommand{\norm}{}\let\norm\relax\DeclarePairedDelimiter{\norm}{\lVert}{\rVert}
\usepackage[table]{xcolor} % [table] : fournit aussi \rowcolors (alternance de lignes)
\usepackage{pagecolor}
\usepackage{tikz}
\usetikzlibrary{arrows.meta,positioning,calc,shapes.geometric,shapes.misc,shapes.arrows,decorations.pathmorphing,decorations.pathreplacing,decorations.markings,patterns,shadows,mindmap,trees,fit,backgrounds,matrix,intersections,angles,quotes}
\usepackage{pgfplots}
\usepackage[version=4]{mhchem} % équations nucléaires \ce{^{235}_{92}U}
\usepackage[european,siunitx]{circuitikz} % schémas électriques
\pgfplotsset{compat=1.18}
\usepgfplotslibrary{fillbetween,groupplots} % aires entre courbes (calcul intégral)
\usepackage{enumitem}
\usepackage{fancyhdr}
% [most] charge toutes les bibliothèques tcolorbox (listings, minted, raster,
% documentation...) et gonfle inutilement la mémoire TeX sur un document long
% (~300 boîtes) : on ne charge que ce qui est réellement utilisé.
\usepackage[skins,breakable]{tcolorbox}
\usepackage{graphicx}
\usepackage{colortbl}
\usepackage{booktabs}
\usepackage{tabularx}
\usepackage{diagbox} % case d'en-tête diagonale des tableaux à double entrée
% Colonnes extensibles centrée / gauche / droite (C, L, R), souvent utilisées par les modèles
\newcolumntype{C}{>{\centering\arraybackslash}X}
\newcolumntype{L}{>{\raggedright\arraybackslash}X}
\newcolumntype{R}{>{\raggedleft\arraybackslash}X}
\usepackage{array}
\usepackage{multicol}
\usepackage{siunitx}
% parse-numbers=false : accepte aussi \SI{2,5 \times 10^{-3}}{...}
\sisetup{locale=FR,output-decimal-marker={,},group-minimum-digits=4,parse-numbers=false}
\DeclareSIUnit\ohm{\ensuremath{\symup{\Omega}}} % U+2126 absent de la police
\DeclareSIUnit\division{div} % réglages d'oscilloscope (ms/div, V/div)
\DeclareSIUnit\tr{tr} % tours (vitesse de rotation, ex. \SI{1500}{\tr\per\minute})
\DeclareSIUnit\ch{ch} % cheval-vapeur (puissance moteur, unité usuelle hors SI)
\DeclareSIUnit\franccfa{FCFA} % franc CFA (coûts, contexte camerounais)
\DeclareSIUnit\dioptre{\ensuremath{\delta}} % vergence d'une lentille (dioptrie)
\usepackage{qrcode}
% \degree : commande fréquemment utilisée par les modèles pour un angle ou une
% température en dehors de \SI{}{} (non fournie par les paquets chargés).
\providecommand{\degree}{^{\circ}}
% \qed (fin de démonstration), utilisé par les modèles sans amsthm
\providecommand{\qed}{\hfill$\square$}
% \barre : double barre des valeurs interdites dans un tableau de variations (tkz-tab)
\providecommand{\barre}{\ensuremath{\|}}
% environnement proof (amsthm non chargé) utilisé par les modèles
\ifdefined\proof\else\newenvironment{proof}[1][Démonstration]{\par\noindent\textit{#1.}\ }{\qed\par}\fi
\usepackage{lastpage}
\usepackage{hyperref}
\hypersetup{hidelinks}

% ============================================================
% Palette « Pop » OnBuch+ — mêmes hex que la classe P (plus_ui.dart)
% ============================================================
\definecolor{poporange}{HTML}{F59321}
\definecolor{popdark}{HTML}{E07A0C}
\definecolor{popcream}{HTML}{FFF6E8}
\definecolor{popink}{HTML}{1C1714}
\definecolor{poppurple}{HTML}{7A5AE0}
\definecolor{popgreen}{HTML}{2E9E6A}
\definecolor{poppink}{HTML}{E0457B}
\definecolor{popblue}{HTML}{2F7DE1}
\definecolor{popgold}{HTML}{C98A12}
\definecolor{popmuted}{HTML}{8A7F76}
\definecolor{popline}{HTML}{EADFD2}
\definecolor{popgray}{HTML}{8A7F76}
\colorlet{popgrey}{popgray}
% Alias tolérants (le modèle nomme parfois les couleurs différemment)
\colorlet{popwhite}{white}
\colorlet{popblack}{popink}
\colorlet{popred}{poppink}
\colorlet{popyellow}{popgold}
% Teintes claires pour fonds de tableaux / zones
\colorlet{popblueL}{popblue!10!white}
\colorlet{poporangeL}{poporange!14!white}
\colorlet{popgreenL}{popgreen!12!white}
\colorlet{poppurpleL}{poppurple!12!white}
\colorlet{poppinkL}{poppink!12!white}
\colorlet{popgoldL}{popgold!14!white}

\pagecolor{popcream}

% ============================================================
% Typographie
% ============================================================
\defaultfontfeatures{Ligatures=TeX}
\setmainfont{PlusJakartaSans-Medium.ttf}[Path=../../../fonts/,BoldFont=PlusJakartaSans-Bold.ttf]
% Maths en sans-serif (Fira Math) avec chiffres et lettres droites de Plus
% Jakarta, pour que les formules se fondent dans le texte.
\usepackage[math-style=ISO,bold-style=ISO]{unicode-math}
\setmathfont{FiraMath-Regular.otf}[Path=../../../fonts/]
\setmathfont{PlusJakartaSans-Medium.ttf}[Path=../../../fonts/,range={up/num,up/{Latin,latin},\mathup}]
\usepackage{icomma} % virgule décimale sans espace en mode math
% Caractères mathématiques Unicode absents des polices de texte (Plus Jakarta,
% Archivo Black) : rendus par leur équivalent mathématique, en texte comme en
% formule — sinon ils s'affichent en carré vide (ex. « Arithmétique dans ℤ »).
\usepackage{newunicodechar}
\newunicodechar{ℝ}{\ensuremath{\mathbb{R}}}
\newunicodechar{ℤ}{\ensuremath{\mathbb{Z}}}
\newunicodechar{ℂ}{\ensuremath{\mathbb{C}}}
\newunicodechar{ℕ}{\ensuremath{\mathbb{N}}}
\newunicodechar{ℚ}{\ensuremath{\mathbb{Q}}}
\newunicodechar{𝔻}{\ensuremath{\mathbb{D}}}
\newunicodechar{α}{\ensuremath{\alpha}}
\newunicodechar{β}{\ensuremath{\beta}}
\newunicodechar{γ}{\ensuremath{\gamma}}
\newunicodechar{δ}{\ensuremath{\delta}}
\newunicodechar{ε}{\ensuremath{\varepsilon}}
\newunicodechar{θ}{\ensuremath{\theta}}
\newunicodechar{λ}{\ensuremath{\lambda}}
\newunicodechar{μ}{\ensuremath{\mu}}
\newunicodechar{σ}{\ensuremath{\sigma}}
\newunicodechar{τ}{\ensuremath{\tau}}
\newunicodechar{φ}{\ensuremath{\varphi}}
\newunicodechar{ψ}{\ensuremath{\psi}}
\newunicodechar{ω}{\ensuremath{\omega}}
\newunicodechar{Σ}{\ensuremath{\Sigma}}
% majuscules produites par \MakeUppercase dans les titres (α-aminés -> Α)
\newunicodechar{Α}{\ensuremath{\alpha}}
\newunicodechar{Β}{\ensuremath{\beta}}
\newunicodechar{Γ}{\ensuremath{\Gamma}}
\newunicodechar{Δ}{\ensuremath{\Delta}}
\newunicodechar{Ε}{\ensuremath{\varepsilon}}
\newunicodechar{Θ}{\ensuremath{\Theta}}
\newunicodechar{Λ}{\ensuremath{\Lambda}}
\newunicodechar{Μ}{\ensuremath{\mu}}
\newunicodechar{Τ}{\ensuremath{\tau}}
\newunicodechar{Φ}{\ensuremath{\Phi}}
\newunicodechar{Ψ}{\ensuremath{\Psi}}
\newunicodechar{Ω}{\ensuremath{\Omega}}
\newunicodechar{↦}{\ensuremath{\mapsto}}
\newunicodechar{∈}{\ensuremath{\in}}
\newunicodechar{∉}{\ensuremath{\notin}}
\newunicodechar{∧}{\ensuremath{\wedge}}
\newunicodechar{⇔}{\ensuremath{\Leftrightarrow}}
\newunicodechar{⟺}{\ensuremath{\Longleftrightarrow}}
\newunicodechar{⇒}{\ensuremath{\Rightarrow}}
\newunicodechar{⟹}{\ensuremath{\Longrightarrow}}
\newunicodechar{∘}{\ensuremath{\circ}}
\newunicodechar{∪}{\ensuremath{\cup}}
\newunicodechar{∩}{\ensuremath{\cap}}
\newunicodechar{≡}{\ensuremath{\equiv}}
\newunicodechar{⊂}{\ensuremath{\subset}}
\newunicodechar{⊥}{\ensuremath{\perp}}
\newunicodechar{∃}{\ensuremath{\exists}}
\newunicodechar{∀}{\ensuremath{\forall}}
\newunicodechar{∛}{\ensuremath{\sqrt[3]{\,}}}
\newunicodechar{∜}{\ensuremath{\sqrt[4]{\,}}}
\newunicodechar{⃗}{\ensuremath{{}^{\to}}}
\newunicodechar{ⁿ}{\ifmmode^{n}\else\textsuperscript{\ensuremath{n}}\fi}
\newunicodechar{ˣ}{\ifmmode^{x}\else\textsuperscript{\ensuremath{x}}\fi}
\newunicodechar{ᵇ}{\ifmmode^{b}\else\textsuperscript{\ensuremath{b}}\fi}
\newunicodechar{ʳ}{\ifmmode^{r}\else\textsuperscript{\ensuremath{r}}\fi}
\newunicodechar{ᵘ}{\ifmmode^{u}\else\textsuperscript{\ensuremath{u}}\fi}
\newunicodechar{ʸ}{\ifmmode^{y}\else\textsuperscript{\ensuremath{y}}\fi}
\newunicodechar{ᵃ}{\ifmmode^{a}\else\textsuperscript{\ensuremath{a}}\fi}
\newunicodechar{ᵉ}{\ifmmode^{e}\else\textsuperscript{\ensuremath{e}}\fi}
\newunicodechar{ᵏ}{\ifmmode^{k}\else\textsuperscript{\ensuremath{k}}\fi}
\newunicodechar{ᵖ}{\ifmmode^{p}\else\textsuperscript{\ensuremath{p}}\fi}
\newunicodechar{ᴺ}{\ifmmode^{N}\else\textsuperscript{\ensuremath{N}}\fi}
\newunicodechar{ᶜ}{\ifmmode^{c}\else\textsuperscript{\ensuremath{c}}\fi}
\newunicodechar{ʰ}{\ifmmode^{h}\else\textsuperscript{\ensuremath{h}}\fi}
\newunicodechar{ᵠ}{\ifmmode^{q}\else\textsuperscript{\ensuremath{q}}\fi}
\newunicodechar{ₙ}{\ifmmode_{n}\else\textsubscript{\ensuremath{n}}\fi}
\newunicodechar{ₐ}{\ifmmode_{a}\else\textsubscript{\ensuremath{a}}\fi}
\newunicodechar{ᵢ}{\ifmmode_{i}\else\textsubscript{\ensuremath{i}}\fi}
\newunicodechar{ᵣ}{\ifmmode_{r}\else\textsubscript{\ensuremath{r}}\fi}
\newunicodechar{ₖ}{\ifmmode_{k}\else\textsubscript{\ensuremath{k}}\fi}
\newunicodechar{ᵦ}{\ifmmode_{\beta}\else\textsubscript{\ensuremath{\beta}}\fi}
\newunicodechar{ₓ}{\ifmmode_{x}\else\textsubscript{\ensuremath{x}}\fi}
\newfontfamily\popdisplay{ArchivoBlack-Regular.ttf}[Path=../../../fonts/]
\newfontfamily\poplogo{SpaceGrotesk-Bold.ttf}[Path=../../../fonts/]
\newfontfamily\popsemi{PlusJakartaSans-SemiBold.ttf}[Path=../../../fonts/]
\newfontfamily\popxbold{PlusJakartaSans-ExtraBold.ttf}[Path=../../../fonts/]
\newfontfamily\popxboldls{PlusJakartaSans-ExtraBold.ttf}[Path=../../../fonts/,LetterSpace=3.0]

\setlist[enumerate]{leftmargin=1.7em,itemsep=5pt,topsep=4pt}
\setlist[itemize]{leftmargin=1.7em,itemsep=4pt,topsep=4pt,label={\color{poporange}\textbullet}}
\setlength{\parindent}{0pt}
\setlength{\parskip}{5pt}
\renewcommand{\arraystretch}{1.25}

% pgfplots : style Pop par défaut
\pgfplotsset{
  every axis/.append style={axis line style={popink,line width=1pt},tick style={popink},
    grid style={popline,line width=0.5pt},label style={font=\small},tick label style={font=\footnotesize}},
  popcurve/.style={poporange,line width=1.8pt,no markers,smooth},
}
% Même style disponible dans un \draw TikZ simple (hors pgfplots)
\tikzset{popcurve/.style={poporange,line width=1.8pt}}
\tikzset{popgrid/.style={popline,very thin}}
\colorlet{popcurve}{poporange} % le nom du style est parfois utilisé comme couleur

% ============================================================
% Pill / squiggle
% ============================================================
\newcommand{\poppill}[3]{%
  \tikz[baseline=-0.55ex]\node[draw=popink,line width=1.2pt,rounded corners=2.6pt,%
    fill=#2,inner xsep=6.5pt,inner ysep=3pt]{%
    \popxboldls\fontsize{7}{7}\selectfont\color{#3}\MakeUppercase{#1}};%
}
\newcommand{\popsquiggle}[1]{%
  \tikz[baseline=-0.4ex]{
    \draw[#1,line width=2.6pt,line cap=round]
      plot [smooth] coordinates {(0,0.09) (0.34,0.01) (0.68,0.21) (1.02,0.01) (1.36,0.10)};
  }%
}
% Mot-clé surligné (marqueur orange sous le texte)
\newcommand{\cle}[1]{{\popxbold\color{popdark}#1}}

% ============================================================
% En-tête / pied
% ============================================================
\pagestyle{fancy}
\fancyhf{}
\renewcommand{\headrulewidth}{0pt}
\renewcommand{\footrulewidth}{0pt}
\lhead{\raisebox{-0.15ex}{\poplogo\fontsize{13}{13}\selectfont\color{popink}OnBuch\color{poporange}+}}
\rhead{\poppill{\DOCMATIERE\ \textbullet\ \DOCNIVEAU\ \textbullet\ Leçon \DOCLECON}{white}{popink}}
\lfoot{\popxbold\fontsize{7.6}{7.6}\selectfont\color{popmuted}\MakeUppercase{Cours signé OnBuch+}\ \textbullet\ {\popsemi\DOCTITRE}}
\rfoot{\tikz[baseline=-0.6ex]\node[rounded corners=8.5pt,fill=popink,minimum height=17pt,minimum width=17pt,inner xsep=5pt,inner sep=0]{\popxbold\fontsize{8}{8}\selectfont\color{popcream}\thepage\,/\,\pageref*{LastPage}};}

% ============================================================
% Titres
% ============================================================
\newcommand{\chaptitre}[2]{%
  \begin{tcolorbox}[enhanced,colback=poporange,colframe=popink,boxrule=2.6pt,arc=11pt,
      drop shadow=popink,left=15pt,right=15pt,top=12pt,bottom=11pt,before skip=3pt,after skip=18pt]
    \poppill{#2}{popink}{popcream}\par\vspace{9pt}
    {\raggedright\hyphenpenalty=10000\exhyphenpenalty=10000\popdisplay\fontsize{22}{24}\selectfont\color{popcream}\MakeUppercase{#1}\par}
  \end{tcolorbox}
}
% Section numérotée : pastille orange avec le numéro + titre Archivo
\newcounter{popsec}
\newcommand{\coursec}[1]{%
  \stepcounter{popsec}\setcounter{popsub}{0}%
  \par\vspace{16pt}\noindent
  \begin{minipage}{\linewidth}
  \tikz[baseline=-0.7ex]\node[draw=popink,line width=1.6pt,fill=poporange,rounded corners=4pt,minimum size=22pt,inner sep=0]{\popdisplay\fontsize{12}{12}\selectfont\color{popcream}\thepopsec};%
  \hspace{8pt}{\popdisplay\fontsize{15}{17}\selectfont\color{popink}\MakeUppercase{#1}}\par
  \vspace{4pt}\noindent{\color{poporange}\rule{36pt}{3.6pt}}
  \end{minipage}\par\nopagebreak\vspace{8pt}
}
\newcounter{popsub}
\newcommand{\courssub}[1]{%
  \stepcounter{popsub}%
  \par\vspace{9pt}\noindent{\popxbold\fontsize{11.5}{13}\selectfont\color{popink}{\color{poporange}\thepopsec.\thepopsub}\hspace{6pt}#1}\par\nopagebreak\vspace{4pt}
}

% ============================================================
% Boîtes pédagogiques — même recette : fond blanc, bordure épaisse colorée,
% ombre dure encre, badge plein. Titre optionnel [#1].
% ============================================================
\tcbset{popbase/.style={breakable,enhanced,colback=white,boxrule=2.3pt,arc=9pt,
  shadow={3pt}{-3pt}{0pt}{popink},left=14pt,right=14pt,top=11pt,bottom=11pt,before skip=11pt,after skip=15pt,
  fontupper=\popsemi\fontsize{10.6}{14.5}\selectfont\color{popink},
  fontlower=\popsemi\fontsize{10.6}{14.5}\selectfont\color{popink}}}
% Coche dessinée (le glyphe ✓ n'existe pas dans les polices du document)
\newcommand{\popcheck}[1]{\tikz[baseline=-0.5ex]\draw[#1,line width=1.6pt,line cap=round,line join=round](-0.13,0.0)--(-0.04,-0.09)--(0.13,0.1);}
\providecommand{\coche}{\popcheck{popgreen}}
\providecommand{\resultat}[1]{\par\smallskip\noindent\fcolorbox{popgreen}{popgreenL}{\parbox{\dimexpr\linewidth-2\fboxsep-2\fboxrule\relax}{\textbf{Résultat.} #1}}\par\smallskip}
\newcommand{\popboxhead}[3]{\poppill{#1}{#2}{white}\ifx\relax#3\relax\else\hspace{6pt}{\popxbold\fontsize{10.6}{12}\selectfont\color{popink}#3}\fi\par\vspace{7pt}}

\newtcolorbox{aretenir}[1][]{popbase,colframe=popgreen,before upper={\popboxhead{À retenir}{popgreen}{#1}}}
\newtcolorbox{exemplebox}[1][]{popbase,colframe=poporange,before upper={\popboxhead{Exemple}{poporange}{#1}}}
\newtcolorbox{definition}[1][]{popbase,colframe=popblue,before upper={\popboxhead{Définition}{popblue}{#1}}}
\newtcolorbox{propriete}[1][]{popbase,colframe=poppurple,before upper={\popboxhead{Propriété}{poppurple}{#1}}}
\newtcolorbox{methode}[1][]{popbase,colframe=popgold,before upper={\popboxhead{Méthode}{popgold}{#1}}}
\newtcolorbox{attention}[1][]{popbase,colframe=poppink,before upper={\popboxhead{Attention}{poppink}{#1}}}
\newtcolorbox{experience}[1][]{popbase,colframe=popblue,colback=popblueL,before upper={\popboxhead{Expérience}{popblue}{#1}}}
% « Le savais-tu ? » : carte violette pleine (contexte, culture, Cameroun)
\newtcolorbox{savaistu}[1][]{popbase,colback=poppurple,colframe=popink,
  fontupper=\popsemi\fontsize{10.4}{14.2}\selectfont\color{white},
  before upper={\poppill{Le savais-tu\,?}{white}{poppurple}\ifx\relax#1\relax\else\hspace{6pt}{\popxbold\color{white}#1}\fi\par\vspace{7pt}}}
% Exercice résolu : énoncé en haut, \tcblower puis la solution sur fond crème
\newcounter{popexo}\newcounter{popexr}
\newtcolorbox{exoresolu}[1][]{popbase,colframe=poporange,colbacklower=popcream,
  segmentation style={popink,line width=1.2pt,dashed},
  before upper={\stepcounter{popexr}\popboxhead{Exercice résolu \thepopexr}{poporange}{#1}},
  before lower={\poppill{Solution}{popink}{popcream}\par\vspace{6pt}}}
% Exercice (non résolu en place) — la correction est dans la partie Corrigés
\newtcolorbox{exercice}[1][]{popbase,colframe=popink,
  before upper={\stepcounter{popexo}\popboxhead{Exercice \thepopexo}{popink}{#1}}}
% Corrigé
\newtcolorbox{corrige}[1][]{popbase,colframe=popgreen,colback=popgreenL,
  before upper={\popboxhead{Corrigé}{popgreen}{#1}}}
% Objectifs / prérequis / bilan
\newtcolorbox{objectifs}{popbase,colframe=popink,colback=poporangeL,
  before upper={\popboxhead{Objectifs de la leçon}{popink}{}}}
\newtcolorbox{prerequis}{popbase,colframe=popink,colback=popblueL,
  before upper={\popboxhead{Prérequis}{popblue}{}}}
\newtcolorbox{bilan}{popbase,colframe=popink,colback=white,boxrule=2.8pt,
  before upper={\popboxhead{Fiche bilan}{poporange}{L'essentiel à connaître pour l'examen}}}
% Figure encadrée avec légende
\newtcolorbox{popfigure}[1][]{enhanced,breakable=false,colback=white,colframe=popline,boxrule=1.4pt,arc=8pt,
  left=10pt,right=10pt,top=10pt,bottom=8pt,before skip=10pt,after skip=12pt,halign=center,
  lower separated=false,halign lower=center,
  fontlower=\popsemi\fontsize{9}{11.5}\selectfont\color{popmuted},
  #1}
\newcommand{\legende}[1]{\par\vspace{6pt}{\popsemi\fontsize{9}{11.5}\selectfont\color{popmuted}#1}}

% ============================================================
% Signature OnBuch+
% ============================================================
\newcommand{\poplogoinline}[1]{{\poplogo\fontsize{#1}{#1}\selectfont\color{popink}OnBuch\color{poporange}+}}
\newcommand{\signatureonbuch}{%
  \begin{tcolorbox}[enhanced,colback=popink,colframe=popink,boxrule=2.2pt,arc=10pt,
      drop shadow=poporange,left=14pt,right=14pt,top=10pt,bottom=10pt,before skip=0pt,after skip=0pt]
    \tikz[baseline=-0.7ex]\node[circle,fill=popgreen,draw=popcream,line width=1.3pt,minimum size=22pt,inner sep=0]{\popcheck{white}};%
    \hspace{9pt}\begin{minipage}[c]{0.62\linewidth}
      {\popxbold\fontsize{12}{13}\selectfont\color{popcream}Signé\ \textbullet\ {\poplogo OnBuch\color{poporange}+}}\par\vspace{2pt}
      {\raggedright\popsemi\fontsize{8}{10.5}\selectfont\color{popline}\MakeUppercase{Équipe pédagogique OnBuch+}\\\MakeUppercase{Conforme au programme officiel MINESEC}\par}
    \end{minipage}\hfill
    {\poplogo\fontsize{22}{22}\selectfont\color{popcream}OnBuch\color{poporange}+}
  \end{tcolorbox}%
}

% ============================================================
% Page de garde
% #1 titre · #2 matière · #3 niveau · #4 module · #5 n° leçon · #6 durée ·
% #7 description / objectifs (texte ou itemize)
% ============================================================
\newcommand{\pagegarde}[7]{%
  \thispagestyle{empty}
  \newgeometry{a4paper,margin=1.7cm,top=1.6cm,bottom=1.4cm}%
  \begin{tikzpicture}[remember picture,overlay]
    \fill[poporange!18] ([xshift=-3.2cm,yshift=-3.6cm]current page.north east) circle (4.6cm);
    \fill[poppurple!14] ([xshift=2.4cm,yshift=7.6cm]current page.south west) circle (3.8cm);
  \end{tikzpicture}%
  {\poplogo\fontsize{30}{30}\selectfont\color{popink}OnBuch\color{poporange}+}\hfill
  \raisebox{0.6ex}{\poppill{Cours signé \textbullet\ Premium}{popink}{popcream}}\par
  \vspace{1pt}\popsquiggle{poppurple}\par
  \vspace{8pt}
  \begin{tcolorbox}[enhanced,colback=poporange,colframe=popink,boxrule=2.8pt,arc=13pt,
      drop shadow=popink,left=18pt,right=18pt,top=12pt,bottom=13pt,after skip=10pt]
    \poppill{#2 \textbullet\ #3}{popink}{popcream}\hspace{5pt}\poppill{#4}{white}{popink}\par\vspace{8pt}
    {\popdisplay\fontsize{14}{14}\selectfont\color{popink}LEÇON #5}\par\vspace{4pt}
    {\raggedright\hyphenpenalty=10000\exhyphenpenalty=10000\popdisplay\fontsize{25}{27}\selectfont\color{popcream}\MakeUppercase{#1}\par}
    \vspace{8pt}
    {\popxbold\fontsize{9.5}{9.5}\selectfont\color{popink}\MakeUppercase{Durée conseillée : #6}}
  \end{tcolorbox}
  \begin{tcolorbox}[enhanced,colback=white,colframe=popink,boxrule=2.2pt,arc=10pt,
      drop shadow=popink,left=16pt,right=16pt,top=10pt,bottom=10pt,after skip=8pt]
    \poppill{Dans cette leçon}{poppurple}{white}\par\vspace{5pt}
    \popsemi\fontsize{9.3}{12.6}\selectfont\color{popink}#7
  \end{tcolorbox}
  \noindent\begin{minipage}[t]{0.487\linewidth}
    \begin{tcolorbox}[enhanced,colback=popgreen,colframe=popink,boxrule=2.2pt,arc=10pt,
        drop shadow=popink,left=11pt,right=11pt,top=10pt,bottom=10pt,height=3.1cm,valign=center]
      \tcbox[colback=white,colframe=popink,boxrule=1.3pt,arc=5pt,boxsep=0pt,nobeforeafter,
        left=3pt,right=3pt,top=3pt,bottom=3pt,valign=center]{\qrcode[height=1.9cm]{https://play.google.com/store/apps/details?id=com.luvvix.onbuch&hl=fr}}%
      \hspace{8pt}\begin{minipage}[c]{0.5\linewidth}
        {\popdisplay\fontsize{11}{12.5}\selectfont\color{white}\MakeUppercase{Télécharge OnBuch}}\par\vspace{4pt}
        {\raggedright\popsemi\fontsize{8.4}{11}\selectfont\color{white}Gratuit \textbullet\ Google Play\\Tuteur IA, annales, concours, résultats\par}
      \end{minipage}
    \end{tcolorbox}
  \end{minipage}\hfill
  \begin{minipage}[t]{0.487\linewidth}
    \begin{tcolorbox}[enhanced,colback=poppurple,colframe=popink,boxrule=2.2pt,arc=10pt,
        drop shadow=popink,left=11pt,right=11pt,top=10pt,bottom=10pt,height=3.1cm,valign=center]
      \tikz[baseline=-0.7ex]\node[circle,fill=white,draw=popink,line width=1.3pt,minimum size=1.6cm,inner sep=0]{\popdisplay\fontsize{24}{24}\selectfont\color{poppurple}+};%
      \hspace{8pt}\begin{minipage}[c]{0.6\linewidth}
        {\popdisplay\fontsize{11}{12.5}\selectfont\color{white}\MakeUppercase{Passe à OnBuch+}}\par\vspace{4pt}
        {\raggedright\popsemi\fontsize{8.4}{11}\selectfont\color{white}Tous les cours signés, Tuteur IA illimité, fiches PDF — onglet OnBuch+ de l'app\par}
      \end{minipage}
    \end{tcolorbox}
  \end{minipage}
  \vspace{8pt}
  \popcontenu
  \vfill
  \signatureonbuch
  \restoregeometry
  \newpage
}

% Rangée « Ce cours contient » (page de garde)
\newcommand{\poptile}[3]{% #1 couleur #2 gros texte #3 légende
  \begin{tcolorbox}[enhanced,colback=#1,colframe=popink,boxrule=2pt,arc=9pt,shadow={2.5pt}{-2.5pt}{0pt}{popink},
      left=6pt,right=6pt,top=9pt,bottom=9pt,halign=center,valign=center,height=1.75cm,width=\linewidth,nobeforeafter]
    {\popdisplay\fontsize{12.5}{13}\selectfont\color{white}\MakeUppercase{#2}}\par\vspace{3pt}
    {\centering\popsemi\fontsize{7.8}{9.5}\selectfont\color{white}#3\par}
  \end{tcolorbox}}
\newcommand{\popcontenu}{%
  \par\noindent{\popxboldls\fontsize{7.5}{7.5}\selectfont\color{popmuted}CE COURS SIGNÉ CONTIENT}\par\vspace{7pt}
  \noindent\begin{minipage}[t]{0.235\linewidth}\poptile{popblue}{Cours}{complet, illustré, pas à pas}\end{minipage}\hfill
  \begin{minipage}[t]{0.235\linewidth}\poptile{poporange}{Méthodes}{et exercices résolus}\end{minipage}\hfill
  \begin{minipage}[t]{0.235\linewidth}\poptile{poppink}{Exercices}{type examen + corrigés}\end{minipage}\hfill
  \begin{minipage}[t]{0.235\linewidth}\poptile{popgreen}{Fiche bilan}{l'essentiel à retenir}\end{minipage}\par}

% Sommaire
\newtcolorbox{sommaire}{enhanced,colback=white,colframe=popink,boxrule=2.2pt,arc=10pt,
  drop shadow=popink,left=18pt,right=18pt,top=13pt,bottom=13pt,before skip=6pt,after skip=20pt,
  before upper={\poppill{Sommaire}{poppurple}{white}\par\vspace{9pt}\popsemi\fontsize{10.6}{14.5}\selectfont\color{popink}}}

% Signature de fin de cours — poussée tout en bas de la dernière page
\newcommand{\findecours}{%
  \par\vfill\nopagebreak
  \begin{center}\popsquiggle{poporange}\end{center}\vspace{4pt}
  \signatureonbuch
}
\AtBeginDocument{\def\llbracket{\mathopen{[\mkern-3.5mu[}}\def\rrbracket{\mathclose{]\mkern-3.5mu]}}} % intervalles d'entiers (glyphe ⟦ absent de la police)
% Glyphes absents de Fira Math : redéfinis à partir de glyphes disponibles
\AtBeginDocument{%
  \def\setminus{\mathbin{\backslash}}\let\smallsetminus\setminus
  \def\vdots{\mathord{\vbox{\baselineskip3.2pt\lineskiplimit0pt\kern4pt\hbox{.}\hbox{.}\hbox{.}}}}%
  \def\ddots{\mathinner{\mkern1mu\raise6.5pt\hbox{.}\mkern2mu\raise3.6pt\hbox{.}\mkern2mu\raise0.7pt\hbox{.}\mkern1mu}}%
  \def\triangle{\mathord{\scalebox{0.9}{$\symup{\Delta}$}}}%
  \def\nabla{\mathord{\raisebox{\depth}{\scalebox{1}[-1]{$\symup{\Delta}$}}}}%
  \def\checkmark{\popcheck{popdark}}%
  \def\star{\mathbin{\ast}}\def\bigstar{\mathord{\ast}}%
  \def\diamond{\mathbin{\scalebox{0.6}{$\Diamond$}}}%
  \def\bigcup{\mathop{\mathchoice{\raisebox{-0.3em}{\scalebox{1.7}{$\cup$}}}{\scalebox{1.25}{$\cup$}}{\cup}{\cup}}\displaylimits}%
  \def\bigcap{\mathop{\mathchoice{\raisebox{-0.3em}{\scalebox{1.7}{$\cap$}}}{\scalebox{1.25}{$\cap$}}{\cap}{\cap}}\displaylimits}%
  \def\lesseqqgtr{\mathrel{\lessgtr}}\def\bigoplus{\mathop{\oplus}\displaylimits}%
}
\providecommand{\ssi}{si et seulement si} % abréviation fréquente
\AtBeginDocument{\providecommand{\wideparen}{\overparen}} % arc de cercle

%% FIGFIX-BEGIN v5 — correctifs globaux des figures (docs/onbuch-generation/tools/figfix)
\usepackage{adjustbox}\usepackage{etoolbox}\tcbuselibrary{raster}
% toute figure TikZ/pgfplots plus large que la ligne est réduite à la largeur disponible
\BeforeBeginEnvironment{tikzpicture}{\begin{adjustbox}{max width=\linewidth}}
\AfterEndEnvironment{tikzpicture}{\end{adjustbox}}
% légendes pgfplots lisibles par-dessus les courbes
\pgfplotsset{every axis/.append style={legend style={fill=white,fill opacity=0.92,text opacity=1,draw=black!15,inner sep=3pt}}}
% étiquettes d'arêtes (syntaxe edge["..."]) avec fond blanc
\tikzset{every edge quotes/.append style={fill=white,inner sep=1.2pt}}
%% FIGFIX-END
\begin{document}

\pagegarde{Application de la biotechnologie dans la production de biens et services}{SVT}{1ère TI}{Module 4 : La Biotechnologie}{NCT13}{3 h}{Cette leçon présente la biotechnologie, ses définitions et ses grands domaines d'application (agroalimentaire, santé, environnement). Elle distingue les biotechnologies traditionnelles (fermentation) des biotechnologies modernes (sélection assistée, biotechnologies enzymatiques) et ouvre la réflexion sur leurs enjeux économiques et éthiques.
\vspace{4pt}
\begin{multicols}{2}\small
\begin{itemize}
\item À la fin de cette leçon, je sais définir la biotechnologie et citer ses principaux domaines d'application
\item À la fin de cette leçon, je sais décrire le principe de la fermentation et citer des exemples camerounais de produits fermentés
\item À la fin de cette leçon, je sais décrire le principe d'un procédé biotechnologique appliqué à la production de biens ou de services
\item À la fin de cette leçon, je sais comparer une biotechnologie traditionnelle et une biotechnologie moderne selon des critères précis
\item À la fin de cette leçon, je sais expliquer le principe de la sélection assistée et des biotechnologies enzymatiques
\item À la fin de cette leçon, je sais identifier les enjeux économiques d'une application biotechnologique donnée
\end{itemize}\end{multicols}}

\begin{sommaire}
\begin{multicols}{2}
\begin{enumerate}
\item Définition et domaines d'application de la biotechnologie
\item Les biotechnologies traditionnelles : la fermentation
\item Les biotechnologies modernes : sélection assistée et biotechnologies enzymatiques
\item Comparaison entre biotechnologies traditionnelles et modernes
\item Enjeux économiques et éthiques de la biotechnologie
\item Activité d'intégration
\item Méthodes et exercices résolus
\item Exercices
\item Corrigés des exercices
\item Fiche bilan
\end{enumerate}
\end{multicols}
\end{sommaire}

\begin{objectifs}
\begin{itemize}
\item À la fin de cette leçon, je sais définir la biotechnologie et citer ses principaux domaines d'application
\item À la fin de cette leçon, je sais décrire le principe de la fermentation et citer des exemples camerounais de produits fermentés
\item À la fin de cette leçon, je sais décrire le principe d'un procédé biotechnologique appliqué à la production de biens ou de services
\item À la fin de cette leçon, je sais comparer une biotechnologie traditionnelle et une biotechnologie moderne selon des critères précis
\item À la fin de cette leçon, je sais expliquer le principe de la sélection assistée et des biotechnologies enzymatiques
\item À la fin de cette leçon, je sais identifier les enjeux économiques d'une application biotechnologique donnée
\item À la fin de cette leçon, je sais discuter des enjeux éthiques d'une application biotechnologique en argumentant
\item À la fin de cette leçon, je sais construire un tableau comparatif et un schéma de principe d'un procédé biotechnologique
\end{itemize}
\end{objectifs}

\begin{prerequis}
\begin{itemize}
\item La cellule : structure et rôle des organites (Seconde)
\item Les microorganismes : bactéries et levures (Seconde)
\item Les enzymes : nature protéique et rôle de catalyseur biologique
\item Les notions de reproduction et d'hérédité (croisements, sélection)
\item Les besoins alimentaires et les grandes familles d'aliments
\end{itemize}
\end{prerequis}

\newpage

\coursec{Définition et domaines d'application de la biotechnologie}

\textbf{Situation d'accroche.} Amina, élève en Première C à Bafoussam, passe ses vacances chez son grand-père. Le matin, elle observe que le vin de palme récolté la veille fermente tout seul : des bulles se forment, une odeur caractéristique se dégage. Au marché, elle achète un yaourt dont l'étiquette mentionne des « ferments lactiques ». Enfin, en visitant la coopérative agricole du village, elle découvre avec surprise que les jeunes plants de bananiers plantain ne viennent pas d'un champ, mais d'un laboratoire où ils sont produits par culture in vitro. Trois productions, trois univers apparemment très différents. Pourtant, Amina a l'intuition qu'un principe commun les relie.

\textbf{Problématique.} Qu'y a-t-il de commun entre ces trois productions ? Comment les êtres vivants ou leurs constituants peuvent-ils être mis au service de la production de biens et de services utiles à l'Homme ? Et quels avantages ou risques ces techniques présentent-elles pour un pays comme le Cameroun ?

\courssub{Qu'est-ce que la biotechnologie ?}

\textbf{Observation et intuition.} Dans les trois cas observés par Amina, un être vivant ou un de ses constituants effectue un « travail » à la place de l'Homme : des levures transforment le sucre de la sève de palmier en alcool ; des bactéries lactiques transforment le lait en yaourt ; des cellules de bananier, placées dans un milieu de culture, régénèrent des plants entiers. L'Homme ne fabrique pas directement le produit : il \emph{oriente} et \emph{exploite} une activité biologique naturelle. C'est précisément l'idée de la biotechnologie.

\begin{definition}[Biotechnologie]
La \cle{biotechnologie} est l'ensemble des techniques qui utilisent des organismes vivants (microorganismes, cellules végétales ou animales) ou leurs constituants (enzymes, gènes) pour produire des \cle{biens} ou des \cle{services} utiles à l'Homme.
\end{definition}

\textbf{Analyse de la définition : les trois éléments clés.} Toute application biotechnologique repose sur l'association de trois éléments indissociables :

\begin{itemize}
\item un \cle{agent biologique} : l'organisme vivant (levure, bactérie, cellule végétale) ou le constituant cellulaire (enzyme, gène) qui réalise la transformation ;
\item un \cle{procédé technique} : la méthode mise en œuvre par l'Homme pour encadrer l'activité de cet agent (fermentation, culture in vitro, sélection assistée, réaction enzymatique) ;
\item une \cle{finalité de production} : le bien (boisson, aliment, médicament, plant) ou le service (dépollution, diagnostic, traitement des eaux) obtenu.
\end{itemize}

\begin{methode}[Identifier une application biotechnologique]
Pour reconnaître et décrire une application biotechnologique, procède en trois étapes :
\begin{enumerate}
\item Repérer l'\textbf{agent biologique} utilisé (quel organisme ? quelle molécule d'origine vivante ?).
\item Identifier le \textbf{procédé} mis en œuvre (quelle technique encadre l'activité de cet agent ?).
\item Nommer le \textbf{bien ou service} obtenu (à quoi sert le produit final ?).
\end{enumerate}
Si l'un de ces trois éléments manque (par exemple un procédé purement chimique, sans agent vivant), il ne s'agit pas d'une biotechnologie.
\end{methode}

\begin{exemplebox}[Le vin de palme : une biotechnologie du quotidien]
La production de vin de palme au Cameroun illustre parfaitement la méthode : l'\textbf{agent biologique} est constitué de levures naturelles (champignons microscopiques du genre \emph{Saccharomyces}) présentes dans la sève ; le \textbf{procédé} est la fermentation alcoolique, qui transforme les sucres en éthanol et en dioxyde de carbone selon l'équation \ce{C6H12O6 -> 2C2H5OH + 2CO2} ; le \textbf{bien} obtenu est une boisson alcoolisée consommée dans tout le pays.
\end{exemplebox}

\begin{attention}[Biotechnologie n'est pas synonyme de génie génétique]
Erreur fréquente au Probatoire : réduire la biotechnologie à la « manipulation des gènes ». Le \cle{génie génétique} (modification directe de l'ADN) n'est qu'\emph{une} des techniques des biotechnologies \emph{modernes}. La fermentation du vin de palme, vieille de plusieurs millénaires, est une biotechnologie qui ne fait intervenir aucune manipulation de gènes. La biotechnologie est donc un concept beaucoup plus large que le génie génétique.
\end{attention}

\courssub{Les grands domaines d'application de la biotechnologie}

Les biotechnologies interviennent dans de nombreux secteurs. On retient classiquement trois grands domaines, que l'on code souvent par des couleurs : l'agroalimentaire, la santé et l'environnement.

\begin{popfigure}
\begin{center}
\begin{tikzpicture}[font=\small]
\node[draw=popdark, fill=popblueL, rounded corners, thick, minimum width=3.6cm, minimum height=1.1cm, align=center] (bio) at (0,0) {\textbf{BIOTECHNOLOGIE}\\ organismes vivants\\ ou leurs constituants};
\node[draw=popgreen, fill=popgreenL, rounded corners, thick, minimum width=3.4cm, align=center] (agro) at (-6,0) {\textbf{Agroalimentaire}};
\node[draw=poppink, fill=poppinkL, rounded corners, thick, minimum width=3.4cm, align=center] (sante) at (0,2.8) {\textbf{Santé}};
\node[draw=poporange, fill=poporangeL, rounded corners, thick, minimum width=3.4cm, align=center] (env) at (6,0) {\textbf{Environnement}};
\draw[-{Stealth}, thick, popgreen] (bio) -- (agro);
\draw[-{Stealth}, thick, poppink] (bio) -- (sante);
\draw[-{Stealth}, thick, poporange] (bio) -- (env);
\node[align=left, font=\footnotesize, anchor=north] at (-6,-0.7) {$\bullet$ boissons fermentées,\\ yaourts, pain\\ $\bullet$ variétés améliorées\\ (maïs, manioc, plantain)};
\node[align=left, font=\footnotesize, anchor=west] at (-1.9,2.1) {$\bullet$ vaccins, antibiotiques, insuline\\ $\bullet$ diagnostic médical};
\node[align=left, font=\footnotesize, anchor=north] at (6,-0.7) {$\bullet$ traitement des eaux usées\\ $\bullet$ biogaz à partir\\ des déchets organiques};
\end{tikzpicture}
\end{center}
\legende{Les trois grands domaines d'application de la biotechnologie, avec deux exemples chacun. L'agent biologique est au cœur de chaque application.}
\end{popfigure}

\textbf{Le domaine agroalimentaire.} C'est le domaine le plus ancien et le plus familier au Cameroun. On y distingue deux types d'applications :
\begin{itemize}
\item la \emph{fabrication d'aliments et de boissons fermentées} : vin de palme, bière de maïs ou de mil (« bil-bil »), yaourts et fromages (bactéries lactiques), pain (levure \emph{Saccharomyces cerevisiae} qui fait lever la pâte) ;
\item l'\emph{amélioration des variétés végétales} : sélection de variétés de maïs, de manioc ou de bananier plantain plus productives ou résistantes aux maladies, multiplication de plants sains par culture in vitro dans les stations de recherche (IRAD, CARBAP).
\end{itemize}

\textbf{Le domaine de la santé.} Les biotechnologies permettent de produire des substances impossibles à obtenir par synthèse chimique simple :
\begin{itemize}
\item les \cle{vaccins}, préparés à partir d'agents pathogènes tués, atténués ou de leurs fragments ;
\item les \cle{antibiotiques}, comme la pénicilline produite par le champignon \emph{Penicillium} ;
\item l'\cle{insuline} humaine, produite depuis les années 1980 par des bactéries dans lesquelles on a introduit le gène humain de l'insuline, pour traiter le diabète ;
\item le \cle{diagnostic médical} : tests rapides (paludisme, grossesse) utilisant des anticorps, molécules produites par des cellules du système immunitaire.
\end{itemize}

\textbf{Le domaine de l'environnement.} Les microorganismes sont de formidables « ouvriers de la dépollution » :
\begin{itemize}
\item le \cle{traitement des eaux usées} : dans les stations d'épuration, des bactéries dégradent la matière organique polluante (boues activées) ;
\item la \cle{dépollution des sols} (bioremédiation) : certains microorganismes dégradent les hydrocarbures après une pollution pétrolière ;
\item la production de \cle{biogaz} : la fermentation des déchets organiques (fumier, résidus de cuisine) par des bactéries en l'absence d'oxygène produit du méthane, utilisable pour la cuisson ou l'éclairage — une piste précieuse pour les zones rurales camerounaises.
\end{itemize}

\courssub{Biotechnologies traditionnelles et biotechnologies modernes}

\textbf{Démarche d'investigation.} Comparons deux productions : le vin de palme du grand-père d'Amina et l'insuline produite en usine.
\begin{itemize}
\item \textbf{Observation} : le vin de palme fermente spontanément, sans aucune connaissance scientifique préalable ; l'insuline exige un laboratoire, la connaissance de l'ADN et des techniques de pointe.
\item \textbf{Hypothèse} : il existe deux grandes générations de biotechnologies, distinguées par le niveau de connaissances scientifiques mobilisé.
\item \textbf{Données historiques} : la fermentation est pratiquée depuis plus de 6000 ans (bière en Mésopotamie, pain en Égypte antique) ; la nature microbienne de la fermentation n'est démontrée qu'en 1857 par Louis Pasteur ; le mot « biotechnologie » n'apparaît qu'en 1919 ; la production d'insuline par génie génétique date de 1982.
\item \textbf{Interprétation} : pendant des millénaires, l'Homme a utilisé les êtres vivants de façon \emph{empirique}, sans comprendre les mécanismes ; depuis le XX\textsuperscript{e} siècle, il les utilise de façon \emph{raisonnée}, fondée sur la microbiologie, la biochimie et la génétique.
\item \textbf{Conclusion} : on distingue les biotechnologies \emph{traditionnelles} et les biotechnologies \emph{modernes}.
\end{itemize}

\begin{definition}[Traditionnelle ou moderne ?]
Les \cle{biotechnologies traditionnelles} sont des pratiques anciennes et empiriques, fondées sur l'expérience transmise de génération en génération (fermentations artisanales, sélection paysanne des semences). Les \cle{biotechnologies modernes} sont fondées sur les connaissances scientifiques récentes (microbiologie, biologie moléculaire) : sélection assistée, culture in vitro, biotechnologies enzymatiques, génie génétique.
\end{definition}

\begin{popfigure}
\begin{center}
\begin{tikzpicture}[font=\footnotesize]
\draw[-{Stealth}, very thick, popdark] (0,0) -- (13.5,0) node[below left=2pt and 0pt]{temps};
\foreach \x/\date/\txt/\col in {
  1/{-4000}/{Fermentation ancestrale\\ (bière, pain, vin)}/popgreen,
  4/{1857}/{Pasteur : rôle des\\ microorganismes}/popblue,
  6.5/{1919}/{Mot « biotechnologie »\\ (K. Ereky)}/poppurple,
  9/{années 1950}/{Sélection assistée des\\ variétés végétales}/poporange,
  11.5/{1982}/{Insuline par génie génétique\\ (biotechnologies modernes)}/poppink}
{
  \fill[\col] (\x,0) circle (2.5pt);
  \draw[\col, thick] (\x,0) -- (\x,0.9);
  \node[above, align=center, \col] at (\x,0.9) {\textbf{\date}\\ \txt};
}
\node[below, align=center, popgreen] at (2.5,-0.35) {\textbf{Biotechnologies traditionnelles}};
\draw[popgreen, thick] (0.3,-0.25) -- (5.2,-0.25);
\node[below, align=center, poppink] at (10.2,-0.35) {\textbf{Biotechnologies modernes}};
\draw[poppink, thick] (5.4,-0.25) -- (13.2,-0.25);
\end{tikzpicture}
\end{center}
\legende{Frise chronologique : de la fermentation ancestrale aux biotechnologies modernes. La frontière se situe avec l'essor de la microbiologie et de la biologie moléculaire.}
\end{popfigure}

\begin{savaistu}[Un mot récent pour une pratique millénaire]
Le mot « biotechnologie » n'a été forgé qu'en 1919 par l'ingénieur hongrois Károly Ereky. Pourtant, l'Homme pratique la fermentation depuis plus de 6000 ans : les Sumériens brassaient déjà la bière, et les Égyptiens de l'Antiquité fabriquaient du pain levé. Au Cameroun, la fermentation du vin de palme, du « sha'a » (maïs fermenté) ou du manioc pour le bâton de manioc s'inscrit dans cette longue tradition, bien antérieure à la science qui l'explique aujourd'hui.
\end{savaistu}

\begin{exoresolu}[Reconnaître une application biotechnologique]
Énoncé. Pour chacune des productions suivantes, indique s'il s'agit d'une application biotechnologique ; si oui, précise l'agent biologique, le procédé et le bien obtenu, ainsi que le domaine concerné : (a) fabrication de yaourt au marché de Bafoussam ; (b) production d'engrais chimique par synthèse industrielle ; (c) production de biogaz à partir de fumier.
\tcblower
Solution détaillée.
\begin{itemize}
\item (a) \textbf{Oui}. Agent biologique : bactéries lactiques (\emph{Lactobacillus}, \emph{Streptococcus thermophilus}). Procédé : fermentation lactique du lait. Bien obtenu : le yaourt. Domaine : agroalimentaire.
\item (b) \textbf{Non}. La synthèse d'engrais est un procédé purement chimique : aucun organisme vivant ni constituant du vivant n'intervient. Ce n'est donc pas une biotechnologie.
\item (c) \textbf{Oui}. Agent biologique : bactéries méthanogènes. Procédé : fermentation anaérobie (méthanisation) des déchets organiques. Bien obtenu : le biogaz (méthane), source d'énergie. Domaine : environnement.
\end{itemize}
\end{exoresolu}

\begin{aretenir}[L'essentiel de la section]
\begin{itemize}
\item La biotechnologie utilise des organismes vivants ou leurs constituants (enzymes, gènes) pour produire des biens et des services.
\item Toute application biotechnologique associe un \textbf{agent biologique}, un \textbf{procédé} et une \textbf{finalité de production}.
\item Trois grands domaines : agroalimentaire (fermentations, variétés améliorées), santé (vaccins, antibiotiques, insuline, diagnostic), environnement (eaux usées, dépollution, biogaz).
\item On distingue les biotechnologies \textbf{traditionnelles} (empiriques, anciennes) et \textbf{modernes} (fondées sur la science récente) ; le génie génétique n'est qu'une technique parmi d'autres des biotechnologies modernes.
\end{itemize}
\end{aretenir}

\coursec{Les biotechnologies traditionnelles : la fermentation}

Dans la section précédente, nous avons défini la \cle{biotechnologie} comme l'ensemble des techniques qui utilisent des organismes vivants (ou leurs constituants : cellules, enzymes) pour produire des biens et des services utiles à l'être humain. Nous avons aussi distingué deux grandes catégories : les \cle{biotechnologies traditionnelles}, pratiquées empiriquement depuis des millénaires, et les \cle{biotechnologies modernes}, nées des progrès récents de la biologie. Il est temps d'entrer dans le vif du sujet avec la plus ancienne et la plus répandue des biotechnologies traditionnelles : la \cle{fermentation}.

\courssub{Une situation concrète pour commencer : le secret du vin de palme}

\textbf{Observation.} À Batibo, à Bafou ou à Ebolowa, le palmiculteur récolte le matin une sève de palmier limpide et sucrée. Le soir même, cette sève est devenue trouble, mousseuse, légèrement piquante : c'est du \emph{vin de palme}. Le lendemain, il est encore plus fort en alcool. Personne n'a ajouté d'alcool ni chauffé le liquide. Que s'est-il passé ?

\textbf{Hypothèse.} Des êtres vivants microscopiques, présents dans l'air ou dans le récipient de récolte, transforment le sucre de la sève en alcool.

\textbf{Vérification expérimentale.} On observe une goutte de vin de palme frais au microscope optique : on y découvre des cellules ovoïdes, certaines en train de bourgeonner : ce sont des \cle{levures} (champignons unicellulaires du genre \emph{Saccharomyces}). Si l'on fait bouillir la sève pour tuer ces microorganismes et qu'on la conserve ensuite dans un récipient stérile fermé, elle ne fermente pas : elle reste sucrée. En revanche, si l'on ajoute volontairement des levures à la sève stérilisée puis refroidie, la fermentation reprend, avec dégagement de bulles de gaz.

\textbf{Interprétation et conclusion.} Ce sont bien les levures vivantes qui transforment le sucre de la sève en alcool, en dégageant un gaz. Cette transformation biologique, réalisée par des microorganismes, est une fermentation. L'être humain n'a rien inventé : il a appris, au fil des générations, à exploiter et à contrôler un phénomène naturel.

\courssub{Définition de la fermentation}

\begin{definition}[La fermentation]
La \cle{fermentation} est une \textbf{transformation de matière organique} (le plus souvent des \textbf{sucres}) réalisée par des \textbf{microorganismes vivants} (levures ou bactéries), qui en tirent de l'\textbf{énergie} en l'\textbf{absence d'oxygène} (anaérobiose) ou en présence très limitée d'oxygène. Elle produit des substances nouvelles (alcool, acides, gaz) et une faible quantité d'énergie utilisée par le microorganisme pour vivre et se multiplier.
\end{definition}

Trois idées sont essentielles dans cette définition :

\begin{itemize}
\item La fermentation est une réaction \textbf{biologique} : elle est catalysée par des \cle{enzymes} produites par des cellules vivantes. Sans microorganismes vivants, pas de fermentation.
\item Elle se déroule \textbf{sans oxygène} (ou avec très peu) : c'est une voie de production d'énergie dite \emph{anaérobie}, beaucoup moins rentable en énergie que la respiration cellulaire, ce qui explique que le microorganisme doive transformer de grandes quantités de sucre.
\item Elle produit des \textbf{composés nouveaux} (éthanol, acide lactique, dioxyde de carbone...) qui modifient profondément le goût, la texture et la conservation de l'aliment.
\end{itemize}

\begin{attention}[Fermentation n'est pas pourriture]
On entend parfois : « ce jus a tourné, il est pourri ». Erreur fréquente ! La \textbf{pourriture} (décomposition) est une dégradation \emph{incontrôlée} de la matière organique par des microorganismes variés, souvent en présence d'oxygène, qui produit des substances toxiques ou nauséabondes. La \textbf{fermentation} est au contraire une transformation \emph{contrôlée} par des microorganismes \emph{utiles} (levures, bactéries lactiques), dans des conditions maîtrisées, qui donne des produits comestibles, sains et souvent meilleurs que l'aliment de départ. Au Probatoire, ne confonds jamais ces deux termes.
\end{attention}

\courssub{La fermentation alcoolique : le sucre devient alcool}

\textbf{Principe.} Les \cle{levures} (notamment \emph{Saccharomyces cerevisiae}), en anaérobiose, dégradent le glucose pour en extraire de l'énergie. Les déchets de cette dégradation sont l'\textbf{éthanol} (alcool éthylique) et le \textbf{dioxyde de carbone}.

\begin{propriete}[Équation globale de la fermentation alcoolique (admise)]
La fermentation alcoolique du glucose par les levures suit l'équation globale :
\[ \ce{C6H12O6 -> 2 C2H5OH + 2 CO2} + \text{énergie} \]
Cette équation est \textbf{admise} : la démonstration biochimique détaillée (glycolyse, étapes de fermentation) n'est \textbf{pas exigible} au Probatoire en Première C et TI. Retiens simplement : \emph{une molécule de glucose donne deux molécules d'éthanol et deux molécules de dioxyde de carbone}.
\end{propriete}

\begin{popfigure}
\begin{center}
\begin{tikzpicture}[font=\small, >=Stealth]
% Réactifs
\draw[fill=popblueL, rounded corners=4pt] (-5.2,0.6) rectangle (-1.6,2.2);
\node[align=center] at (-3.4,1.4) {\textbf{Substrat}\\ glucose \ce{C6H12O6}\\ (sève, jus, moût)};
\draw[fill=poppinkL, rounded corners=4pt] (-5.2,-2.2) rectangle (-1.6,-0.6);
\node[align=center] at (-3.4,-1.4) {\textbf{Microorganisme}\\ levures vivantes\\ (\emph{Saccharomyces})};
% Cuve
\draw[fill=popcream, thick, rounded corners=6pt] (-0.8,-2.4) rectangle (2.8,2.4);
\node[align=center] at (1,0.9) {\textbf{Cuve de fermentation}};
\node[align=center, text=popdark] at (1,-0.4) {anaérobiose\\ (peu ou pas d'\ce{O2})};
\node[align=center, text=popdark] at (1,-1.5) {température\\ \SIrange{25}{30}{\celsius}};
% Flèches entrées
\draw[->, very thick, popblue] (-1.6,1.4) -- (-0.8,1.4);
\draw[->, very thick, poppink] (-1.6,-1.4) -- (-0.8,-1.4);
% Produits
\draw[fill=popgreenL, rounded corners=4pt] (4.6,0.6) rectangle (8.2,2.2);
\node[align=center] at (6.4,1.4) {\textbf{Éthanol}\\ $2\,\ce{C2H5OH}$\\ (boisson alcoolisée)};
\draw[fill=popgoldL, rounded corners=4pt] (4.6,-2.2) rectangle (8.2,-0.6);
\node[align=center] at (6.4,-1.4) {\textbf{Dioxyde de carbone}\\ $2\,\ce{CO2}$ (bulles)\\ + énergie pour la levure};
% Flèches sorties
\draw[->, very thick, popgreen] (2.8,1.4) -- (4.6,1.4);
\draw[->, very thick, popgold] (2.8,-1.4) -- (4.6,-1.4);
\end{tikzpicture}
\end{center}
\legende{Principe de la fermentation alcoolique : les levures transforment le glucose en éthanol et en dioxyde de carbone, en anaérobiose et à température favorable. L'énergie libérée sert à la survie et à la multiplication des levures.}
\end{popfigure}

\textbf{Évolution au cours du temps.} Au début de la fermentation du vin de palme, les levures se multiplient et consomment le sucre : la concentration en alcool augmente d'abord lentement (phase de latence), puis rapidement (phase de fermentation active), avant de plafonner. Pourquoi ce plafond ? Parce que le sucre s'épuise et surtout parce que l'alcool accumulé devient \emph{toxique pour les levures elles-mêmes} : au-delà d'environ 12 à 15~\% d'alcool en volume, les levures meurent et la fermentation s'arrête. C'est pourquoi les boissons simplement fermentées (vin, bière, vin de palme) ne dépassent jamais naturellement ce taux ; les alcools plus forts exigent une distillation.

\begin{popfigure}
\begin{center}
\begin{tikzpicture}
\begin{axis}[
  width=0.85\linewidth, height=6cm,
  xlabel={Temps de fermentation (heures)},
  ylabel={Teneur en alcool (\% en volume)},
  xmin=0, xmax=48, ymin=0, ymax=8,
  xtick={0,6,12,18,24,30,36,42,48},
  ytick={0,1,2,3,4,5,6,7,8},
  grid=both, grid style={popline!40},
  legend style={at={(0.03,0.97)}, anchor=north west, font=\small}
]
\addplot[popcurve, domain=0:48, samples=200] {6.5/(1+exp(-0.22*(x-18)))};
\addlegendentry{Vin de palme (courbe schématique)}
% Phases
\draw[dashed, popmuted] (axis cs:12,0) -- (axis cs:12,8);
\draw[dashed, popmuted] (axis cs:30,0) -- (axis cs:30,8);
\node[font=\footnotesize, text=popdark] at (axis cs:6,7.4) {latence};
\node[font=\footnotesize, text=popdark] at (axis cs:21,7.4) {fermentation active};
\node[font=\footnotesize, text=popdark] at (axis cs:39,7.4) {ralentissement};
\end{axis}
\end{tikzpicture}
\end{center}
\legende{Évolution schématique de la teneur en alcool au cours de la fermentation du vin de palme : phase de latence (multiplication des levures), phase active (production rapide d'éthanol), puis ralentissement et plafonnement dus à l'épuisement du sucre et à la toxicité de l'alcool pour les levures.}
\end{popfigure}

\begin{exemplebox}[Un exemple chiffré pour bien raisonner]
D'après l'équation $\ce{C6H12O6 -> 2 C2H5OH + 2 CO2}$, une mole de glucose (masse molaire $M = \SI{180}{g/mol}$) donne théoriquement deux moles d'éthanol ($M = \SI{46}{g/mol}$) et deux moles de dioxyde de carbone ($M = \SI{44}{g/mol}$). Ainsi, à partir de \textbf{\SI{180}{g} de glucose}, la fermentation alcoolique \emph{théorique} produit environ \textbf{\SI{92}{g} d'éthanol} ($2 \times 46$) et \textbf{\SI{88}{g} de dioxyde de carbone} ($2 \times 44$). En pratique, le rendement réel est inférieur, car une partie du sucre sert à la croissance des levures.
\end{exemplebox}

\begin{exoresolu}[Calculer une production théorique d'éthanol]
Énoncé : un producteur de boisson fermentée dispose d'un moût contenant \SI{900}{g} de glucose. On admet que la fermentation alcoolique suit l'équation $\ce{C6H12O6 -> 2 C2H5OH + 2 CO2}$ et qu'elle est totale. Données : $M(\ce{C6H12O6}) = \SI{180}{g/mol}$ ; $M(\ce{C2H5OH}) = \SI{46}{g/mol}$ ; $M(\ce{CO2}) = \SI{44}{g/mol}$. Calculer la masse théorique d'éthanol et de dioxyde de carbone produites.
\tcblower
Solution détaillée :
\begin{enumerate}
\item Quantité de glucose : $n = \dfrac{m}{M} = \dfrac{900}{180} = \SI{5}{mol}$.
\item D'après l'équation, \SI{1}{mol} de glucose donne \SI{2}{mol} d'éthanol et \SI{2}{mol} de \ce{CO2}. Donc $n(\text{éthanol}) = n(\ce{CO2}) = 2 \times 5 = \SI{10}{mol}$.
\item Masse d'éthanol : $m = 10 \times 46 = \SI{460}{g}$.
\item Masse de dioxyde de carbone : $m = 10 \times 44 = \SI{440}{g}$.
\end{enumerate}
Vérification : $460 + 440 = 900$~g, masse du glucose consommé : la conservation de la masse est respectée. Conclusion : la fermentation de \SI{900}{g} de glucose produit théoriquement \textbf{\SI{460}{g} d'éthanol} et \textbf{\SI{440}{g} de \ce{CO2}}.
\end{exoresolu}

\courssub{La fermentation lactique : le sucre devient acide lactique}

\textbf{Observation.} Laisse du lait frais dans un récipient propre à température ambiante pendant un jour ou deux (comme le font les éleveurs peuls du Nord-Cameroun pour préparer le \emph{kossam}) : il caille, épaissit et prend un goût agréablement acidulé. Ce n'est pas une décomposition : c'est la \cle{fermentation lactique}.

\textbf{Principe.} Des \cle{bactéries lactiques} (genres \emph{Lactobacillus}, \emph{Streptococcus}, \emph{Lactococcus}) transforment le \textbf{lactose} (sucre du lait) en \textbf{acide lactique}. L'accumulation d'acide lactique abaisse le pH du milieu : les protéines du lait coagulent (le lait « caille ») et la plupart des microorganismes indésirables, qui ne supportent pas l'acidité, sont bloqués.

\[ \text{lactose} \xrightarrow{\text{bactéries lactiques}} \text{acide lactique} + \text{énergie} \]

Contrairement à la fermentation alcoolique, la fermentation lactique ne dégage \textbf{pas de gaz} : c'est un critère simple pour distinguer les deux.

\textbf{Applications concrètes :}
\begin{itemize}
\item les \textbf{yaourts} et laits fermentés industriels (ensemencement contrôlé du lait par des souches sélectionnées de \emph{Lactobacillus bulgaricus} et \emph{Streptococcus thermophilus}) ;
\item le \textbf{lait caillé} traditionnel et le \emph{kossam} (lait fermenté consommé dans le Grand Nord camerounais, souvent mélangé à de la bouillie de mil) ;
\item les \textbf{légumes lactofermentés} (choucroute, olives, cornichons) : le sel ajouté sélectionne les bactéries lactiques au détriment des microorganismes de décomposition.
\end{itemize}

\courssub{Les fermentations dans l'alimentation camerounaise}

La fermentation est au cœur de nombreuses traditions alimentaires camerounaises. Quelques exemples à connaître :

\begin{center}
\begin{tabularx}{\linewidth}{|l|X|X|}
\hline
\rowcolor{popblueL} \textbf{Produit} & \textbf{Substrat et microorganismes} & \textbf{Type de fermentation} \\
\hline
Vin de palme & Sève de palmier (sucres) ; levures naturelles & Alcoolique \\
\hline
Bière de maïs « bili-bili » & Maïs ou mil germé (amidon transformé en sucres) ; levures & Alcoolique \\
\hline
Bâton de manioc « bobolo », « miondo » & Tubercules de manioc ; bactéries lactiques et levures & Mixte (lactique et alcoolique) \\
\hline
« Kossam », lait caillé & Lait (lactose) ; bactéries lactiques & Lactique \\
\hline
\end{tabularx}
\end{center}

\begin{savaistu}[La fermentation rend le manioc moins toxique]
Le manioc \emph{amer}, très cultivé au Cameroun, contient des composés cyanogènes qui peuvent libérer du \textbf{cyanure}, substance toxique. Le trempage prolongé et la fermentation des tubercules (étapes de la fabrication du \emph{bobolo} et du \emph{miondo}) permettent aux enzymes des microorganismes de dégrader une grande partie de ces composés toxiques. La fermentation traditionnelle est donc aussi une technique de \textbf{détoxification} alimentaire, mise au point empiriquement bien avant que la chimie n'en explique le mécanisme !
\end{savaistu}

\courssub{À quoi sert la fermentation ? Les quatre grands rôles}

\begin{aretenir}[Rôles de la fermentation dans la production alimentaire]
\begin{enumerate}
\item \textbf{Conservation des aliments} : l'alcool, l'acide lactique et l'acidification du milieu empêchent le développement des microorganismes de décomposition et de nombreux agents pathogènes. Le lait caillé se conserve bien plus longtemps que le lait frais en l'absence de réfrigérateur.
\item \textbf{Amélioration du goût et de la texture} : arômes nouveaux, saveur acidulée ou alcoolisée, texture épaissie (yaourt) ou attendrie (manioc fermenté).
\item \textbf{Amélioration de la digestibilité} : les microorganismes « prédigèrent » une partie des sucres et des protéines ; le lactose du yaourt, partiellement dégradé, est mieux toléré par les personnes qui digèrent mal le lait.
\item \textbf{Enrichissement nutritionnel et détoxification} : certaines fermentations augmentent la teneur en vitamines (notamment du groupe B) et éliminent des composés toxiques (cyanures du manioc amer).
\end{enumerate}
\end{aretenir}

\courssub{Les conditions nécessaires à une fermentation réussie}

\begin{methode}[Décrire le principe d'un procédé de fermentation]
Pour décrire correctement un procédé de fermentation (compétence exigible au Probatoire), suis toujours la même chaîne logique, en cinq maillons :
\begin{enumerate}
\item \textbf{Substrat} : quelle matière organique est transformée ? (glucose de la sève, lactose du lait, amidon du maïs...)
\item \textbf{Microorganisme} : quel agent vivant réalise la transformation ? (levures, bactéries lactiques...)
\item \textbf{Conditions} : anaérobiose ou oxygène limité, température favorable (souvent \SIrange{25}{40}{\celsius}), milieu humide adapté, parfois salé ou acidifié.
\item \textbf{Produits obtenus} : éthanol + \ce{CO2}, acide lactique, arômes...
\item \textbf{Intérêt} : conservation, goût, digestibilité, valeur économique du produit.
\end{enumerate}
Exemple rédigé : « Le vin de palme est obtenu par fermentation alcoolique des sucres de la sève de palmier (substrat) par des levures naturelles (microorganisme), en anaérobiose à la température ambiante (conditions) ; il en résulte de l'éthanol et du dioxyde de carbone (produits), ce qui conserve la sève et lui donne son goût caractéristique (intérêt). »
\end{methode}

En résumé, quatre conditions sont indispensables : des \textbf{microorganismes vivants} en nombre suffisant, un \textbf{substrat} (sucre) disponible, une \textbf{température favorable} (ni trop froide, ce qui bloque les enzymes, ni trop chaude, ce qui tue les cellules), et un \textbf{milieu adapté} (humidité, acidité, absence d'oxygène selon le cas). Si l'une de ces conditions fait défaut, la fermentation ne démarre pas, s'arrête, ou dérive vers une contamination.

\courssub{Les limites des biotechnologies traditionnelles}

\begin{experience}[Mettre en évidence la variabilité d'une fermentation traditionnelle]
\textbf{Protocole.} Trois groupes d'élèves préparent chacun un jus d'ananas sucré dans des bouteilles identiques, ensemencé avec un peu de vin de palme frais, et le laissent fermenter \SI{48}{h} à température ambiante. Le groupe A bouche hermétiquement sa bouteille ; le groupe B la couvre d'un simple tissu ; le groupe C place sa bouteille près d'une fenêtre très chaude.
\textbf{Observations.} La bouteille A fermente correctement (odeur alcoolisée franche, bulles). La bouteille B présente par endroits une odeur de vinaigre : des bactéries acétiques de l'air ont transformé une partie de l'éthanol en acide acétique en présence d'oxygène. La bouteille C a fermenté vite puis s'est arrêtée : la chaleur excessive a tué une partie des levures.
\textbf{Interprétation.} Le même procédé « traditionnel » donne des résultats très différents selon les conditions : le produit final est \textbf{variable en qualité et en rendement}.
\textbf{Conclusion.} Sans maîtrise scientifique des microorganismes et des conditions, la fermentation traditionnelle reste empirique et aléatoire.
\end{experience}

Cette expérience illustre les trois grandes limites des biotechnologies traditionnelles :

\begin{itemize}
\item \textbf{Des procédés empiriques} : transmis oralement, ils reposent sur l'habitude (« on a toujours fait ainsi ») et non sur la connaissance des microorganismes et des paramètres physico-chimiques.
\item \textbf{Un rendement et une qualité variables} : la flore microbienne naturelle diffère d'un récipient, d'une saison ou d'une région à l'autre ; deux cuvées de vin de palme ne sont jamais identiques.
\item \textbf{Des risques de contamination} : des microorganismes indésirables (bactéries acétiques, moisissures, voire agents pathogènes) peuvent se développer si l'hygiène ou les conditions ne sont pas maîtrisées, altérant le produit ou le rendant dangereux.
\end{itemize}

\begin{attention}[Ce qu'on attend de toi au Probatoire]
Deux pièges classiques : (1) écrire que la fermentation alcoolique « consomme de l'oxygène » — c'est faux, elle se déroule en \textbf{anaérobiose} ; (2) affirmer que « plus la fermentation dure, plus l'alcool augmente indéfiniment » — c'est faux, la teneur en alcool \textbf{plafonne} car l'éthanol finit par tuer les levures et le sucre s'épuise. Pense aussi à toujours citer le microorganisme responsable : une fermentation sans microorganisme nommé est une réponse incomplète.
\end{attention}

\begin{exercice}[Pour t'entraîner]
Le « kossam » est un lait fermenté consommé dans le Grand Nord du Cameroun.
\begin{enumerate}
\item En suivant la méthode en cinq maillons (substrat $\rightarrow$ microorganisme $\rightarrow$ conditions $\rightarrow$ produits $\rightarrow$ intérêt), décris le principe de fabrication du kossam.
\item Explique pourquoi le kossam se conserve plus longtemps que le lait frais à température ambiante.
\item Cite deux différences entre la fermentation du kossam et celle du vin de palme.
\end{enumerate}
\end{exercice}

\begin{corrige}[Exercice : le kossam]
\begin{enumerate}
\item Substrat : le lactose du lait. Microorganisme : des bactéries lactiques naturellement présentes ou apportées par un ferment (un peu de lait déjà caillé). Conditions : anaérobiose relative, température ambiante favorable, récipient propre. Produits : acide lactique (et arômes), qui fait coaguler les protéines du lait. Intérêt : conservation du lait, goût acidulé apprécié, meilleure digestibilité.
\item L'acide lactique produit abaisse le pH du milieu ; cette acidité inhibe la croissance des microorganismes de décomposition et de nombreux agents pathogènes, qui ne supportent pas un milieu acide. Le lait fermenté se conserve donc mieux que le lait frais.
\item Deux différences : (a) le microorganisme responsable est une bactérie lactique pour le kossam, une levure pour le vin de palme ; (b) les produits sont l'acide lactique (sans dégagement gazeux) pour le kossam, contre éthanol + \ce{CO2} (avec bulles) pour le vin de palme.
\end{enumerate}
\end{corrige}

En définitive, la fermentation est la biotechnologie traditionnelle par excellence : simple, ancienne, profondément ancrée dans les pratiques alimentaires camerounaises, mais limitée par son caractère empirique. Comment l'être humain a-t-il dépassé ces limites ? C'est précisément l'objet de la section suivante, consacrée aux biotechnologies modernes : la sélection assistée et les biotechnologies enzymatiques.

\coursec{Les biotechnologies modernes : sélection assistée et biotechnologies enzymatiques}

Dans la section précédente, nous avons vu que la fermentation est une biotechnologie \cle{traditionnelle} : l'Homme y exploite des microorganismes entiers (levures, bactéries) sans toujours comprendre les mécanismes moléculaires en jeu. Depuis le XX\textsuperscript{e} siècle, les progrès de la génétique, de la biologie moléculaire et de la biochimie ont permis de passer à des biotechnologies \cle{modernes} : au lieu d'utiliser l'organisme entier « au hasard », on \textbf{oriente} les croisements grâce aux connaissances génétiques (\emph{sélection assistée}) ou on \textbf{isole} la molécule active elle-même, l'enzyme, pour réaliser une réaction précise (\emph{biotechnologies enzymatiques}). Cette section présente ces deux grandes familles de biotechnologies modernes et leurs applications concrètes, notamment au Cameroun.

\courssub{La sélection assistée : orienter les croisements grâce à la génétique}

\textbf{Une situation concrète pour commencer.} Dans la région de l'Extrême-Nord du Cameroun, les paysans constatent que certaines parcelles de maïs résistent mieux à la sécheresse que d'autres, et que certains plants de manioc échappent à la mosaïque (maladie virale qui détruit les feuilles). Comment obtenir, de façon fiable et rapide, des variétés qui combinent \emph{bon rendement} et \emph{résistance} ? La réponse traditionnelle — replanter les graines des meilleurs plants (sélection massale) — fonctionne, mais elle est lente, imprécise et ne permet pas de dissocier les caractères liés. La sélection assistée apporte une démarche dirigée et accélérée.

\begin{definition}[La sélection assistée par marqueurs (SAM)]
La \cle{sélection assistée} (ou \emph{sélection assistée par marqueurs}, SAM) est une biotechnologie moderne qui consiste à \textbf{choisir et croiser de façon dirigée} des organismes (végétaux ou animaux) présentant des \textbf{caractères d'intérêt} (rendement élevé, résistance aux maladies, tolérance à la sécheresse, qualité nutritive), en s'aidant des connaissances en génétique et de \textbf{marqueurs moléculaires} (séquences d'ADN polymorphes, comme les SNP ou SSR, statistiquement associées au gène d'intérêt), afin d'obtenir des \textbf{variétés ou races améliorées} sans avoir à phénotyper le caractère cible à l'âge adulte.
\end{definition}

\textbf{Principe et démarche.} Le caractère recherché (par exemple la résistance à la mosaïque du manioc) est porté par un ou plusieurs gènes (QTL). La démarche d'investigation du sélectionneur suit les étapes suivantes :

\begin{enumerate}
\item \textbf{Observation} : dans une population ségrégante (issue d'un croisement entre une lignée donneuse du caractère et une lignée élite), certains individus présentent le caractère d'intérêt (ex. : un plant de manioc reste sain dans une parcelle infestée).
\item \textbf{Hypothèse} : ce caractère est héréditaire, c'est-à-dire transmissible à la descendance selon les lois de Mendel (ou polygénique).
\item \textbf{Vérification / Génotypage} : on extrait l'ADN de jeunes plantules. L'identification de \emph{marqueurs moléculaires} liés au gène d'intérêt (liaison génétique) permet de repérer très tôt (stade plantule, sans attendre l'âge adulte ni l'expression phénotypique) les individus qui portent l'allèle favorable : c'est la \cle{sélection assistée par marqueurs}.
\item \textbf{Sélection et multiplication} : seuls les descendants qui combinent les caractères recherchés (génotype favorable + bon fond génétique) sont conservés, multipliés (souvent par culture \emph{in vitro}) et diffusés aux agriculteurs après essais multilocalisés.
\end{enumerate}

\begin{popfigure}
\begin{center}
\begin{tikzpicture}[
  etape/.style={draw=popblue, fill=popblueL, rounded corners=3pt, text width=3.0cm, align=center, minimum height=1.2cm, font=\small},
  fleche/.style={-{Stealth[length=3mm]}, thick, popdark}]
\node[etape, align=center] (p1) at (0,0){\textbf{Parents choisis}\\ Lignée donneuse (résistance) $\times$ Lignée élite (rendement)};
\node[etape, fill=poporangeL, draw=poporange, align=center] (p2) at (4.8,0){\textbf{Croisement dirigé}\\ Pollinisation / accouplement contrôlé $\rightarrow$ F1};
\node[etape, fill=poppurpleL, draw=poppurple, align=center] (p3) at (9.6,0){\textbf{Population ségrégante (F2, BC)}\\ Individus variés phénotypiquement};
\node[etape, fill=popgreenL, draw=popgreen, align=center] (p4) at (4.8,-2.6){\textbf{Génotypage (SAM)}\\ Repérage allèles favorables via marqueurs ADN};
\node[etape, fill=popgoldL, draw=popgold, align=center] (p5) at (0,-2.6){\textbf{Sélection \& Multiplication}\\ Variété améliorée homogène};
\draw[fleche] (p1) -- (p2);
\draw[fleche] (p2) -- (p3);
\draw[fleche] (p3.south) |- (p4.east);
\draw[fleche] (p4) -- (p5);
\draw[fleche, dashed, popmuted] (p5.north) -- (p1.south) node[midway, right, font=\scriptsize, popmuted]{Récurrence / Nouveaux cycles};
\end{tikzpicture}
\end{center}
\legende{Étapes de la sélection assistée par marqueurs (SAM) : les parents porteurs d'allèles d'intérêt sont croisés ; dans la descendance ségrégante, le génotypage précoce (marqueurs ADN) permet de sélectionner les individus recombinants porteurs de la combinaison idéale, qui sont ensuite multipliés.}
\end{popfigure}

\begin{exemplebox}[Variétés améliorées diffusées au Cameroun grâce à la SAM]
La sélection assistée a permis d'obtenir et de diffuser au Cameroun, sous l'égide de l'IRAD (Institut de Recherche Agricole pour le Développement) et de partenaires internationaux (IITA, CIRAD) :
\begin{itemize}
\item des variétés de \textbf{maïs} à haut rendement, à cycle court (90-100 jours), tolérantes à la sécheresse et à la striga (plante parasite), adaptées aux zones soudano-sahéliennes (Extrême-Nord, Nord) ;
\item des variétés de \textbf{manioc} résistantes à la mosaïque africaine du manioc (CMD) et à la pourriture brune des racines, à teneur élevée en matière sèche (\textgreater 30\%) ;
\item des variétés de \textbf{plantain} (ex. hybrides FHIA) à gros régimes, cycle court et tolérantes au cercosporiose noire (Sigatoka) ;
\item des \textbf{races animales améliorées} : croisement de zébus locaux (Kapsiki, Gudali - résistants à la chaleur, aux tiques, trypanotolérants) avec des races spécialisées (Holstein, Montbéliarde pour le lait ; Charolaise pour la viande) pour obtenir des animaux à la fois rustiques et productifs (double aptitude).
\end{itemize}
\end{exemplebox}

\begin{savaistu}[La culture in vitro (micropropagation) : un levier indispensable]
Complément utile au Probatoire : la \cle{culture in vitro} (ou micropropagation) consiste à prélever un petit fragment de végétal (méristème apical, bourgeon axillaire) et à le placer sur un \textbf{milieu de culture stérile} (gélosé type Murashige \& Skoog) contenant eau, sels minéraux, sucres (saccharose), vitamines et \textbf{régulateurs de croissance} (auxines, cytokinines). Le fragment se développe en un plant entier, génétiquement identique au parent (clone). Ce procédé permet de :
\begin{itemize}
\item \textbf{Multiplier rapidement et en grand nombre} des plants sains (ex. : 10 000 plants de bananier à partir d'un seul méristème en 6 mois) ;
\item \textbf{Assainir} le matériel végétal : la thermothérapie combinée à la culture de méristème élimine les virus (ex. : BBTV sur bananier, CMV sur manioc) ;
\item \textbf{Conserver} le matériel génétique (banque de gènes à croissance lente ou cryoconservation).
\end{itemize}
Au Cameroun, elle est utilisée massivement pour produire des plants vitro de \textbf{bananiers/plantains} (projet PNDRT, IRAD, structures privées) et de \textbf{palmiers à huile} (variétés Tenera à haut rendement huileux). La culture in vitro est souvent associée à la sélection assistée : une fois le meilleur génotype identifié, on le multiplie massivement pour la diffusion.
\end{savaistu}

\courssub{Les biotechnologies enzymatiques : utiliser des enzymes isolées dans l'industrie}

\textbf{Observation de départ.} Un jus de goyave fraîchement pressé est trouble ; après quelques heures en présence d'une préparation commerciale, il devient limpide. Une viande dure devient tendre après trempage dans du jus de papaye verte (latex). Une tache de sauce sur un vêtement disparaît avec une lessive « enzymatique » lavant à 30-40 °C. Dans chaque cas, ce ne sont pas des organismes vivants entiers qui agissent, mais des \textbf{protéines isolées, purifiées} : les \cle{enzymes}.

\begin{definition}[Biotechnologie enzymatique]
Une \cle{biotechnologie enzymatique} utilise des \textbf{enzymes extraites et purifiées d'organismes vivants} (bactéries, champignons filamenteux, tissus animaux ou végétaux) pour \textbf{catalyser une réaction chimique précise} utile à la production d'un bien ou d'un service. L'enzyme, en tant que \emph{catalyseur biologique}, accélère la réaction sans être consommée ni modifiée de façon permanente.
\end{definition}

\begin{attention}[Distinction fondamentale : Enzyme vs Microorganisme]
Erreur fréquente et sanctionnée au Probatoire : confondre l'enzyme avec le microorganisme qui la produit.
\begin{itemize}
\item Une enzyme est une \textbf{protéine} (macromolécule), c'est-à-dire une \textbf{molécule inerte}, produite par un organisme vivant ; elle n'est pas elle-même vivante : elle ne se nourrit pas, ne se reproduit pas, ne respire pas.
\item Dans la \textbf{fermentation} (biotechnologie traditionnelle), on utilise l'organisme \emph{entier et vivant} (métabolisme complet).
\item Dans la \textbf{biotechnologie enzymatique}, on utilise la \emph{molécule isolée} (activité catalytique unique).
\end{itemize}
Ne jamais écrire : « les enzymes se multiplient dans le milieu », « les enzymes meurent à haute température ». On écrit : « l'enzyme est \emph{dénaturée} (perte de sa conformation tridimensionnelle native) par la chaleur ou un pH extrême ».
\end{attention}

\begin{methode}[Décrire le principe d'un procédé enzymatique industriel]
Pour décrire correctement un procédé enzymatique dans une copie, précise toujours les quatre éléments suivants :
\begin{enumerate}
\item l'\textbf{enzyme} utilisée (nom précis, ex. : $\alpha$-amylase, pectinase, chymosine) ;
\item le \textbf{substrat} transformé (ex. : amidon, pectine, caséine $\kappa$) ;
\item le(s) \textbf{produit(s)} obtenu(s) (ex. : dextrines/maltose/glucose, galacturonates, para-$\kappa$-caséine + macropeptide) ;
\item l'\textbf{application industrielle} visée (ex. : brassage, clarification jus, fromagerie, détergence).
\end{enumerate}
\end{methode}

\textbf{Principaux exemples d'enzymes utilisées dans l'industrie (connaissances exigible Probatoire).}

\begin{center}
\begin{tabularx}{\linewidth}{|l|X|X|X|}
\hline
\rowcolor{popblueL} \textbf{Enzyme (classe EC)} & \textbf{Substrat transformé} & \textbf{Produit(s) principal(aux)} & \textbf{Application industrielle majeure} \\
\hline
\textbf{Amylases} (hydrolases) & Amidon (maïs, sorgho, manioc) & Dextrines, maltose, glucose (sirop) & Brasserie (saccharification), amidonnerie (sirop de glucose/fructose), boulangerie, textile (désencollage) \\
\hline
\textbf{Chymosine} (protéase aspartique) & Caséine $\kappa$ (protéine du lait) & Para-$\kappa$-caséine (coagulum) + Glycomacropeptide & Fabrication du fromage (coagulation du lait) \\
\hline
\textbf{Pectinases} (hydrolases : pectinestérase, polygalacturonase, pectate lyase) & Pectine (paroi cellulaire fruits) & Oligogalacturonides, galacturonate & Clarification jus de fruits (mangue, goyave, ananas), extraction huile de palme, macération vinicole \\
\hline
\textbf{Protéases} (sérine, cystéine, métalloprotéases) & Protéines (collagène, kératine, caséine) & Peptides, acides aminés & Détergents (lessives basses températures), attendrissage viande, dépilage/détrempe cuirs, hydrolysats protéiques \\
\hline
\textbf{Cellulases / Hémicellulases} & Cellulose, hémicelluloses & Glucose, xylose, oligosaccharides & Bioéthanol 2G (biomasse lignocellulosique), alimentation animale, papier (bioblanchiment) \\
\hline
\end{tabularx}
\end{center}

\begin{exemplebox}[Les amylases en brasserie et amidonnerie : l'étape de saccharification]
En brasserie artisanale ou industrielle (utilisant maïs, sorgho, riz), l'amidon (polymère de glucose $\alpha$-1,4/$\alpha$-1,6) ne peut pas être utilisé directement par les levures (\emph{Saccharomyces cerevisiae}), car c'est une molécule trop grosse pour traverser la membrane plasmique. Les \textbf{amylases} (endogènes du malt ou exogènes commerciales fongiques/bactériennes) le découpent en \textbf{sucres simples fermentescibles} (glucose, maltose, maltotriose) selon les réactions d'hydrolyse :
\[ \text{Amidon} + n\,\text{H}_2\text{O} \xrightarrow{\text{amylases}} \text{dextrines} \rightarrow \text{maltose} + \text{glucose} \]
Ces sucres seront ensuite transformés en éthanol et $\text{CO}_2$ par les levures lors de la fermentation alcoolique : \ce{C6H12O6 -> 2 C2H5OH + 2 CO2}. L'action des amylases (saccharification) est donc l'\textbf{étape préalable indispensable} à la production de bière ou de bioéthanol à partir de céréales/féculents.
\end{exemplebox}

\begin{savaistu}[La chymosine : de l'estomac du veau aux fermenteurs recombinants]
La \textbf{chymosine} (enzyme coagulante de la présure) était autrefois extraite de la caillette (abomasum), c'est-à-dire de l'estomac des jeunes veaux non sevrés, ce qui posait des problèmes de coût, de disponibilité saisonnière, de variabilité qualitative et d'acceptabilité (végétariens, prescriptions religieuses \emph{halal}/\emph{kasher}). Depuis les années 1990, elle est le plus souvent \textbf{produite par des microorganismes GMOs} (\emph{Aspergillus niger}, \emph{Kluyveromyces lactis}, \emph{Escherichia coli} K12) cultivés en fermenteurs, dans lesquels le gène \emph{chymosin} bovin a été inséré par recombinant DNA technology. L'enzyme recombinante (Fermentation-Produced Chymosin, FPC) est structurellement identique à l'enzyme native, plus pure, moins chère et répond aux exigences \emph{halal}/\emph{kasher}/\emph{vegetarian}. C'est l'exemple historique et emblématique de la transition entre biotechnologie traditionnelle (extraction tissulaire) et biotechnologie moderne (production recombinante).
\end{savaistu}

\textbf{Avantages des enzymes industrielles par rapport aux procédés chimiques classiques.} Par rapport à l'utilisation d'acides concentrés ($\text{H}_2\text{SO}_4$ pour l'hydrolyse d'amidon), de bases fortes ou de températures/pressions extrêmes, l'utilisation d'enzymes présente des avantages majeurs :
\begin{itemize}
\item \textbf{Spécificité absolue} : l'enzyme n'agit que sur un substrat précis (ex. : $\alpha$-1,4 liaison pour l'amylase), ce qui limite drastiquement les produits secondaires indésirables (couleur, goût, toxines) et préserve les composés d'intérêt (arômes, vitamines).
\item \textbf{Vitesse catalytique élevée} : nombre de tours catalytiques ($k_{\text{cat}}$) de $10^3$ à $10^7\,\text{s}^{-1}$ ; une enzyme transforme des milliers de molécules de substrat par seconde.
\item \textbf{Conditions douces (économie d'énergie)} : température modérée (30-60 °C généralement), pression atmosphérique, pH proche de la neutralité ou faiblement acide/basique.
\item \textbf{Respect de l'environnement} : les enzymes sont biodégradables (protéines), non toxiques, et remplacent des réactifs chimiques agressifs, corrosifs et polluants (effluents moins chargés).
\item \textbf{Contrôlabilité} : l'activité s'arrête par simple inactivation thermique (pasteurisation) ou changement de pH.
\end{itemize}

\courssub{Production industrielle des enzymes : le fermenteur et l'amont/aval}

Les enzymes industrielles sont produites par des \textbf{microorganismes producteurs} (souches sélectionnées, souvent améliorées par mutagenèse classique ou ingénierie métabolique/GMO : \emph{Bacillus subtilis}, \emph{Aspergillus niger}, \emph{Trichoderma reesei}) cultivés en grand volume dans des cuves appelées \cle{fermenteurs} (ou bioréacteurs). Le procédé global comprend l'\textbf{amont} (production de la biomasse/enzyme) et l'\textbf{aval} (récupération/purification).

\begin{popfigure}
\begin{center}
\begin{tikzpicture}[font=\small, scale=0.95, transform shape]
% Cuve principale
\draw[thick, popdark, fill=popblueL!50] (0,0) -- (0,4.5) arc (180:0:2.5) -- (5,0) arc (0:-180:2.5) -- cycle;
% Niveau liquide
\draw[thick, popdark, fill=popblue!20] (0,1.0) -- (0,4.5) arc (180:0:2.5) -- (5,1.0) arc (0:-180:2.5) -- cycle;
% Agitateur (arbre + pales)
\draw[thick, popdark] (2.5,5.8) -- (2.5,1.5);
\draw[thick, popdark, fill=popmuted!40] (1.8,1.5) -- (3.2,1.5) -- (3.5,1.0) -- (1.5,1.0) -- cycle; % pale bas
\draw[thick, popdark, fill=popmuted!40] (1.8,3.2) -- (3.2,3.2) -- (3.5,2.7) -- (1.5,2.7) -- cycle; % pale haut
% Moteur
\draw[thick, popdark, fill=popgray!30] (1.8,5.8) rectangle (3.2,6.3);
\node[align=center, font=\tiny] at (2.5,6.05) {Moteur};
% Arrivées (gauche)
\draw[-{Stealth[length=2.5mm]}, thick, popgreen] (-1.8,3.8) -- (0.05,3.8);
\node[left, align=right, popgreen, font=\small] at (-1.9,3.8) {Apport stérile\\ nutriments (C, N, P, sels)};
\draw[-{Stealth[length=2.5mm]}, thick, poporange] (-1.8,2.8) -- (0.05,2.8);
\node[left, align=right, poporange, font=\small] at (-1.9,2.8) {Air stérile\\ (aération / $\text{O}_2$)};
% Sonde température/pH (droite haut)
\draw[thick, popgold] (4.8,5.5) -- (4.8,2.0);
\draw[popgold, fill=popgold] (4.8,2.0) circle (0.08);
\node[right, align=left, popgold, font=\small] at (5.0,5.5) {Sondes\\ Temp / pH / pO2};
% Contrôle (jackets)
\draw[thick, poppurple, dashed] (-0.2,0) -- (-0.2,4.5) arc (180:0:2.7) -- (5.2,4.5) -- (5.2,0);
\node[left, align=right, poppurple, font=\small] at (-0.5,2.2) {Double enveloppe\\ (eau glacée/vapeur)\\ Thermorégulation};
% Sortie / Prélèvement (droite bas)
\draw[-{Stealth[length=2.5mm]}, thick, poppurple] (5.05,1.2) -- (6.8,1.2);
\node[right, align=left, poppurple, font=\small] at (6.9,1.2) {Prélèvement bouillie\\ $\rightarrow$ Séparation\\ (centrifugation/filtration)};
% Légende interne
\node[align=center, font=\small] at (2.5,2.7) {Culture de microorganismes\\ producteurs (biomasse + enzyme\\ intra/extra-cellulaire)};
\end{tikzpicture}
\end{center}
\legende{Schéma d'un fermenteur industriel à cuve agitée (type \emph{stirred-tank reactor}) : les microorganismes producteurs se développent en milieu liquide stérile, nourris en substrats carbonés/azotés et en air stérile, sous agitation mécanique permanente pour l'homogénéité et le transfert d'$\text{O}_2$ ; la température et le pH sont contrôlés en continu via sondes et double enveloppe ; le milieu de fermentation (bouillie) est ensuite prélevé pour l'extraction (rupture cellulaire si enzyme intracellulaire, ou centrifugation/filtration directe si extracellulaire) et la purification (précipitation, chromatographie, ultrafiltration, lyophilisation).}
\end{popfigure}

\begin{attention}[Conditions opératoires et stabilité enzymatique]
Une enzyme est une protéine \textbf{thermolabile} et \textbf{sensible au pH} : une \textbf{température trop élevée} (généralement \textgreater 50-60 °C pour les enzymes mésophiles) ou un \textbf{pH extrême} (loin de l'optimum) la \textbf{dénature} (perte irréversible de sa conformation tridimensionnelle native $\rightarrow$ perte d'activité). En industrie comme au laboratoire, il faut donc maintenir les conditions optimales de température ($T_{\text{opt}}$) et de pH ($\text{pH}_{\text{opt}}$) pendant la réaction. Ne jamais écrire qu'une enzyme « est tuée » ou « meurt » par la chaleur : une molécule ne meurt pas, elle est \emph{dénaturée} (ou inactivée thermiquement).
\end{attention}

\begin{exoresolu}[Décrire un procédé enzymatique : clarification de jus de mangue]
\textbf{Énoncé.} Une unité de transformation agroalimentaire au Cameroun produit du jus de mangue clarifié. Le jus pressé est trouble. On ajoute une préparation enzymatique commerciale. Après incubation à 45 °C pendant 2 h, le jus devient limpide. On pasteurise ensuite le jus.
\begin{enumerate}
\item Identifier la biotechnologie mise en jeu et justifier.
\item Décrire le principe du procédé en précisant l'enzyme, le substrat, le(s) produit(s) et l'application.
\item Citer deux avantages de ce procédé enzymatique par rapport à une clarification chimique (ex. : gélatine, bentonite, PVPP) ou physique seule (centrifugation).
\item Expliquer pourquoi l'étape de pasteurisation finale est indispensable vis-à-vis de l'enzyme.
\end{enumerate}
\tcblower
\textbf{Solution détaillée.}
\begin{enumerate}
\item Il s'agit d'une \textbf{biotechnologie enzymatique}. Justification : on utilise une \textbf{enzyme isolée et purifiée} (préparation commerciale) pour catalyser une réaction précise (dégradation de la pectine), et non un organisme vivant entier (ce qui serait de la fermentation).
\item Principe du procédé :
\begin{itemize}
\item \textbf{Enzyme} : \textbf{pectinase} (mélange de pectinestérase, polygalacturonase, pectate lyase) d'origine fongique (\emph{Aspergillus niger} généralement) ;
\item \textbf{Substrat} : la \textbf{pectine} (hétéropolysaccharide riche en acide galacturonique), composant majeur de la paroi cellulaire et de la lamelle moyenne des fruits, responsable de la turbidité (colloïdes pectiques) et de la viscosité ;
\item \textbf{Produits} : \textbf{oligogalacturonides} et \textbf{galacturonate} (solubles, non colloïdaux) $\rightarrow$ disparition de la turbidité (jus limpide) et baisse de la viscosité (meilleur rendement de pressage/filtration) ;
\item \textbf{Application industrielle} : clarification et augmentation du rendement d'extraction des jus de fruits tropicaux (mangue, goyave, ananas, fruit de la passion).
\end{itemize}
\item Avantages (au moins deux) :
\begin{itemize}
\item \textbf{Spécificité} : seule la pectine est hydrolysée ; les arômes, vitamines (vitamine C), pigments (caroténoïdes) et sucres sont préservés (qualité organoleptique et nutritionnelle supérieure).
\item \textbf{Conditions douces} : réaction à 45-50 °C, pH 3,5-4,5 (pH naturel du jus) $\rightarrow$ économie d'énergie vs traitement thermique seul ou chimique.
\item \textbf{Écologique} : l'enzyme est biodégradable, pas de résidus chimiques dans l'effluent ni le produit final (pas d'allergènes type gélatine/bentonite).
\item \textbf{Rendement} : dépectinisation $\rightarrow$ meilleur débourbage, filtration plus rapide, rendement jus +10-15\%.
\end{itemize}
\item La pasteurisation (ex. 90-95 °C, 15-30 s) \textbf{inactive (dénature) définitivement l'enzyme pectinase résiduelle}. Sans cette étape, l'enzyme resterait active dans la bouteille (même à 4 °C, activité résiduelle lente) et continuerait à dégrader la pectine $\rightarrow$ déstabilisation du jus (dépôt, perte de corps, trouble secondaire) et risque de fermentation si levures présentes. C'est une étape de \textbf{stabilisation enzymatique}.
\end{enumerate}
\end{exoresolu}

\begin{aretenir}[L'essentiel de la section : Biotechnologies Modernes]
\begin{itemize}
\item \textbf{Sélection assistée (SAM)} : utilisation de \textbf{marqueurs moléculaires ADN} (SNP, SSR) liés aux gènes d'intérêt pour sélectionner précocement (stade plantule) les meilleurs recombinants. Permet d'obtenir rapidement des variétés végétales (maïs, manioc, plantain, riz, coton) et races animales (bovins, porcins) améliorées (rendement + résistance/adaptation). Couplée à la \textbf{culture in vitro} (micropropagation de méristèmes) pour multiplication massive de plants sains (assainis virus) et diffusion (bananier, palmier à huile, pomme de terre, ananas au Cameroun).
\item \textbf{Biotechnologies enzymatiques} : utilisation d'\textbf{enzymes purifiées} (protéines catalytiques, \emph{non vivantes}) pour des réactions ciblées : \textbf{amylases} (saccharification amidon $\rightarrow$ brasserie, sirops, bioéthanol), \textbf{chymosine} (coagulation lait $\rightarrow$ fromage), \textbf{pectinases} (clarification jus), \textbf{protéases} (détergents, cuir, alimentation), \textbf{cellulases} (bioéthanol 2G, fourrages).
\item \textbf{Production} : culture de souches microbiennes productrices (souvent GMOs pour chymosine, amylases, protéases) en \textbf{fermenteurs} (contrôle $T$, pH, p$\text{O}_2$, agitation, stérilité) $\rightarrow$ extraction (rupture/centrifugation) $\rightarrow$ purification (concentration, formulation).
\item \textbf{Atouts enzymes} : spécificité, vitesse, conditions douces (énergie), biodégradables (environnement), contrôlables (inactivation thermique).
\item \textbf{Vigilance vocabulaire} : Enzyme = protéine (molécule) $\neq$ Microorganisme (vivant). Dénaturation $\neq$ mort. SAM $\neq$ Transgenèse (pas d'introduction de gène étranger dans le génome de la variété finale en SAM classique).
\end{itemize}
\end{aretenir}

\coursec{Comparaison entre biotechnologies traditionnelles et modernes}

Dans la section précédente, nous avons découvert les biotechnologies modernes : la sélection assistée, qui choisit les meilleurs organismes grâce à des marqueurs moléculaires, et les biotechnologies enzymatiques, qui exploitent des enzymes purifiées pour conduire des réactions précises. Avant cela, nous avions étudié la fermentation, pratiquée depuis des millénaires sans qu'on en connaisse le mécanisme. Une question s'impose alors naturellement : qu'est-ce qui distingue réellement une biotechnologie « traditionnelle » d'une biotechnologie « moderne » ? L'une est-elle supérieure à l'autre ? C'est à ces questions, essentielles pour le Probatoire, que cette section est consacrée.

\courssub{Pourquoi comparer ? Une situation concrète pour commencer}

Sur un marché de Douala, une vendeuse propose du vin de palme fraîchement récolté : la sève de palmier fermente spontanément grâce aux levures présentes dans l'air et sur l'écorce de l'arbre. À quelques kilomètres de là, une brasserie industrielle produit chaque jour des milliers de bouteilles de bière en utilisant des souches de levures sélectionnées en laboratoire, dans des cuves dont la température et l'acidité sont contrôlées en permanence.

Dans les deux cas, on utilise des \cle{microorganismes vivants} (des levures) pour transformer des sucres en alcool : c'est bien de la biotechnologie. Pourtant, les deux procédés diffèrent profondément. Comparer ces deux démarches, c'est comprendre ce que le progrès scientifique a réellement apporté — et ce qu'il n'a pas remplacé.

\begin{methode}[Comment comparer deux biotechnologies ?]
Pour comparer rigoureusement une biotechnologie traditionnelle et une biotechnologie moderne (savoir-faire exigible au Probatoire), il faut procéder en étapes :
\begin{enumerate}
\item Choisir des \cle{critères de comparaison identiques} pour les deux procédés : ancienneté, fondement scientifique, agents utilisés, maîtrise du procédé, rendement, coût.
\item Décrire chaque procédé selon chacun de ces critères, sans en oublier aucun.
\item Présenter les résultats dans un \cle{tableau à double entrée} (critères en lignes, types de biotechnologie en colonnes).
\item Conclure en dégageant les \cle{points communs} (ce qui fait que les deux sont des biotechnologies) et en discutant la complémentarité des deux approches.
\end{enumerate}
L'erreur à éviter : comparer deux procédés sur des critères différents, ce qui rend la comparaison sans valeur scientifique.
\end{methode}

\courssub{Les critères de comparaison, un par un}

\textbf{Ancienneté.}\par
Les biotechnologies traditionnelles sont \cle{millénaires} : la fermentation du pain, du vin, de la bière ou du lait est attestée depuis plusieurs milliers d'années (Égypte antique, Mésopotamie). Au Cameroun, la fabrication du vin de palme, du « kaffir » ou du lait caillé remonte à des générations. Les biotechnologies modernes sont, elles, très récentes : elles naissent au \textsc{xx}\textsuperscript{e} siècle, avec la découverte des microorganismes (Pasteur, vers 1860), puis de l'ADN et des techniques de biologie moléculaire (seconde moitié du \textsc{xx}\textsuperscript{e} siècle), permettant la sélection assistée par marqueurs et le génie génétique.

\textbf{Fondement scientifique.}\par
La biotechnologie traditionnelle repose sur un savoir-faire \cle{empirique} : on sait \emph{que} ça marche, sans savoir \emph{pourquoi}. Le palmier-tapster sait que la sève devient alcoolisée en quelques heures, mais il ignore que des levures transforment le glucose en éthanol. La biotechnologie moderne repose au contraire sur des connaissances scientifiques précises : la \cle{microbiologie} (nature et physiologie des microorganismes) et la \cle{biologie moléculaire} (ADN, gènes, enzymes, marqueurs moléculaires). On connaît l'agent responsable, la réaction chimique qu'il catalyse et les conditions optimales de son activité.

\textbf{Agents utilisés.}\par
Dans les procédés traditionnels, on utilise des microorganismes \cle{naturels}, présents spontanément dans l'environnement (levures de l'air, bactéries lactiques du lait cru) : la flore microbienne est donc variable d'une production à l'autre. Dans les procédés modernes, on utilise des agents \cle{sélectionnés} (souches pures de levures ou de bactéries choisies pour leur efficacité, obtenues par sélection classique ou sélection assistée) voire \cle{modifiés} (organismes génétiquement modifiés), ou encore des \cle{enzymes purifiées} extraites de ces organismes.

\textbf{Maîtrise du procédé.}\par
Le procédé traditionnel est peu contrôlé : température ambiante, durée variable, contamination possible par des microorganismes indésirables (d'où un vin de palme qui « tourne » rapidement en vinaigre sous l'action de bactéries acétiques). Le procédé moderne est \cle{contrôlé} : température, pH, aération, stérilité et concentration en substrat sont surveillés en continu dans des fermenteurs instrumentés. La reproductibilité est donc bien meilleure : chaque lot de production est identique au précédent.

\textbf{Rendement.}\par
Le rendement d'un procédé traditionnel est \cle{variable} : il dépend du climat, de la saison, de la flore microbienne du jour. Le rendement d'un procédé moderne est \cle{élevé et régulier}, car l'agent sélectionné est optimisé et les conditions sont maintenues à leur valeur optimale.

\textbf{Coût.}\par
La biotechnologie traditionnelle exige un \cle{faible investissement} : quelques récipients, un savoir-faire transmis oralement. Elle est accessible à tous, y compris en milieu rural. La biotechnologie moderne exige un \cle{investissement lourd} : laboratoires, fermenteurs, personnel qualifié, chaîne du froid pour les ferments. En revanche, à grande échelle, le coût unitaire de production devient faible grâce aux rendements élevés.

\courssub{Le tableau comparatif synthétique}

La figure ci-dessous rassemble ces six critères dans un tableau comparatif, tel qu'on attend de toi au Probatoire.

\begin{popfigure}
\begin{tikzpicture}[x=1cm,y=1cm]
% Dimensions : 3 colonnes
\fill[popblueL] (0,0) rectangle (3.4,-0.7);
\fill[poporangeL] (3.4,0) rectangle (9.2,-0.7);
\fill[popgreenL] (9.2,0) rectangle (15,-0.7);
\node[anchor=west,font=\bfseries] at (0.1,-0.35) {Critère};
\node[font=\bfseries] at (6.3,-0.35) {Biotechnologie traditionnelle};
\node[font=\bfseries] at (12.1,-0.35) {Biotechnologie moderne};
% Lignes
\foreach \y/\crit/\trad/\mod in {
 -1.4/{Ancienneté}/{Millénaire (préhistoire, Antiquité)}/{Récente (\textsc{xx}\textsuperscript{e} siècle)},
 -2.4/{Fondement}/{Empirique : savoir-faire transmis}/{Scientifique : microbiologie, biologie moléculaire},
 -3.4/{Agents utilisés}/{Microorganismes naturels, flore spontanée}/{Souches sélectionnées ou modifiées, enzymes purifiées},
 -4.4/{Maîtrise du procédé}/{Faible : conditions ambiantes, peu contrôlées}/{Forte : paramètres contrôlés (T°, pH, stérilité)},
 -5.4/{Rendement}/{Variable, souvent faible}/{Élevé et reproductible},
 -6.4/{Coût}/{Faible investissement, accessible à tous}/{Investissement lourd (laboratoire, fermenteurs)}
}{
 \node[anchor=west,font=\bfseries\small] at (0.1,\y) {\crit};
 \node[anchor=west,text width=5.5cm,font=\small] at (3.5,\y) {\trad};
 \node[anchor=west,text width=5.5cm,font=\small] at (9.3,\y) {\mod};
}
% Grille
\draw[popline] (0,0) rectangle (15,-6.9);
\draw[popline] (3.4,0) -- (3.4,-6.9);
\draw[popline] (9.2,0) -- (9.2,-6.9);
\foreach \y in {-0.7,-1.9,-2.9,-3.9,-4.9,-5.9} {\draw[popline] (0,\y) -- (15,\y);}
\end{tikzpicture}
\legende{Tableau comparatif des biotechnologies traditionnelles et modernes selon six critères : ancienneté, fondement scientifique, agents utilisés, maîtrise du procédé, rendement et coût.}
\end{popfigure}

\begin{propriete}[Rendement et reproductibilité (admise)]
Les biotechnologies modernes offrent en général un \cle{meilleur rendement} et une meilleure \cle{reproductibilité} que les procédés traditionnels, grâce à la sélection des agents biologiques et au contrôle des conditions de culture. Ce gain de performance s'obtient au prix d'\cle{investissements plus lourds} (équipements, énergie, personnel qualifié). Cette propriété est admise : sa démonstration n'est pas exigible au Probatoire, mais tu dois savoir la justifier par les critères du tableau comparatif.
\end{propriete}

\courssub{Points communs : ce qui unit les deux approches}

Au-delà de leurs différences, les deux types de biotechnologie partagent une même essence, ce qui justifie qu'on les réunisse sous un même nom :

\begin{itemize}
\item Dans les deux cas, on utilise des \cle{organismes vivants} (microorganismes, plantes, animaux) ou leurs \cle{constituants} (enzymes, cellules, gènes).
\item Dans les deux cas, la finalité est la \cle{production de biens ou de services} utiles à l'Homme : aliments, boissons, médicaments, traitement des déchets.
\item Dans les deux cas, le procédé repose sur des \cle{réactions biologiques} (le plus souvent des réactions enzymatiques, comme la fermentation alcoolique : \ce{C6H12O6 -> 2C2H5OH + 2CO2}).
\end{itemize}

\begin{popfigure}
\begin{tikzpicture}[x=1cm,y=1cm]
% Cercles de Venn
\fill[poporange!35] (0,0) circle (3.1);
\fill[popblue!30] (4.4,0) circle (3.1);
\draw[poporange,thick] (0,0) circle (3.1);
\draw[popblue,thick] (4.4,0) circle (3.1);
% Titres
\node[font=\bfseries,color=popdark] at (-1.3,2.6) {Traditionnelle};
\node[font=\bfseries,color=popdark] at (5.7,2.6) {Moderne};
% Spécifique traditionnelle
\node[text width=3.4cm,align=center,font=\small] at (-1.2,0.9) {Empirique, millénaire};
\node[text width=3.4cm,align=center,font=\small] at (-1.2,-0.1) {Flore microbienne naturelle};
\node[text width=3.4cm,align=center,font=\small] at (-1.2,-1.1) {Faible coût, rendement variable};
% Spécifique moderne
\node[text width=3.4cm,align=center,font=\small] at (5.6,0.9) {Fondée sur la science};
\node[text width=3.4cm,align=center,font=\small] at (5.6,-0.1) {Souches sélectionnées, enzymes};
\node[text width=3.4cm,align=center,font=\small] at (5.6,-1.1) {Rendement élevé, coût lourd};
% Intersection
\node[text width=2.6cm,align=center,font=\small\bfseries,color=popdark] at (2.2,0.7) {Organismes vivants ou leurs constituants};
\node[text width=2.6cm,align=center,font=\small\bfseries,color=popdark] at (2.2,-0.7) {Production de biens et services};
\end{tikzpicture}
\legende{Diagramme de Venn comparant biotechnologies traditionnelle et moderne. L'intersection rassemble les points communs : l'utilisation d'organismes vivants ou de leurs constituants, dans le but de produire des biens ou des services.}
\end{popfigure}

\courssub{Complémentarité : quand le moderne améliore le traditionnel}

Il serait faux d'opposer les deux approches comme deux rivales. Dans la pratique, les biotechnologies modernes \cle{améliorent} le plus souvent des procédés traditionnels ancestraux, sans changer leur principe biologique.

\begin{exemplebox}[Du vin de palme artisanal à la brasserie industrielle]
\textbf{Procédé traditionnel : le vin de palme.} La sève récoltée au sommet du palmier contient des sucres. Des levures naturelles (présentes dans l'air, sur l'écorce, dans le récipient) réalisent spontanément la fermentation alcoolique : \ce{C6H12O6 -> 2C2H5OH + 2CO2}. Le produit est obtenu en quelques heures, sans aucun équipement, mais sa teneur en alcool varie, il se conserve mal et s'acidifie vite (fermentation acétique par des bactéries indésirables).

\textbf{Procédé moderne : la brasserie industrielle.} Le principe biologique est \emph{identique} : des levures (\emph{Saccharomyces cerevisiae}) transforment les sucres du moût en éthanol. Mais on utilise des \cle{souches de levures sélectionnées} pour leur rendement et leurs qualités gustatives, introduites en quantité connue dans un moût stérilisé, dans des fermenteurs où température, pH et oxygénation sont contrôlés. Résultat : un produit constant, de teneur en alcool maîtrisée, conservable et produit en grande quantité.

\textbf{Même exemple dans la laiterie :} le yaourt artisanal repose sur des ferments naturels ; l'industrie utilise les mêmes bactéries lactiques (\emph{Lactobacillus delbrueckii} subsp. \emph{bulgaricus}, \emph{Streptococcus thermophilus}) mais sous forme de ferments sélectionnés et dosés avec précision. La biotechnologie moderne n'a donc pas inventé la fermentation : elle l'a \emph{comprise}, puis \emph{optimisée}.
\end{exemplebox}

\begin{attention}[« Moderne » ne veut pas dire « toujours meilleur »]
Ne rédige jamais qu'une biotechnologie moderne est supérieure « en tout point ». Pour une \cle{petite production locale} (vin de palme du village, lait caillé familial, manioc fermenté), le procédé traditionnel reste \cle{parfaitement adapté} : il ne coûte presque rien, ne demande ni électricité ni laboratoire, et répond aux goûts des consommateurs locaux. Le procédé moderne n'est pertinent que lorsqu'on vise une \cle{production à grande échelle}, régulière et standardisée. Au Probatoire, une discussion nuancée des enjeux selon le contexte est toujours valorisée.
\end{attention}

\begin{savaistu}[Tradition et modernité main dans la main au Cameroun]
De nombreuses industries agroalimentaires camerounaises combinent savoir-faire traditionnel et biotechnologie moderne. Les brasseries du pays utilisent des ferments sélectionnés tout en valorisant des matières premières locales (maïs, sorgho) issues de recettes de boissons fermentées traditionnelles. De même, des unités de transformation du manioc s'appuient sur la fermentation traditionnelle du « bâton de manioc », tout en y introduisant des ferments sélectionnés pour sécuriser et régulariser la production. La recherche camerounaise travaille aussi à sélectionner les souches de levures du vin de palme pour en améliorer la conservation.
\end{savaistu}

\begin{exoresolu}[Construire un tableau comparatif]
\textbf{Énoncé.} Un élève affirme : « La fabrication du vin de palme et la production industrielle de bière n'ont rien en commun, car l'une est traditionnelle et l'autre moderne. » À l'aide de quatre critères de comparaison, montre que cette affirmation est à la fois exacte sur les différences et fausse sur l'essentiel.
\tcblower
\textbf{Solution détaillée.}

On applique la méthode de comparaison : critères identiques pour les deux procédés, puis conclusion sur les points communs.

\begin{center}
\begin{tabularx}{\linewidth}{|l|X|X|}
\hline
\rowcolor{popblueL} \textbf{Critère} & \textbf{Vin de palme (traditionnel)} & \textbf{Bière industrielle (moderne)} \\
\hline
Agents utilisés & Levures naturelles de l'air et de l'écorce, flore variable & Souches de \emph{Saccharomyces} sélectionnées, dose connue \\
\hline
Maîtrise du procédé & Conditions ambiantes, non contrôlées ; risque d'acidification & Température, pH et stérilité contrôlés en fermenteur \\
\hline
Rendement & Variable selon la saison et la flore & Élevé et constant d'un lot à l'autre \\
\hline
Coût & Quasi nul : récipient et savoir-faire & Investissement lourd : cuves, laboratoire, personnel \\
\hline
\end{tabularx}
\end{center}

\textbf{Conclusion.} Les deux procédés diffèrent bien par leurs agents, leur maîtrise, leur rendement et leur coût : l'élève a raison sur les différences. Mais il a tort sur l'essentiel : dans les deux cas, ce sont des \cle{levures vivantes} qui réalisent la \emph{même} réaction biologique, la fermentation alcoolique \ce{C6H12O6 -> 2C2H5OH + 2CO2}, pour produire un bien utile à l'Homme. Ce sont donc deux biotechnologies de même nature biologique, l'une empirique, l'autre scientifiquement optimisée.
\end{exoresolu}

\begin{aretenir}[L'essentiel de la section]
\begin{itemize}
\item Comparer deux biotechnologies exige des \cle{critères identiques} : ancienneté, fondement, agents, maîtrise, rendement, coût.
\item La biotechnologie traditionnelle est \cle{empirique}, utilise la flore naturelle, coûte peu, mais son rendement est variable.
\item La biotechnologie moderne est fondée sur la \cle{microbiologie et la biologie moléculaire}, utilise des agents sélectionnés ou modifiés, et offre un rendement élevé et reproductible, au prix d'investissements lourds.
\item Point commun fondamental : dans les deux cas, des \cle{organismes vivants ou leurs constituants} produisent des biens ou des services.
\item Les procédés modernes \cle{améliorent} le plus souvent des procédés traditionnels (ferments sélectionnés pour la bière ou le yaourt) ; « moderne » ne signifie pas « meilleur dans tous les contextes ».
\end{itemize}
\end{aretenir}

\begin{exercice}[À toi de jouer]
Une coopérative de producteurs de lait de l'Ouest camerounais hésite entre deux options pour fabriquer du yaourt : (a) conserver la méthode traditionnelle avec ensemencement par un reste de la production précédente ; (b) acheter des ferments lactiques sélectionnés et un incubateur thermostaté.
\begin{enumerate}
\item Construis un tableau comparant les deux options selon quatre critères.
\item Conseille la coopérative en distinguant le cas d'une petite production villageoise et celui d'une production destinée aux supermarchés de Yaoundé.
\end{enumerate}
\end{exercice}

\begin{corrige}[Exercice : À toi de jouer]
\textbf{1. Tableau comparatif.}
\begin{center}
\begin{tabularx}{\linewidth}{|l|X|X|}
\hline
\rowcolor{popblueL} \textbf{Critère} & \textbf{Option (a) traditionnelle} & \textbf{Option (b) moderne} \\
\hline
Agents & Flore lactique naturelle, variable d'un lot à l'autre & Ferments sélectionnés, souches pures et dosées \\
\hline
Maîtrise & Température ambiante, fermentation aléatoire & Température contrôlée par l'incubateur (environ 43 °C) \\
\hline
Rendement et qualité & Variables, risque d'échec ou de contamination & Élevés, goût et texture constants \\
\hline
Coût & Très faible (aucun équipement) & Investissement en ferments et en matériel \\
\hline
\end{tabularx}
\end{center}
\textbf{2. Conseil.} Pour une petite production villageoise vendue rapidement sur place, l'option (a) reste adaptée : coût quasi nul, savoir-faire local, pas de besoin en électricité fiable. Pour approvisionner les supermarchés de Yaoundé, l'option (b) s'impose : il faut un produit régulier, standardisé, sûr et conservable, ce que seuls des ferments sélectionnés et une température contrôlée garantissent. L'investissement sera amorti par le volume des ventes. On retrouve le principe : « moderne » est pertinent selon le contexte et l'échelle de production.
\end{corrige}

\coursec{Enjeux économiques et éthiques de la biotechnologie}

Après avoir comparé les biotechnologies traditionnelles et modernes, il est indispensable d'évaluer les conséquences de leur déploiement à l'échelle industrielle ou agricole. Toute innovation technique s'accompagne de bénéfices potentiels, mais aussi de risques et de questions de société. Cette section propose une analyse structurée, illustrée par des exemples camerounais, pour vous permettre de \cle{discuter} les enjeux d'une application biotechnologique donnée, comme l'exige le programme.

\courssub{Enjeux économiques positifs}

Les biotechnologies génèrent des activités industrielles et agricoles créatrices de valeur ajoutée et d'emplois, contribuant à la souveraineté nationale.

\begin{itemize}
    \item \textbf{Industries agroalimentaires :} les brasseries (production de bière par fermentation alcoolique), les laiteries (yaourts, fromages par fermentation lactique) et les unités de transformation du manioc en \cle{farine de foufou} ou en \cle{attieke} reposent sur des ferments sélectionnés (\emph{Saccharomyces cerevisiae}, bactéries lactiques). Au Cameroun, les \cle{Brasseries du Cameroun} (UCB, SOCABO) et de nombreuses PME locales emploient des milliers de personnes.
    \item \textbf{Semenciers et sélection assistée :} l'utilisation de marqueurs moléculaires (\cle{SAM}, \emph{Marker Assisted Selection}) accélère la création de variétés de maïs, de riz ou de plantain à haut rendement et résistantes aux maladies (ex. résistance au \emph{Fusarium} pour le bananier/plantain). Cela réduit les importations de semences et renforce la souveraineté alimentaire.
    \item \textbf{Santé humaine :} la production d'\cle{insuline humaine recombinante}, de vaccins (hépatite B, COVID-19 à ARNm, vaccins recombinants) et d'enzymes thérapeutiques par des micro-organismes modifiés (\emph{E. coli}, \emph{S. cerevisiae}, cellules CHO) diminue la dépendance aux importations coûteuses et améliore l'accès aux soins.
    \item \textbf{Valorisation des déchets et énergie :} la méthanisation de déchets agricoles, d'effluents d'abattoirs ou de déchets de marchés produit du \cle{biogaz} (énergie renouvelable, riche en \ce{CH4}) et du \cle{digestat} (engrais organique riche en azote minéral). C'est un levier d'assainissement urbain et d'économie circulaire.
\end{itemize}

\begin{exemplebox}[Production de biogaz à partir des déchets de marché à Yaoundé]
À Yaoundé, le marché de Mfoundi génère quotidiennement plusieurs tonnes de déchets organiques (épluchures, fruits abîmés). Un projet pilote de méthanisation a été installé : les déchets sont broyés, mélangés à de l'eau et introduits dans un digesteur anaérobie de \SI{500}{m^3}. La fermentation anaérobie produit un biogaz riche en méthane (\approx \SI{60}{\percent} \ce{CH4}) qui alimente un groupe électrogène de \SI{50}{kW}, éclairant le marché la nuit. Le digestat, après compostage aérobie, est vendu comme amendement organique aux maraîchers de la périphérie.
\textbf{Bénéfices :} énergie locale renouvelable, réduction du volume des déchets mis en décharge, production d'engrais organique, création d'emplois (collecte, exploitation, maintenance).
\textbf{Risques techniques :} investissement initial élevé (\approx \SI{150}{millions} FCFA), nécessité d'une maintenance technique pointue (contrôle pH, température, C/N), gestion des odeurs et des fuites de méthane (gaz à effet de serre, pouvoir de réchauffement global 28 fois celui du \ce{CO2} sur 100 ans).
\textbf{Encadrement :} respect des normes de sécurité (étanchéité, détection de fuites \ce{CH4}/\ce{H2S}), formation des opérateurs, contrôle régulier par le Ministère de l'Environnement, de la Protection de la Nature et du Développement Durable (MINEPDED).
\end{exemplebox}

\courssub{Enjeux économiques négatifs ou risques}

Malgré les opportunités, plusieurs freins économiques structurels existent, particulièrement dans les pays en développement.

\begin{itemize}
    \item \textbf{Coût élevé des équipements et de la R\&D :} les fermenteurs industriels (acier inoxydable, contrôle stérile, instrumentation), les séquenceurs haut débit, les plateformes de criblage enzymatique ou de protéomique nécessitent des investissements de plusieurs centaines de millions à milliards de FCFA, souvent inaccessibles aux PME sans appui public (Fonds de développement, partenariats public-privé) ou internationaux.
    \item \textbf{Dépendance aux intrants importés :} souches microbiennes brevetées (collections type ATCC, DSMZ), milieux de culture spécialisés (peptones, extraits de levure, facteurs de croissance), enzymes purifiées de restriction/modification, kits de diagnostic moléculaire — la plupart sont importés. Cela crée une vulnérabilité face aux fluctuations de change, aux barrières douanières et aux ruptures d'approvisionnement (ex. crise COVID-19).
    \item \textbf{Concurrence avec les petits producteurs :} l'industrialisation de la fermentation (ex. production de vin de palme, de \cle{bili-bili} ou de \cle{ngoundja} à grande échelle avec ferments standardisés) peut marginaliser les producteurs artisanaux, détenteurs de savoir-faire traditionnels et de biodiversité microbienne locale, si aucune mesure d'accompagnement (label d'origine, appui à la coopérative, certification qualité) n'est prévue.
    \item \textbf{Risque de monopole et propriété intellectuelle :} la brevetabilité des séquences géniques, des procédés d'obtention ou des variétés végétales (certificat d'obtention végétale - COV) peut concentrer le marché entre les mains de quelques firmes multinationales, limitant l'accès aux innovations pour les pays du Sud et les petits agriculteurs (droit de l'agriculteur à ressemer).
\end{itemize}

\courssub{Enjeux éthiques et juridiques}

Les biotechnologies modernes (ADN recombinant, édition de génome, thérapie génique) soulèvent des questions de fond sur le vivant, la justice distributive et la responsabilité intergénérationnelle.

\begin{definition}[Biosécurité]
La \cle{biosécurité} désigne l'ensemble des mesures législatives, réglementaires et techniques visant à prévenir les risques liés à l'utilisation des organismes génétiquement modifiés (OGM) et des micro-organismes pathogènes pour la santé humaine, animale et l'environnement, conformément au principe de précaution.
\end{definition}

\begin{itemize}
    \item \textbf{Brevetabilité du vivant et APA :} déposer un brevet sur un gène, une souche ou une variété soulève la question de l'appropriation privée d'un patrimoine génétique souvent issu de la biodiversité collective et des savoirs traditionnels associés (ex. plantes médicinales camerounaises comme \emph{Prunus africana}, \emph{Ancistrocladus korupensis}). Le \cle{Protocole de Nagoya} (2010) sur l'Accès aux ressources génétiques et le Partage juste et équitable des Avantages (APA) découlant de leur utilisation impose le consentement préalable en connaissance de cause (PIC) et des conditions mutuellement convenues (MAT).
    \item \textbf{Accès équitable aux innovations de santé :} les thérapies géniques (coût > \SI{1}{million} USD/traitement), les vaccins à ARNm ou les semences améliorées tolérantes à la sécheresse ne doivent pas rester l'apanage des pays riches. Les licences volontaires, les pools de brevets (\emph{Medicines Patent Pool}), les accords de transfert de technologie (ex. hub ARNm OMS en Afrique du Sud) et les licences d'office (flexibilités ADPIC) sont des leviers juridiques.
    \item \textbf{Risques sur la biodiversité et flux de gènes :} la dissémination accidentelle d'OGM (pollinisation croisée avec des variétés locales ou parents sauvages, fuite de micro-organismes modifiés en milieu ouvert) peut perturber les écosystèmes locaux (effet sur la faune auxiliaire, transfert horizontal de gènes de résistance aux antibiotiques). Le principe de précaution impose des essais confinés (niveaux L1 à L3) et une surveillance post-commercialisation (monitoring).
    \item \textbf{Acceptabilité sociale et justice environnementale :} au Cameroun comme ailleurs, la perception publique est influencée par des craintes sanitaires (allergénicité, toxicité), des considérations culturelles (alimentation traditionnelle, sacredness of seeds) et la méfiance envers les multinationales (« terminator seeds »). Une communication transparente, basée sur des données scientifiques indépendantes et une participation citoyenne (conférences de consensus), est indispensable.
\end{itemize}

\begin{savaistu}[Cadre réglementaire camerounais : Biosécurité et Nagoya]
Le Cameroun a ratifié le \cle{Protocole de Carthagène} sur la prévention des risques biotechnologiques (loi n°2003/006 du 21 avril 2003) et le \cle{Protocole de Nagoya} (loi n°2012/004 du 19 avril 2012). La loi de 2003 fixe le régime de la biosécurité : elle crée un \cle{Comité National de Biosécurité (CNB)} placé auprès du MINRESI, chargé d'évaluer les dossiers de demande d'essais en champ confiné, de mise sur le marché ou d'importation d'OGM. Tout organisme modifié doit faire l'objet d'une évaluation des risques environnementaux (ERA) et sanitaires avant toute utilisation. Le point focal national CBD (Convention sur la Diversité Biologique) assure la coordination.
\end{savaistu}

\courssub{Enjeux sanitaires et environnementaux : sécurité des produits et confinement}

La sécurité des produits (médicaments, aliments, engrais) et la protection de l'environnement sont des conditions \emph{sine qua non} du développement durable des biotechnologies.

\begin{itemize}
    \item \textbf{Sécurité des produits biopharmaceutiques :} tout médicament issu de biotechnologie (protéines recombinantes, anticorps monoclonaux, vaccins) doit subir des tests rigoureux de toxicité (aiguë, chronique), d'allergénicité, d'immunogénicité, de stabilité et de pureté (absence de contaminants : ADN hôte, protéines hôte, endotoxines/LPS, virus adventices). L'exemple de l'insuline recombinante montre que la purification par chromatographies successives (échange d'ions, exclusion stérique, phase inverse) élimine les contaminants pour garantir l'innocuité selon les monographies de la Pharmacopée.
    \item \textbf{Confinement biologique et physique :} les micro-organismes producteurs (bactéries, levures, champignons filamenteux) utilisés en fermentation industrielle sont souvent des souches \cle{auxotrophes} (dépendantes d'un nutriment absent dans la nature) ou porteuses de systèmes \emph{kill-switch} (suicide conditionnel) pour limiter leur survie en cas de fuite. Les effluents de fermentation sont stérilisés (thermique \SI{121}{\celsius} 20 min ou chimique) avant rejet.
    \item \textbf{Niveaux de confinement (L1--L4) :} la réglementation (OMS, directive UE, loi camerounaise) impose des niveaux de confinement (Laboratoire L1 à L4, serre P1 à P3) selon la dangerosité de l'organisme (pathogénicité, dissémination, existence de prophylaxie). Les laboratoires camerounais de référence (Centre Pasteur du Cameroun - CPC, IRAD, LANAVET, universités) appliquent ces normes pour manipuler des agents pathogènes (ex. \emph{Mycobacterium tuberculosis}, virus Ebola) ou des OGM.
\end{itemize}

\courssub{Biotechnologies au service du développement durable}

Bien encadrées, les biotechnologies offrent des solutions concrètes pour atteindre les Objectifs de Développement Durable (ODD 2, 3, 6, 7, 12, 13, 15).

\begin{itemize}
    \item \textbf{Réduction des intrants chimiques de synthèse (ODD 2, 15) :} les \cle{biopesticides} (ex. \emph{Bacillus thuringiensis} var. \emph{kurstaki} produisant des toxines Cry insecticides spécifiques, \emph{Trichoderma} spp. antagonistes fongiques) et les \cle{biofertilisants} (bactéries fixatrices d'azote \emph{Azotobacter}, \emph{Rhizobium} ; champignons mycorhiziens à arbuscules \emph{Glomus} spp. ; bactéries solubilisant le phosphate \emph{Pseudomonas}) remplacent partiellement les pesticides et engrais chimiques, préservant la microbiologie des sols, la qualité de l'eau et la santé des applicateurs.
    \item \textbf{Dépollution et bioremédiation (ODD 6, 14, 15) :} la \cle{bioremédiation} utilise des micro-organismes (bactéries \emph{Pseudomonas}, \emph{Alcanivibrio} ; champignons \emph{Phanerochaete}) pour dégrader les hydrocarbures (marées noires, sols pollués par le pétrole), les métaux lourds (bioaccumulation, biosorption, biominéralisation) ou les pesticides (dégradation enzymatique par hydrolases, oxydoréductases). Au Cameroun, des essais de biostimulation de sols contaminés par des déchets pétroliers à Kribi (bassin sédimentaire) ont montré une réduction de \SI{80}{\percent} des hydrocarbures totaux en 6 mois.
    \item \textbf{Économie circulaire et bioéconomie (ODD 7, 9, 12) :} la transformation de déchets agro-industriels (bagasse de canne, tourteaux de coton/soja/arachide, drêches de brasserie, écorces de cacao) en bioéthanol (2G), en enzymes industrielles (cellulases, xylanases, amylases), en biomatériaux (bioplastiques PHA, PLA) ou en protéines unicellulaires pour l'alimentation animale valorise des flux aujourd'hui perdus ou brûlés, réduisant l'empreinte carbone.
\end{itemize}

\begin{methode}[Discuter les enjeux d'une application biotechnologique -- Grille Probatoire]
Pour structurer une argumentation rigoureuse (épreuve écrite, oral, exposé), suivez ces 4 temps obligatoires :
\begin{enumerate}
    \item \textbf{Description technique du procédé :} nommez l'organisme producteur (espèce, souche), le gène ou l'enzyme d'intérêt, le type de biotechnologie (traditionnelle / moderne : ADN recombinant, édition de génome, fermentation de précision), le produit final (molécule, service) et l'échelle (artisanale, pilote, industrielle).
    \item \textbf{Bénéfices attendus (multi-critères) :} 
    \begin{itemize}
        \item \textit{Économiques :} emplois directs/indirects, revenus, substitution importations, balance commerciale, valeur ajoutée locale. Chiffrez (rendement, coût unitaire, volume marché).
        \item \textit{Sociaux/Santé :} accessibilité (prix, disponibilité), sécurité alimentaire, couverture vaccinale, lutte contre maladies négligées.
        \item \textit{Environnementaux :} réduction \ce{CO2}, économie circulaire, préservation biodiversité, moindre éco-toxicité.
    \end{itemize}
    \item \textbf{Risques, limites et incertitudes :} 
    \begin{itemize}
        \item \textit{Techniques :} rendement instable, contamination, coût de purification, échelle (scale-up).
        \item \textit{Économiques :} CAPEX/OPEX élevés, dépendance intrants/brevets, risque monopole, viabilité économique sans subvention.
        \item \textit{Éthiques/Juridiques :} brevetabilité, APA (Nagoya), équité d'accès, consentement communautés locales.
        \item \textit{Santé/Environnement :} dissémination OGM, transfert horizontal gènes, toxicologie, résistance (antibiotiques, insecticides).
    \end{itemize}
    \item \textbf{Conditions d'encadrement et conclusion argumentée :} 
    \begin{itemize}
        \item Cadre légal : avis CNB, évaluation risques (ERA), conformité Cartagena/Nagoya, code propriété intellectuelle (OAPI).
        \item Gouvernance : comités d'éthique (institutionnel, national), participation citoyenne, traçabilité (étiquetage OGM).
        \item Mesures techniques : confinement (L1-L3), \emph{kill-switch}, auxotrophie, surveillance post-mise sur le marché (pharmacovigilance, monitoring environnemental).
        \item \textbf{Conclusion nuancée :} « Sous réserve de [conditions], cette application est \textbf{souhaitable / à réorienter / à rejeter} car [argument principal pondéré]. »
    \end{itemize}
\end{enumerate}
\end{methode}

\begin{attention}[Pièges à éviter dans une discussion d'enjeux -- Probatoire]
\begin{itemize}
    \item \textbf{Position absolue :} ne pas écrire « c'est totalement bon » ou « c'est totalement mauvais ». L'exigence est une \cle{argumentation nuancée} : reconnaître les avantages \textbf{et} les risques, puis pondérer.
    \item \textbf{Confusion des types :} ne pas mélanger \cle{biotechnologie traditionnelle} (fermentation spontanée ou inoculée, sélection massale) et \cle{biotechnologie moderne} (OGM, enzymatique, SAM, thérapie génique) ; les enjeux éthiques, juridiques et techniques diffèrent radicalement.
    \item \textbf{Oubli des acteurs :} citer explicitement les \textbf{acteurs concernés} et leurs intérêts : producteurs (rural/urbain), industriels (PME/grands groupes), consommateurs/patients, État (régulateur, financeur), recherche publique/privée, ONG, communautés locales (détentrices savoirs traditionnels).
    \item \textbf{Négligence du cadre légal :} une discussion sans référence au \textbf{droit} (loi biosécurité, protocole Cartagena, protocole Nagoya, accord ADPIC/OAPI, code travail, code environnement) est incomplète et perd des points.
    \item \textbf{Approximations scientifiques :} « les OGM sont dangereux » (vague) vs « le risque de flux de gènes du maïs Bt vers le maïs paysan par pollinisation croisée est documenté et nécessite une distance d'isolement ou une stérilité mâle » (précis, scientifique).
\end{itemize}
\end{attention}

\begin{exoresolu}[Cas d'étude : Production d'insuline humaine recombinante par \emph{E. coli} K12 modifié au Cameroun]
\textbf{Énoncé :} Une firme pharmaceutique souhaite implanter au Cameroun une unité de production d'insuline humaine recombinante à partir d'une souche d'\emph{Escherichia coli} K12 portant le gène de la proinsuline humaine. Discutez les enjeux de ce projet pour le pays.

\textbf{Solution détaillée :}
\begin{enumerate}
    \item \textbf{Description du procédé :} Biotechnologie moderne (ADN recombinant, fermentation de précision). Le gène codant la proinsuline humaine (synthétisé \emph{in vitro}, codons optimisés pour \emph{E. coli}) est inséré dans un vecteur plasmidique à promoteur fort (ex. \emph{tac}, \emph{T7}), transformé dans \emph{E. coli} K12 (souche auxotrophe, non pathogène, \emph{recA-} pour stabilité). Fermentation en mode \emph{fed-batch} en cuve inox \SI{10000}{L} (milieu minéral défini, induction IPTG). La proinsuline s'accumule en \cle{corps d'inclusion} (inclusion bodies). Lyse cellulaire, solubilisation (urée/guanidine), repliage oxydatif (formation ponts disulfure A7-B7, A20-B19, A6-A11), clivage enzymatique (trypsine/carboxypeptidase B) $\rightarrow$ insuline mature. Purification chromatographique (échange d'ions, exclusion stérique, phase inverse). Formulation finale (zinc, phénol, m-crésol, pH neutre). Contrôle qualité Pharmacopée (HPLC, SDS-PAGE, puissance biologique, stérilité, endotoxines $<$ \SI{10}{UI/mg}).
    \item \textbf{Bénéfices attendus :} 
    \begin{itemize}
        \item \textit{Santé (ODD 3) :} insuline locale, coût de production réduit (\approx \SI{30}{\percent} moins cher que l'importé), disponibilité continue pour les \approx \num{600000} diabétiques camerounais (prévalence \approx \SI{6}{\percent}), réduction des ruptures de stock.
        \item \textit{Économie (ODD 8, 9) :} création de \approx 150 emplois directs hautement qualifiés (biologistes, ingénieurs procédés, QC/QA, maintenance), formation par transfert de technologie, stimulation sous-traitance locale (emballage, logistique chaîne de froid).
        \item \textit{Souveraineté (ODD 17) :} réduction facture importation médicaments (\approx \SI{5}{milliards} FCFA/an pour l'insuline), sécurité d'approvisionnement en cas de crise mondiale.
    \end{itemize}
    \item \textbf{Risques et limites :} 
    \begin{itemize}
        \item \textit{Techniques :} contamination bactérienne/phagique (arrêt production), instabilité plasmidique (perte gène), formation corps d'inclusion (rendement repliage \approx \SI{40-60}{\percent}), endotoxines LPS (pyrogènes) exigeant purification drastique.
        \item \textit{Économiques :} CAPEX initial \approx \SI{12}{milliards} FCFA (site stérile, fermenteurs, purification, QC, utilités). OPEX élevé (milieux, résines, consommables, énergie, expertise expatriée). Dépendance matières premières importées (milieux, résines, filtres, IPTG).
        \item \textit{Éthiques/Juridiques :} brevet sur procédé/souche (propriété firme). Risque monopole si pas d'accord licence volontaire ou clause prix préférentiel marché public. Nécessité APA (Nagoya) si souche hôte ou éléments génétiques (promoteur, signal peptide) proviennent biodiversité camerounaise.
        \item \textit{Environnementaux :} rejet éventuel \emph{E. coli} recombinant dans effluents. Confinement L2 obligatoire (décret 2003/006). Traitement eaux usées : stérilisation thermique continue (\SI{130}{\celsius}) + validation inactivation (absence UFC/100mL).
    \end{itemize}
    \item \textbf{Conditions d'encadrement et conclusion argumentée :} 
    \begin{itemize}
        \item Dossier complet soumis au CNB (MINRESI) : évaluation risques environnementaux/sanitaires, plan confinement, plan gestion urgences.
        \item Contrat État-Firme : clause transfert technologie effectif (formation 5 ans, accréditation ISO 9001/13485 site), clause prix plafond pour commande publique (CAMEBU), clause partage bénéfices (APA) si applicable.
        \item Plan surveillance post-commercialisation : pharmacovigilance (effets indésirables), monitoring environnemental effluents (PCR gène insuline, survie souche), audit annuel indépendant.
        \item Implication associations patients diabétiques (ex. Association Camerounaise des Diabétiques) dans suivi accessibilité/prix.
    \end{itemize}
    \textbf{Conclusion :} Ce projet est \textbf{souhaitable} car il répond à un besoin de santé publique majeur (insuline = médicament essentiel OMS) et crée de la valeur locale hautement qualifiée, \textbf{à condition} que l'État négocie fermement des clauses contraignantes de transfert de technologie, de prix équitable et de partage des bénéfices, et que la biosécurité (confinement L2, traitement effluents) soit rigoureusement appliquée et contrôlée par le CNB.
\end{enumerate}
\end{exoresolu}

\begin{popfigure}
\begin{tikzpicture}[scale=0.9, every node/.style={font=\small}]
    % Balance beam
    \draw[popdark, line width=2pt] (-5,0) -- (5,0);
    \draw[popdark, line width=2pt] (0,0) -- (0,1.5);
    \draw[popdark, line width=1.5pt] (-5,0) -- (-5,-0.5) -- (-4.5,-0.5);
    \draw[popdark, line width=1.5pt] (5,0) -- (5,-0.5) -- (4.5,-0.5);
    % Pans
    \draw[popblue, fill=popblueL, rounded corners=3pt] (-6.5,-0.8) rectangle (-3.5,-0.2) node[midway, align=center] {Avantages\\économiques};
    \draw[poporange, fill=poporangeL, rounded corners=3pt] (3.5,-0.8) rectangle (6.5,-0.2) node[midway, align=center] {Risques \&\\questions éthiques};
    % Support
    \draw[popdark, line width=3pt] (0,-0.5) -- (0,-1.5);
    \draw[popdark, line width=3pt] (-0.5,-1.5) -- (0.5,-1.5);
    % Labels
    \node[popdark, align=center] at (-5,1.2) {Emplois, revenus,\\souveraineté, énergie,\\santé accessible};
    \node[popdark, align=center] at (5,1.2) {Coût R\&D, dépendance,\\brevets, biodiversité,\\acceptabilité, monopole};
\end{tikzpicture}
\legende{Balance schématique des enjeux : avantages socio-économiques (gauche) vs risques techniques, économiques, éthiques et environnementaux (droite) d'une application biotechnologique. L'équilibre dépend de la gouvernance.}
\end{popfigure}

\begin{popfigure}
\begin{tikzpicture}[scale=0.85, every node/.style={font=\small},
    decision/.style={draw, rounded corners=4pt, fill=popblueL, minimum width=3.5cm, minimum height=1cm, align=center, text width=3.2cm},
    outcome/.style={draw, rounded corners=4pt, fill=popgreenL, minimum width=3.8cm, minimum height=1cm, align=center, text width=3.5cm},
    risk/.style={draw, rounded corners=4pt, fill=poporangeL, minimum width=3.8cm, minimum height=1cm, align=center, text width=3.5cm},
    frame/.style={draw, rounded corners=4pt, fill=poppurpleL, minimum width=4.5cm, minimum height=1cm, align=center, text width=4.2cm},
    arrow/.style={-Stealth, popdark, thick}
]
    \node[decision, align=center] (app) at (0,3){Application\\biotechnologique\\identifiée};
    \node[outcome, align=center] (ben) at (-4.5,1){Bénéfices attendus\\ (emplois, santé,\\environnement, souveraineté)};
    \node[risk, align=center] (ris) at (4.5,1){Risques potentiels\\ (coût, dépendance,\\biodiversité, éthique, monopole)};
    \node[frame, align=center] (enc) at (0,-1.2){Conditions d'encadrement\\ (Biosécurité CNB, APA Nagoya,\\Transfert tech., Prix équitable,\\Participation citoyenne)};
    \draw[arrow] (app.south west) -- (ben.north);
    \draw[arrow] (app.south east) -- (ris.north);
    \draw[arrow] (ben.south) -- (enc.north west);
    \draw[arrow] (ris.south) -- (enc.north east);
    \node[popdark, align=center, font=\bfseries\small] at (0,-2.8) {Décision argumentée :\\Déploiement conditionné / Réorientation / Abandon};
\end{tikzpicture}
\legende{Arbre de décision pour l'évaluation intégrée d'une application biotechnologique : de l'identification du procédé à la décision finale, en passant par l'analyse multicritères des bénéfices, des risques et des conditions juridico-techniques d'encadrement.}
\end{popfigure}

\begin{aretenir}[Synthèse des concepts clés pour le Probatoire]
\begin{itemize}
    \item \textbf{Biotechnologie :} toute application technique utilisant des organismes vivants ou leurs composants (enzymes, ADN) pour produire biens/services.
    \item \textbf{Traditionnelle vs Moderne :} Traditionnelle = fermentation spontanée/empirique, sélection massale/phénotypique. Moderne = ADN recombinant, édition génomique (CRISPR-Cas9), SAM, fermentation de précision, thérapie cellulaire/génique.
    \item \textbf{Biosécurité (Cartagena) :} prévention risques OGM/OMP. \textbf{APA (Nagoya) :} justice dans partage bénéfices ressources génétiques/savoirs traditionnels.
    \item \textbf{Enjeux :} Économiques (emplois, souveraineté, coût, dépendance), Éthiques (brevet vivant, équité, consentement), Sanitaires/Environnementaux (sécurité produit, confinement, flux gènes).
    \item \textbf{Argumentation :} Description $\rightarrow$ Bénéfices (chiffrés) $\rightarrow$ Risques (précis) $\rightarrow$ Encadrement (légal/technique) $\rightarrow$ Conclusion nuancée.
\end{itemize}
\end{aretenir}

\begin{exercice}[Entraînement Probatoire -- Sujet type]
\textbf{Sujet :} La société \og GreenEnergy Cameroon \fg{} propose d'installer une unité industrielle de production de bioéthanol de 2\up{ème} génération (2G) à partir de la bagasse de canne à sucre (résidu sucrerie SOSUCAM) et de paille de maïs, en utilisant un cocktail enzymatique (cellulases, hémicellulases) produit par un champignon filamenteux \emph{Trichoderma reesei} modifié par CRISPR-Cas9 pour sur-exprimer les gènes \emph{cbh1} et \emph{egl2}. Le bioéthanol sera mélangé à l'essence (E10).
\\
\textbf{Travail demandé :} En vous appuyant sur vos connaissances et la méthode vue en cours, discutez les enjeux de ce projet pour le Cameroun. Structurez votre réponse en 4 parties : 1. Description du procédé ; 2. Bénéfices attendus ; 3. Risques et limites ; 4. Conditions d'encadrement et conclusion argumentée.
\end{exercice}

\begin{corrige}[Exercice -- Bioéthanol 2G au Cameroun]
\textbf{1. Description du procédé :} Biotechnologie moderne (enzymatique + édition de génome). Substrats : bagasse (cellulose \approx 40\%, hémicellulose \approx 25\%, lignine \approx 20\%) et paille maïs. Prétraitement physico-chimique (vapeur explosion / acide dilué) $\rightarrow$ hydrolyse enzymatique (cocktail \emph{T. reesei} $\Delta$cre1 + \emph{cbh1/egl2} surexprimés) $\rightarrow$ fermentation \emph{S. cerevisiae} (tolérant inhibiteurs : furfural, HMF, acide acétique) $\rightarrow$ distillation/déshydratation $\rightarrow$ bioéthanol carburant (E10). OGM : \emph{T. reesei} modifié (CRISPR) $\rightarrow$ confinement L1/L2 requis pour production enzymes.
\textbf{2. Bénéfices :} \textit{Économiques :} valorisation déchets (bagasse brûlée \rightarrow énergie), création filière bioénergie, emplois ruraux (collecte paille), substitution importation essence (\approx \SI{300}{millions} L/an), recettes fiscales. \textit{Environnementaux (ODD 7, 13) :} réduction GES (\approx \SI{70}{\percent} vs essence), économie circulaire, pas de concurrence alimentaire (2G vs 1G maïs/canne). \textit{Sociaux :} revenus complémentaires planteurs maïs/canne.
\textbf{3. Risques :} \textit{Techniques :} récalcitrance lignine (rendement hydrolyse < 90\%), coût enzymes encore élevé (\approx \SI{0.10-0.30}{USD/L} éthanol), inhibiteurs fermentation, échelle industrielle non encore mature (usines 2G rares). \textit{Économiques :} CAPEX très élevé (\approx \SI{20-30}{milliards} FCFA), prix pétrole volatile (compétitivité), dépendance technologie/brevets (enzymes, levure, prétraitement). \textit{Éthiques/Juridiques :} OGM \emph{T. reesei} (CRISPR) $\rightarrow$ réglementation Cartagena/Nagoya (APA si souche parentale biodiversité locale). Accaparement résidus agricoles (paille) vs paillage/sol/organique $\rightarrow$ dégradation sols si exportation totale biomasse. \textit{Environnementaux :} fuites enzymes/OGM, consommation eau/énergie prétraitement.
\textbf{4. Encadrement \& Conclusion :} Avis CNB obligatoire (dossier ERA enzymes OGM). Étude d'impact environnemental et social (EIES) complète. Contrat : clause prix plancher planteurs résidus, clause maintien matière organique sols (retour cendres/vinasses), transfert tech enzymes (IRAD/Universités). Surveillance : émissions atmosphériques (COV, NOx), qualité eaux rejet, traçabilité OGM.
\textbf{Conclusion :} Projet \textbf{souhaitable stratégiquement} (souveraineté énergétique, climat, déchets) \textbf{sous réserve} de : viabilité économique prouvée (subvention transitoire / taxe carbone), garantie durabilité sols (taux restitution résidus), maîtrise technologique enzymes (partenariat public-privé), et cadre biosécurité strict pour l'OGM producteur d'enzymes.
\end{corrige}

En résumé, les biotechnologies offrent des leviers puissants pour le développement économique, la santé et la protection de l'environnement au Cameroun. Leur déploiement responsable exige une \cle{évaluation intégrée} — technique, économique, éthique, juridique — et une gouvernance inclusive associant État, secteur privé, recherche, société civile et communautés locales. C'est à cette condition qu'elles contribueront durablement aux Objectifs de Développement Durable.

\newpage

\coursec{Activité d'intégration}

\coursec{Application de la biotechnologie dans la production de biens et services}

\courssub{Objectifs de l'activité}

À l'issue de cette activité d'intégration, tu seras capable de :

\begin{itemize}
\item \cle{décrire} le principe d'un procédé biotechnologique appliqué à la production de biens ou de services, en identifiant l'agent biologique, le substrat, les conditions opératoires et le produit obtenu ;
\item \cle{comparer} une biotechnologie traditionnelle (fermentation lactique) et une biotechnologie de traitement des déchets (méthanisation), en mobilisant les critères vus dans la leçon ;
\item \cle{discuter} les enjeux économiques, sanitaires et éthiques d'une application biotechnologique concrète ;
\item \cle{rédiger} un avis argumenté destiné à un décideur (la direction du lycée), compétence clé de l'approche par les compétences.
\end{itemize}

\courssub{Situation}

Le lycée de Bafoussam (région de l'Ouest) compte \num{1200} élèves. Son club environnement, encadré par le professeur de SVTEEHB, souhaite valoriser les ressources locales et réduire les dépenses de l'établissement. Deux projets sont en compétition pour obtenir le budget annuel du club, fixé à \num{150000} FCFA.

\textbf{Projet 1 : production de yaourt.} La cantine achète chaque semaine \SI{60}{L} de lait frais auprès des éleveurs de la plaine des Mbo, à \SI{600}{FCFA} le litre. Le club propose de transformer une partie de ce lait en yaourt à l'aide de ferments lactiques (sachet de ferments contenant \emph{Lactobacillus delbrueckii} subsp. \emph{bulgaricus} et \emph{Streptococcus thermophilus}, coûtant \SI{500}{FCFA} et permettant d'ensemencer \SI{10}{L} de lait). Le protocole est le suivant : pasteuriser le lait (\SI{85}{\celsius} pendant 15 minutes), le refroidir à \SI{45}{\celsius}, ajouter les ferments, puis incuber à \SI{42}{\celsius}--\SI{44}{\celsius} pendant 4 à 6 heures dans des pots fermés. Le yaourt obtenu serait vendu \SI{250}{FCFA} le pot de \SI{125}{mL} aux élèves, alors que le yaourt industriel coûte \SI{350}{FCFA} en ville.

\textbf{Projet 2 : production de biogaz.} La cantine produit chaque jour environ \SI{40}{kg} de déchets organiques (épluchures de manioc et de plantain, restes de repas, fumier du poulailler scolaire). Ces déchets sont actuellement brûlés ou abandonnés, ce qui pollue l'air et attire les mouches. Le club propose d'installer un petit digesteur (fosse étanche de \SI{2}{m^3}) dans lequel les déchets, mélangés à l'eau, fermentent en l'absence d'oxygène grâce à une communauté de micro-organismes anaérobies (bactéries hydrolytiques, acidogènes, acétogènes et archées méthanogènes). Le biogaz produit (environ \SI{60}{\percent} de méthane \ce{CH4} et \SI{40}{\percent} de dioxyde de carbone \ce{CO2}) alimenterait un réchaud de la cantine. On estime la production à \SI{1,5}{m^3} de biogaz par jour, soit l'équivalent de \SI{3}{kg} de bois de chauffe économisés quotidiennement ; le digestat résiduel servirait d'engrais organique pour le jardin pédagogique. Le coût d'installation du digesteur est estimé à \SI{120000}{FCFA}.

La direction du lycée hésite : elle ne financera qu'un seul projet cette année. Le club environnement te sollicite, toi et ton groupe, pour analyser scientifiquement les deux projets et formuler une recommandation argumentée.

\courssub{Consignes}

Par groupes de 4 à 6 élèves, traitez les tâches suivantes dans l'ordre, en rédigeant soigneusement vos réponses sur une fiche de synthèse.

\begin{enumerate}
\item Pour \textbf{chaque projet}, identifiez et notez dans un tableau : l'\cle{agent biologique} mis en œuvre, le \cle{substrat} (matière première transformée), les \cle{conditions opératoires} essentielles (température, présence ou absence d'oxygène, durée), le \cle{produit} obtenu et le \cle{service} rendu.
\item Écrivez l'équation-bilan simplifiée de la transformation réalisée dans chaque procédé (conversion du lactose pour le projet 1 ; conversion de la matière organique pour le projet 2).
\item Construisez le \textbf{schéma de principe} de chaque procédé : une chaîne d'étapes fléchées allant du substrat au produit, en précisant à chaque étape le rôle des micro-organismes et les conditions imposées.
\item Classez chaque projet dans la catégorie « biotechnologie \cle{traditionnelle} » ou « biotechnologie \cle{moderne} », en justifiant votre classement à l'aide des critères de la leçon (ancienneté du procédé, maîtrise des souches, niveau de technicité). Que remarquez-vous ?
\item Dressez un tableau comparatif des \textbf{bénéfices économiques} (gains, économies, coûts) et des \textbf{questions sanitaires et éthiques} soulevées par chaque projet (hygiène, sécurité, acceptabilité, impact environnemental).
\item Rédigez un \textbf{avis argumenté d'environ 10 lignes} adressé à la direction du lycée, recommandant l'un des deux projets. Votre avis doit s'appuyer sur au moins deux arguments scientifiques et deux arguments économiques ou sanitaires, et mentionner une limite ou un risque du projet choisi ainsi qu'une mesure pour le maîtriser.
\item Restitution : chaque groupe présente son avis en 3 minutes ; un débat collectif permet à la classe de voter le projet retenu, en explicitant les critères du choix.
\end{enumerate}

\courssub{Ressources à mobiliser}

\begin{itemize}
\item La définition de la \cle{biotechnologie} : toute application technologique qui utilise des organismes vivants (ou leurs constituants : cellules, enzymes) pour produire des biens ou des services.
\item Les caractéristiques des biotechnologies \cle{traditionnelles} : procédés anciens, empiriques à l'origine, reposant sur la fermentation par des micro-organismes (bactéries, levures), technicité simple, faible coût.
\item Les caractéristiques des biotechnologies \cle{modernes} : maîtrise des souches (sélection assistée), utilisation d'enzymes purifiées, contrôle précis des conditions, équipements plus élaborés.
\item La fermentation \cle{lactique} : conversion du lactose en acide lactique par des bactéries lactiques, en anaérobiose, avec acidification (chute du pH vers 4,5) qui coagule les protéines du lait (caséines) et empêche le développement des germes pathogènes.
\item La fermentation \cle{méthanique} (méthanisation) : dégradation anaérobie de la matière organique en quatre étapes (hydrolyse, acidogenèse, acétogenèse, méthanogenèse) par un consortium de micro-organismes, aboutissant au biogaz (\ce{CH4} + \ce{CO2}) et à un digestat fertilisant. Les archées méthanogènes sont strictement anaérobies.
\item Les enjeux : économiques (coûts, revenus, emplois), sanitaires (hygiène, pasteurisation, sécurité du gaz), environnementaux (valorisation des déchets, réduction de la déforestation) et éthiques (acceptabilité sociale, encadrement des pratiques).
\end{itemize}

\courssub{Démarche et corrigé commenté}

\begin{exoresolu}[Corrigé commenté de l'activité d'intégration]
\textbf{Étape 1 — Identification des composantes de chaque procédé.}

\begin{center}
\begin{tabularx}{\linewidth}{|l|X|X|}\hline
\rowcolor{popblueL} \textbf{Composante} & \textbf{Projet 1 : yaourt} & \textbf{Projet 2 : biogaz} \\ \hline
Agent biologique & Bactéries lactiques : \emph{L. bulgaricus} et \emph{S. thermophilus} (ferments sélectionnés, symbiose) & Consortium anaérobie : bactéries hydrolytiques, acidogènes, acétogènes ; archées méthanogènes \\ \hline
Substrat & Lait frais (lactose, protéines, lipides) & Déchets organiques (cellulose, hémicellulose, protéines, lipides) + eau \\ \hline
Conditions & Anaérobiose, \SI{42}{\celsius}--\SI{44}{\celsius}, 4--6 h, pots fermés, pH final \textasciitilde 4,5 & Anaérobiose stricte, digesteur étanche, température mésophile (ambiante à \SI{35}{\celsius}), temps de séjour \textasciitilde 20--30 j \\ \hline
Produit & Yaourt (gel acide, acide lactique, arômes) & Biogaz (\ce{CH4} \SI{60}{\percent}, \ce{CO2} \SI{40}{\percent}) + digestat (engrais organique) \\ \hline
Service rendu & Alimentation saine, locale, moins chère ; revenu pour le club & Énergie de cuisson propre, gestion assainie des déchets, engrais gratuit, lutte anti-déforestation \\ \hline
\end{tabularx}
\end{center}

\tcblower

\textbf{Étape 2 — Équations-bilans simplifiées.}

Fermentation lactique (conversion du lactose, disaccharide \ce{C12H22O11}) :
\[ \ce{C12H22O11 + H2O -> 4 CH3CHOHCOOH} \quad (\text{lactose} \rightarrow \text{4 acides lactiques}) \]
\emph{Note :} L'énergie (ATP) est produite pour la bactérie mais n'apparaît pas dans l'équation chimique globale du substrat.

Fermentation méthanique (schéma global de dégradation de la matière organique, assimilée à un glucide \ce{C6H12O6}) :
\[ \ce{C6H12O6 -> 3CH4 + 3CO2} \]
\emph{Note :} Cette équation masque les 4 étapes successives (hydrolyse, acidogenèse, acétogenèse, méthanogenèse) et l'intervention de plusieurs groupes microbiens distincts.

\textbf{Étape 3 — Schémas de principe.}

\begin{popfigure}
\begin{center}
\begin{tikzpicture}[
box/.style={draw=popblue, fill=popblueL, rounded corners, align=center, minimum height=0.9cm, font=\small, text width=2.8cm},
boxg/.style={draw=popgreen!70!black, fill=popgreenL, rounded corners, align=center, minimum height=0.9cm, font=\small, text width=2.8cm},
fleche/.style={-{Stealth[length=2.5mm]}, thick, popdark},
labelstyle/.style={font=\bfseries\small, left=0.3cm}]
% Projet 1
\node[box, align=center] (a1) at (0,0){Lait frais\\ (substrat)};
\node[box, right=0.8cm of a1, align=center] (a2){Pasteurisation\\ \SI{85}{\celsius}, 15 min\\ (destruction germes)}; 
\node[box, right=0.8cm of a2, align=center] (a3){Refroidissement \SI{45}{\celsius}\\ + Ensemencement\\ (ferments sélectionnés)};
\node[box, right=0.8cm of a3, align=center] (a4){Incubation\\ \SIrange{42}{44}{\celsius}, 4--6 h\\ Anaérobiose, acidification};
\node[box, right=0.8cm of a4, align=center] (a5){Yaourt\\ (produit : gel, pH 4,5)};
\draw[fleche] (a1) -- (a2); \draw[fleche] (a2) -- (a3); \draw[fleche] (a3) -- (a4); \draw[fleche] (a4) -- (a5);
% Projet 2
\node[boxg, below=2.5cm of a1, align=center] (b1){Déchets organiques\\ + eau (substrat)};
\node[boxg, right=0.8cm of b1, align=center] (b2){Broyage, mélange\\ Introduction digesteur\\ (mise en anaérobiose)};
\node[boxg, right=0.8cm of b2, align=center] (b3){Hydrolyse \& Acidogenèse\\ Bactéries hydrolytiques/acidogènes\\ Polymères $\rightarrow$ Acides gras volatils};
\node[boxg, right=0.8cm of b3, align=center] (b4){Acétogenèse \& Méthanogenèse\\ Bactéries acétogènes + Archées méthanogènes\\ Acétate/H2/CO2 $\rightarrow$ CH4};
\node[boxg, right=0.8cm of b4, align=center] (b5){Biogaz (réchaud)\\ + Digestat (jardin)};
\draw[fleche] (b1) -- (b2); \draw[fleche] (b2) -- (b3); \draw[fleche] (b3) -- (b4); \draw[fleche] (b4) -- (b5);
\node[labelstyle, popblue] at (a1.west) {Projet 1};
\node[labelstyle, popgreen!60!black] at (b1.west) {Projet 2};
\end{tikzpicture}
\end{center}
\legende{Chaînes de principe des deux procédés biotechnologiques. Projet 1 (bleu) : fermentation lactique en batch (discontinu) avec souches pures sélectionnées. Projet 2 (vert) : méthanisation continue en digesteur étanche avec consortium microbien complexe (bactéries + archées).}
\end{popfigure}

\textbf{Étape 4 — Classement et analyse critique.}

\begin{aretenir}[Classement des projets]
Le \textbf{Projet 1 (yaourt)} relève d'une \textbf{biotechnologie traditionnelle améliorée} (ou « moderne » au sens scolaire). \emph{Justification :} La fermentation lactique est millénaire (lait caillé, \emph{kindirmou} au Cameroun), mais l'usage de \textbf{souches pures sélectionnées} (\emph{L. bulgaricus} + \emph{S. thermophilus}), la \textbf{pasteurisation} systématique et le \textbf{contrôle de la température} d'incubation traduisent une maîtrise scientifique et technologique moderne.

Le \textbf{Projet 2 (biogaz)} relève d'une \textbf{biotechnologie moderne} appliquée à un phénomène naturel ancien. \emph{Justification :} La digestion anaérobie se produit naturellement (marais, panse des ruminants, lac Nyos), mais le \textbf{digesteur étanche}, la \textbf{récupération contrôlée du méthane}, la \textbf{gestion du digestat} et le \textbf{dimensionnement} du procédé pour un besoin énergétique précis constituent une ingénierie moderne.

\emph{Remarque fondamentale :} Les deux procédés reposent sur le \textbf{même fondement biologique} : la fermentation anaérobie par des micro-organismes. La frontière « traditionnel / moderne » ne tient pas au principe métabolique, mais au \textbf{niveau de maîtrise} : pureté des souches, contrôle des paramètres physico-chimiques (T°, pH, anaérobiose), automatisation, traçabilité.
\end{aretenir}

\textbf{Étape 5 — Tableau des enjeux (économiques, sanitaires, éthiques, environnementaux).}

\begin{center}
\begin{tabularx}{\linewidth}{|l|X|X|}\hline
\rowcolor{popblueL} \textbf{Critère} & \textbf{Projet 1 : Yaourt} & \textbf{Projet 2 : Biogaz} \\ \hline
\textbf{Investissement initial} & Faible (\textasciitilde \SI{30000}{FCFA} : marmites, thermomètre, pots, premiers ferments) & Élevé (\SI{120000}{FCFA} : maçonnerie digesteur, tuyauterie, réchaud) \\ \hline
\textbf{Coûts récurrents} & Achat lait (\SI{36000}{FCFA}/sem) + ferments (\SI{3000}{FCFA}/sem) + énergie (bois/gaz pasteurisation) & Eau, main-d'œuvre chargement, maintenance joints/vannes (faible) \\ \hline
\textbf{Revenus / Économies} & Vente \textasciitilde 480 pots/sem $\times$ \SI{250}{FCFA} = \SI{120000}{FCFA}/sem. Marge nette \textasciitilde \SI{80000}{FCFA}/sem. & Économie bois : \SI{3}{kg}/j $\times$ \SI{150}{FCFA/kg} $\times$ 30 j = \SI{13500}{FCFA}/mois. + Engrais (valeur \textasciitilde \SI{5000}{FCFA}/mois). Amortissement \textasciitilde 18--24 mois. \\ \hline
\textbf{Risques sanitaires} & \textbf{Majeurs si défaillance :} Pasteurisation insuffisante (brucellose, salmonellose), contamination post-incubation (moisissures, \emph{Staphylococcus}), rupture chaîne du froid. & Fuite biogaz (inflammation, asphyxie), \ce{H2S} toxique (traces), manipulation digestat (pathogènes résiduels si T° insuffisante). \\ \hline
\textbf{Impact environnemental} & Valorisation lait local, circuit court. Rejet de petit-lait (polluant si non traité). & \textbf{Très positif :} Suppression brûlage déchets (fumées, \ce{CO2}, \ce{CH4} atmosphérique), substitution bois (lutte déforestation), retour au sol (économie circulaire). \\ \hline
\textbf{Acceptabilité / Éthique} & Élevée (produit culturellement ancré, apprécié). Risque : concurrence déloyale artisans locaux ? & À construire : préjugés « gaz d'égout », peur explosion. Nécessite sensibilisation. Éthique : gestion déchets collective vs individuelle. \\ \hline
\end{tabularx}
\end{center}

\textbf{Étape 6 — Exemple d'avis argumenté (modèle pour la rédaction).}

\begin{exemplebox}[Avis argumenté adressé à la direction du lycée]
\textbf{Objet : Recommandation du club environnement — Choix du projet biogaz.}

Monsieur le Proviseur,

Le club environnement recommande à la direction le \textbf{projet de production de biogaz (Projet 2)}.

\textbf{Arguments scientifiques :} (1) La méthanisation exploite un consortium microbien anaérobie (bactéries hydrolytiques/acidogènes/acétogènes + archées méthanogènes) pour dégrader \SI{40}{kg/j} de déchets organiques en biogaz (\ce{CH4}/\ce{CO2}) et digestat, réalisant une \textbf{double valorisation matière-énergie} que ne permet pas le yaourt. (2) Le procédé supprime le brûlage des déchets à l'air libre, source avérée de polluants atmosphériques (particules fines, HAP, \ce{CH4}) et de prolifération de vecteurs (mouches).

\textbf{Arguments économiques et sanitaires :} (1) Malgré un investissement initial de \SI{120000}{FCFA} (dans le budget \SI{150000}{FCFA}), l'économie de \SI{3}{kg} de bois/jour (\textasciitilde \SI{13500}{FCFA}/mois) et la production d'engrais gratuit assurent un amortissement sous 2 ans, puis un gain net récurrent, sans achat de matière première (contrairement au lait/ferments). (2) Sanitairement, le projet élimine un foyer de pollution et de nuisance immédiat pour la communauté scolaire.

\textbf{Limite et mesure de maîtrise :} Le risque principal est la \textbf{manipulation d'un gaz inflammable} (\ce{CH4}) en milieu scolaire. \textbf{Mesure :} Installation du digesteur à \textgreater \SI{10}{m} des bâtiments, formation de deux agents d'entretien aux consignes de sécurité (détection odeur, contrôle joints hebdomadaire, interdiction flamme nue), pose d'un détendeur et brûleur adaptés.

Le projet yaourt, pertinent pédagogiquement et rentable, pourra être autofinancé l'an prochain par les bénéfices de la cantine ou une micro-subvention parents.
\end{exemplebox}
\end{exoresolu}

\begin{attention}[Pièges fréquents dans ce type d'activité]
\begin{itemize}
\item Confondre \textbf{substrat} (ce qui est transformé : le lait, les déchets) et \textbf{agent biologique} (ce qui transforme : les micro-organismes).
\item Affirmer que la fermentation lactique consomme de l'oxygène : c'est une voie \textbf{strictement anaérobie} (tolérance à l'O2 pour certaines souches, mais pas respiration).
\item Écrire l'équation de la fermentation lactique sans l'eau : $\ce{C12H22O11 -> 4 C3H6O3}$ est faux (bilan H/O non respecté) ; il faut $\ce{C12H22O11 + H2O -> 4 C3H6O3}$.
\item Classer le biogaz parmi les biotechnologies « purement modernes » sans nuance : le principe (fermentation) est traditionnel, seule la maîtrise technologique (digesteur, souches) est moderne.
\item Oublier que les \textbf{archées méthanogènes} (et non des bactéries) réalisent l'étape finale de production de \ce{CH4}.
\item Rédiger un avis sans aucun chiffre : un argument économique crédible s'appuie \textbf{obligatoirement} sur les données de l'énoncé.
\item Négliger le \textbf{digestat} comme produit à part entière (engrais) et non simple déchet.
\end{itemize}
\end{attention}

\courssub{Bilan de l'activité}

Cette activité t'a permis de mettre en évidence plusieurs idées essentielles de la leçon NCT13 :

\begin{itemize}
\item Tout procédé biotechnologique se décrit par la même grille d'analyse : \cle{agent biologique} $\rightarrow$ \cle{substrat} $\rightarrow$ \cle{conditions opératoires contrôlées} $\rightarrow$ \cle{produit(s)} et \cle{service(s) rendu(s)}.
\item Deux applications en apparence très différentes (alimentation humaine vs énergie/assainissement) reposent sur le \textbf{même fondement biologique universel} : la fermentation anaérobie par des micro-organismes.
\item La distinction entre biotechnologie \cle{traditionnelle} et \cle{moderne} ne porte pas sur le principe métabolique, mais sur le \textbf{niveau de maîtrise} : sélection et pureté des souches, contrôle précis des paramètres (T°, pH, anaérobiose), ingénierie du réacteur (digesteur, yaourtière), traçabilité.
\item Le choix d'une biotechnologie ne se réduit jamais à la science : il intègre indissociablement des \cle{enjeux économiques} (CAPEX/OPEX, amortissement, revenu), \cle{sanitaires} (hygiène, sécurité, toxicologie), \cle{environnementaux} (bilan carbone, gestion déchets, biodiversité) et \cle{éthiques} (acceptabilité sociale, justice environnementale, propriété du vivant).
\item Enfin, tu as exercé une compétence transversale précieuse pour le Probatoire et la vie professionnelle : transformer une analyse scientifique rigoureuse en \textbf{recommandation argumentée, chiffrée et nuancée} pour un décideur, comme le font les experts dans les filières agroalimentaire (IRAD, SODEPA), énergétique (SONEL, ARSEL) et environnementale (MINEPDED) du Cameroun.
\end{itemize}

\begin{savaistu}[Biotechnologies au Cameroun : exemples concrets]
\begin{itemize}
\item \textbf{Agroalimentaire :} La production industrielle de \emph{vin de palme} (raffia, palmier à huile) et de \emph{bière traditionnelle} (bil-bil, amont) repose sur des fermentations lactiques et alcooliques traditionnelles. Les brasseries modernes (UCB, Brasseries du Cameroun) utilisent des biotechnologies modernes : levures sélectionnées, fermentation contrôlée en cuves inox, enzymologie (amylases, glucanases) pour le maltage.
\item \textbf{Santé :} Le Cameroun produit des vaccins (fièvre jaune, méningite) à l'Institut Pasteur de Yaoundé (biotechnologie moderne : culture cellulaire, purification). Les tests rapides VIH/paludisme (bandelettes immunochromatographiques) utilisent des anticorps monoclonaux (biotechnologie moderne).
\item \textbf{Environnement/Energie :} Le projet \textbf{Lac Nyos} (dégazage par tuyaux auto-amorçants) est une prouesse d'ingénierie géologique/biologique. Des projets de \textbf{methanisation} existent à la SOCAPALM (déchets palmiers) et dans certains élevages modernes (porcheries de l'Ouest) pour l'autonomie énergétique.
\item \textbf{Agriculture :} L'IRAD utilise la \textbf{sélection assistée par marqueurs moléculaires} (biotechnologie moderne) pour améliorer le manioc (résistance mosaïque), le maïs (striga) et le cacao. La \textbf{micropropagation in vitro} (bananier, plantain, ananas) fournit des plants sains aux paysans.
\end{itemize}
\end{savaistu}

\begin{methode}[Rédiger un avis argumenté pour un décideur (compétence transversale)]
\begin{enumerate}
\item \textbf{Objet clair :} Projet recommandé + Destinataire.
\item \textbf{Arguments scientifiques (2 min) :} Mécanisme biologique, efficacité du procédé, innovation technique. Utiliser le vocabulaire précis (agent, substrat, conditions, rendement).
\item \textbf{Arguments économiques/sanitaires/environnementaux (2 min) :} Chiffres clés de l'énoncé (coûts, gains, économies, quantités). Comparer au statu quo.
\item \textbf{Risque majeur + Mesure de maîtrise (1 min) :} Honnêteté intellectuelle. Identifier le point faible (sécurité, coût, acceptabilité) et proposer une solution concrète, réaliste, locale.
\item \textbf{Conclusion ouverte :} Évoquer la complémentarité des projets (l'un n'exclut pas l'autre à terme).
\item \textbf{Style :} Professionnel, respectueux, synthétique (\textasciitilde 10 lignes), sans familiarité excessive.
\end{enumerate}
\end{methode}

\begin{exercice}[Entraînement Probatoire - Type QCM / Réponse courte]
\begin{enumerate}
\item On souhaite produire de l'éthanol de biocarburant à partir de résidus de canne à sucre (bagasse) par fermentation. Précisez : (a) L'agent biologique principal. (b) Le substrat initial et le prétraitement nécessaire. (c) L'équation bilan de la fermentation alcoolique. (d) Pourquoi ce procédé est-il classé en biotechnologie moderne ?
\item \textbf{Vrai ou Faux ?} Justifiez brièvement.
\begin{itemize}
\item[a)] La méthanisation est un procédé aérobie.
\item[b)] Les archées méthanogènes sont des bactéries.
\item[c)] La pasteurisation du lait vise à détruire les ferments lactiques.
\item[d)] Le digestat de méthanisation est un amendement organique valorisable.
\item[e)] La frontière entre biotechnologie traditionnelle et moderne est infranchissable.
\end{itemize}
\item Un élève écrit : « Le yaourt est une biotechnologie traditionnelle car on fait ça depuis longtemps au village. » Corrigez cette affirmation en utilisant les critères du cours.
\end{enumerate}
\end{exercice}

\begin{corrige}[Exercice Entraînement Probatoire]
\begin{enumerate}
\item (a) \emph{Saccharomyces cerevisiae} (levure de boulanger/brasseur), souche sélectionnée pour haut rendement éthanol et tolérance éthanol. (b) Substrat : bagasse (cellulose, hémicellulose, lignine). Prétraitement indispensable : broyage + traitement chimique (acide/dilué) ou enzymatique (cellulases) pour hydrolyser les polymères en sucres fermentescibles (glucose, xylose). (c) $\ce{C6H12O6 -> 2 C2H5OH + 2 CO2}$. (d) Utilisation d'enzymes purifiées (cellulases), souche génétiquement améliorée ou sélectionnée, bioréacteur contrôlé (pH, T°, agitation, anaérobiose stricte), distillation pour purification.
\item a) \textbf{Faux.} Anaérobie stricte (4 étapes sans O2). b) \textbf{Faux.} Les archées forment un domaine distinct des bactéries (procaryotes mais ribosomes, membrane, paroi différents). c) \textbf{Faux.} La pasteurisation détruit la flore contaminante (pathogènes, altérants) ; on \textbf{ajoute} les ferments lactiques \textbf{après} refroidissement. d) \textbf{Vrai.} Riche en azote/minéraux, stabilisé, hygiénisé (si T° suffisante). e) \textbf{Faux.} Continuum : traditional $\rightarrow$ améliorée $\rightarrow$ moderne. Même procédé, maîtrise croissante.
\item Correction : « Le yaourt \emph{villageois} (lait cru, ensemencement naturel, T° ambiante) est une biotechnologie traditionnelle. Le yaourt \emph{du projet} (lait pasteurisé, ferments purs sélectionnés \emph{L. bulgaricus}/\emph{S. thermophilus}, incubation contrôlée 42-44°C) est une biotechnologie traditionnelle \textbf{améliorée / moderne} car il y a maîtrise des souches et des conditions opératoires. »
\end{enumerate}
\end{corrige}

\newpage

\coursec{Méthodes et exercices résolus}

\coursec{Application de la biotechnologie dans la production de biens et services}

\courssub{Fiche méthode 1 : Décrire le principe d'un procédé biotechnologique}

\begin{methode}[Décrire le principe d'un procédé biotechnologique]
Pour décrire correctement le principe d'un procédé biotechnologique (fermentation, production d'enzyme, sélection assistée...), suis toujours les mêmes étapes :
\begin{enumerate}
\item \textbf{Identifier l'agent biologique utilisé} : micro-organisme (levure, bactérie, moisissure), cellule, enzyme isolée, ou plante/animal sélectionné. Nomme-le précisément (ex. : \emph{Saccharomyces cerevisiae}, bactéries lactiques).
\item \textbf{Préciser la matière première (substrat)} : jus de palme, lait, farine de manioc, moût de maïs...
\item \textbf{Décrire la transformation réalisée} : quelle réaction biochimique ? Quel produit obtient-on ? Si possible, écris l'équation bilan (ex. : \ce{C6H12O6 -> 2C2H5OH + 2CO2} pour la fermentation alcoolique).
\item \textbf{Indiquer les conditions du procédé} : anaérobiose ou aérobiose, température, stérilisation, durée.
\item \textbf{Nommer le bien ou le service obtenu} et son utilité (aliment, médicament, dépollution...).
\end{enumerate}
\textbf{Réflexe Probatoire} : une description complète = agent + substrat + transformation + conditions + produit. Il manque un élément, il manque des points.
\end{methode}

\begin{exoresolu}[Principe de la fabrication du vin de palme]
Le vin de palme, boisson très consommée dans les régions du Littoral et du Sud-Ouest, est obtenu à partir de la sève de palmier (\emph{Elaeis guineensis}) laissée en récipient fermé. Au bout de quelques heures, le liquide devient trouble, mousseux et acquiert un goût alcoolisé.

Décris le principe du procédé biotechnologique mis en jeu, en précisant l'agent, le substrat, la transformation et les conditions.
\tcblower
\textbf{Solution détaillée}

\textbf{Étape 1 : l'agent biologique.} La sève de palme contient naturellement des levures, notamment \emph{Saccharomyces cerevisiae}, présentes sur l'écorce et dans l'air. Ce sont des micro-organismes unicellulaires eucaryotes.

\textbf{Étape 2 : le substrat.} La sève est riche en sucres (saccharose, glucose), qui constituent le substrat de la fermentation.

\textbf{Étape 3 : la transformation.} Les levures réalisent une \cle{fermentation alcoolique} : en l'absence de dioxygène, elles dégradent le glucose pour produire de l'énergie (ATP), avec formation d'éthanol et de dioxyde de carbone :
\[ \ce{C6H12O6 -> 2C2H5OH + 2CO2} \]
Le dégagement de \ce{CO2} explique la mousse observée ; l'éthanol explique le goût alcoolisé.

\textbf{Étape 4 : les conditions.} Le procédé se déroule en \cle{anaérobiose} (récipient fermé limitant l'apport d'air), à température ambiante (environ \SI{25}{\celsius} à \SI{30}{\celsius}), sans stérilisation préalable : c'est une fermentation spontanée.

\textbf{Étape 5 : le produit.} On obtient le vin de palme, boisson alcoolisée traditionnelle.

\textbf{Conclusion :} la fabrication du vin de palme est une biotechnologie traditionnelle dont le principe est la fermentation alcoolique des sucres de la sève de palme par des levures, en anaérobiose.
\end{exoresolu}

\courssub{Fiche méthode 2 : Comparer une biotechnologie traditionnelle et une biotechnologie moderne}

\begin{methode}[Comparer deux biotechnologies]
Une comparaison rigoureuse se fait \textbf{critère par critère}, jamais en deux paragraphes séparés. Construis un tableau avec les critères suivants :
\begin{enumerate}
\item \textbf{Ancienneté et origine} : pratique empirique ancestrale vs technique de laboratoire récente.
\item \textbf{Niveau de maîtrise de l'agent biologique} : micro-organismes spontanés, non identifiés vs souches sélectionnées, enzymes purifiées, gènes identifiés.
\item \textbf{Contrôle des conditions} : conditions ambiantes, reproductibilité faible vs paramètres contrôlés (température, pH, stérilité) en bioréacteur.
\item \textbf{Rendement et qualité} : production artisanale variable vs production industrielle standardisée.
\item \textbf{Coût et accessibilité} : faible coût, accessible aux populations rurales vs investissements lourds.
\item \textbf{Exemples} : fermentation du lait caillé, du vin de palme, du « nkui »... vs production d'insuline, sélection assistée par marqueurs moléculaires, enzymes industrielles.
\end{enumerate}
\textbf{Réflexe Probatoire} : conclus toujours en rappelant le point commun : les deux exploitent des \cle{êtres vivants ou leurs constituants} pour produire des biens ou services.
\end{methode}

\begin{exoresolu}[Comparaison fermentation traditionnelle / production moderne d'enzymes]
Au Cameroun, les femmes transforment le lait en lait caillé par fermentation spontanée. Par ailleurs, des industries produisent de la présure microbienne (enzyme coagulante) en bioréacteurs pour la fabrication fromagère.

Compare ces deux biotechnologies selon quatre critères précis, puis dégage leur point commun.
\tcblower
\textbf{Solution détaillée}

On construit le tableau comparatif critère par critère :

\begin{center}
\begin{tabularx}{\linewidth}{|l|X|X|}
\hline
\rowcolor{popblueL} \textbf{Critère} & \textbf{Lait caillé traditionnel} & \textbf{Présure microbienne industrielle} \\
\hline
Agent biologique & Bactéries lactiques spontanées, non identifiées, flore variable & Souche microbienne sélectionnée, enzyme purifiée et caractérisée \\
\hline
Conditions du procédé & Température ambiante, récipient non stérile, reproductibilité faible & Bioréacteur stérile, température et pH contrôlés, reproductibilité élevée \\
\hline
Rendement et qualité & Production artisanale, qualité variable selon les jours & Production industrielle, qualité standardisée et constante \\
\hline
Coût et accessibilité & Très faible coût, accessible en milieu rural & Investissements lourds, produit destiné à l'industrie \\
\hline
\end{tabularx}
\end{center}

\textbf{Justification des différences :} la biotechnologie traditionnelle repose sur un savoir-faire empirique transmis de génération en génération, sans connaissance précise des micro-organismes en jeu. La biotechnologie moderne s'appuie sur la connaissance scientifique des souches et des enzymes, ce qui permet de contrôler finement le procédé et d'augmenter les rendements.

\textbf{Point commun :} dans les deux cas, on exploite l'activité d'\cle{êtres vivants} (bactéries lactiques, champignons producteurs d'enzymes) ou de leurs constituants (enzymes) pour transformer une matière première (le lait) en un bien utile (lait caillé, fromage).

\textbf{Conclusion :} les deux procédés relèvent de la biotechnologie ; ils diffèrent par le degré de maîtrise scientifique et technique, non par le principe fondamental.
\end{exoresolu}

\courssub{Fiche méthode 3 : Expliquer le principe d'une fermentation donnée}

\begin{methode}[Raisonner sur une fermentation]
Face à un exercice portant sur une fermentation, applique la démarche suivante :
\begin{enumerate}
\item \textbf{Identifier le type de fermentation} : alcoolique (levures $\to$ éthanol + \ce{CO2}), lactique (bactéries lactiques $\to$ acide lactique), acétique (bactéries acétiques $\to$ vinaigre).
\item \textbf{Écrire l'équation bilan} si elle est connue :
\begin{itemize}
\item Fermentation alcoolique : \ce{C6H12O6 -> 2C2H5OH + 2CO2}
\item Fermentation lactique : \ce{C6H12O6 -> 2C3H6O3}
\end{itemize}
\item \textbf{Rappeler le rôle énergétique} : la fermentation est une voie de production d'ATP \emph{sans dioxygène} (anaérobiose), avec un rendement énergétique faible comparé à la respiration.
\item \textbf{Relier les observations aux produits} : mousse/gaz $\to$ \ce{CO2} ; goût acide $\to$ acide lactique ou acétique ; odeur d'alcool $\to$ éthanol.
\item \textbf{Citer les applications} : boissons, laitages, conservation des aliments (l'acidification inhibe les germes de putréfaction).
\end{enumerate}
\textbf{Réflexe} : si l'énoncé décrit un dégagement gazeux, pense immédiatement au \ce{CO2} de la fermentation alcoolique.
\end{methode}

\begin{exoresolu}[Fermentation lactique et conservation du lait]
Une coopérative de femmes de Ngaoundéré transforme le lait frais en lait caillé. On constate que le lait caillé, contrairement au lait frais, se conserve plusieurs jours sans se putréfier. Le pH du lait passe de 6,6 à 4,5 au cours de la transformation.

Explique le principe biotechnologique de cette transformation et justifie la conservation prolongée du lait caillé.
\tcblower
\textbf{Solution détaillée}

\textbf{Étape 1 : identification du type de fermentation.} La baisse du pH (de 6,6 à 4,5) traduit une \cle{acidification} du milieu : il s'agit d'une \cle{fermentation lactique}, réalisée par des bactéries lactiques (genres \emph{Lactobacillus}, \emph{Streptococcus}).

\textbf{Étape 2 : équation bilan.} Ces bactéries dégradent le lactose (sucre du lait) en acide lactique, en anaérobiose :
\[ \ce{C6H12O6 -> 2C3H6O3} \]
L'acide lactique formé est responsable de la chute du pH et de la coagulation des protéines du lait (caséines), d'où la consistance « caillée ».

\textbf{Étape 3 : rôle énergétique.} Cette fermentation fournit aux bactéries l'ATP nécessaire à leur métabolisme, en l'absence de dioxygène.

\textbf{Étape 4 : justification de la conservation.} Le pH acide (4,5) crée un milieu défavorable au développement des micro-organismes de putréfaction et de nombreux germes pathogènes, qui ne se multiplient qu'à pH proche de la neutralité. L'acidification agit donc comme un \cle{conservateur naturel}.

\textbf{Conclusion :} la fabrication du lait caillé repose sur la fermentation lactique du lactose par des bactéries lactiques ; l'acide lactique produit acidifie le milieu, ce qui coagule le lait et inhibe les germes de putréfaction, assurant une conservation prolongée.
\end{exoresolu}

\courssub{Fiche méthode 4 : Décrire le principe de la sélection assistée}

\begin{methode}[Expliquer la sélection assistée]
La sélection assistée (par marqueurs moléculaires) est une biotechnologie moderne d'amélioration des plantes et des animaux. Pour la décrire :
\begin{enumerate}
\item \textbf{Rappeler le principe de la sélection classique} : choisir comme reproducteurs les individus présentant le caractère souhaité (rendement, résistance à une maladie), génération après génération.
\item \textbf{Montrer sa limite} : elle est lente, car il faut attendre que le caractère s'exprime (plante adulte, animal en production) pour le détecter.
\item \textbf{Expliquer l'apport des marqueurs moléculaires} : on identifie dans l'ADN des séquences (marqueurs) associées au gène d'intérêt ; on peut alors repérer les individus porteurs du caractère \emph{dès la graine ou le jeune âge}, par simple analyse d'ADN.
\item \textbf{Préciser ce que ce n'est pas} : la sélection assistée \textbf{ne modifie pas} le patrimoine génétique ; elle accélère le choix des reproducteurs. Ne pas la confondre avec la transgenèse (OGM).
\item \textbf{Conclure sur l'intérêt} : gain de temps considérable, création plus rapide de variétés performantes (maïs, manioc, cacao résistants).
\end{enumerate}
\end{methode}

\begin{exoresolu}[Sélection assistée pour un maïs résistant]
Un programme de recherche de l'IRAD (Institut de Recherche Agricole pour le Développement) souhaite créer une variété de maïs résistante à la striure, maladie fréquente dans la zone soudanienne du Cameroun. La sélection classique exige d'attendre la floraison et l'infection naturelle des plants pour repérer les résistants, soit environ 10 ans de croisements. Grâce à des marqueurs moléculaires liés au gène de résistance, les chercheurs repèrent les plants porteurs dès le stade plantule.

Explique le principe de la sélection assistée mise en œuvre et dégage son intérêt par rapport à la sélection classique.
\tcblower
\textbf{Solution détaillée}

\textbf{Étape 1 : principe de la sélection classique.} Elle consiste à croiser des plants et à ne retenir comme reproducteurs que les descendants qui, une fois adultes et exposés à la maladie, manifestent la résistance. C'est un tri sur le \cle{phénotype}, long car le caractère ne s'exprime que tardivement et de façon dépendante des conditions du milieu.

\textbf{Étape 2 : apport des marqueurs moléculaires.} Les chercheurs ont identifié, dans l'ADN du maïs, des séquences spécifiques situées à proximité du gène de résistance et héritées avec lui : ce sont les \cle{marqueurs moléculaires}. En extrayant l'ADN d'une jeune plantule et en recherchant ces marqueurs, on sait si le plant possède le gène de résistance \emph{avant même} que le caractère s'exprime.

\textbf{Étape 3 : ce qui change concrètement.} On ne sélectionne plus sur l'apparence adulte mais directement sur le \cle{génotype}. Seuls les plants porteurs du marqueur sont conservés pour les croisements ; les autres sont éliminés dès la pépinière, ce qui économise temps, espace et main-d'œuvre.

\textbf{Étape 4 : précision essentielle.} Aucun gène étranger n'est introduit : le patrimoine génétique du maïs n'est pas modifié. La sélection assistée n'est donc \textbf{pas} une technique de création d'OGM ; elle accélère simplement une sélection naturelle dirigée.

\textbf{Conclusion :} la sélection assistée par marqueurs moléculaires repose sur la détection précoce, dans l'ADN, de séquences associées au caractère recherché ; elle réduit fortement la durée de création variétale (de l'ordre de 10 ans à quelques années) et permet d'obtenir plus vite un maïs résistant à la striure pour les producteurs camerounais.
\end{exoresolu}

\courssub{Fiche méthode 5 : Décrire une biotechnologie enzymatique}

\begin{methode}[Expliquer l'utilisation industrielle d'une enzyme]
Pour décrire une biotechnologie enzymatique :
\begin{enumerate}
\item \textbf{Définir l'enzyme} : protéine catalysatrice qui accélère une réaction biochimique spécifique sans être consommée.
\item \textbf{Nommer l'enzyme et sa réaction} : substrat $\to$ produit (ex. : amylase : amidon $\to$ sucres simples ; pectinase : pectine dégradée $\to$ jus clarifié).
\item \textbf{Préciser l'origine de l'enzyme} : extraite de micro-organismes (champignons, bactéries) cultivés en bioréacteurs, puis purifiée.
\item \textbf{Indiquer les conditions d'application} : température et pH optimaux ; possibilité d'enzymes \cle{immobilisées} (fixées sur un support) réutilisables, ce qui réduit les coûts.
\item \textbf{Citer l'application} : industrie agroalimentaire (jus, fromages, panification), lessives (protéases, lipases), textile, santé (tests diagnostiques).
\end{enumerate}
\textbf{Réflexe} : l'intérêt clé d'une enzyme = \textbf{spécificité} (une seule réaction ciblée) + \textbf{efficacité} (catalyse rapide à température modérée, donc économie d'énergie).
\end{methode}

\begin{exoresolu}[Les enzymes dans la production de jus de fruits]
Une unité de transformation de mangues à Maroua obtient un jus trouble et peu abondant par simple pressage. En ajoutant une préparation de pectinase avant le pressage, le rendement en jus augmente et le jus devient limpide.

Explique le principe de cette biotechnologie enzymatique et précise ses avantages.
\tcblower
\textbf{Solution détaillée}

\textbf{Étape 1 : le problème technologique.} La pulpe de mangue contient de la \cle{pectine}, polysaccharide qui lie les cellules entre elles et retient l'eau : elle rend la pulpe visqueuse, piège le jus et provoque le trouble.

\textbf{Étape 2 : l'enzyme et sa réaction.} La \cle{pectinase} est une enzyme produite industriellement par des champignons microscopiques (ex. : \emph{Aspergillus niger}) cultivés en bioréacteurs. Elle catalyse spécifiquement l'hydrolyse de la pectine :
\[ \ce{\text{pectine} ->[$\text{pectinase}$] \text{galacturonates solubles}} \]
La dégradation de la pectine libère le jus retenu dans la pulpe et élimine les particules responsables du trouble.

\textbf{Étape 3 : conditions d'application.} L'enzyme agit à température modérée (environ \SI{40}{\celsius} à \SI{50}{\celsius}) et à un pH légèrement acide, conditions douces compatibles avec la qualité gustative du fruit.

\textbf{Étape 4 : avantages du procédé.}
\begin{itemize}
\item \textbf{Spécificité} : seule la pectine est dégradée ; les arômes et vitamines sont préservés.
\item \textbf{Rendement} : le volume de jus extrait augmente nettement.
\item \textbf{Économie d'énergie} : la réaction se fait à température modérée, sans chauffage intense.
\item \textbf{Qualité} : le jus obtenu est limpide, plus stable et plus apprécié des consommateurs.
\end{itemize}

\textbf{Conclusion :} l'ajout de pectinase est une biotechnologie enzymatique moderne : une enzyme microbienne purifiée catalyse la dégradation de la pectine de la pulpe, ce qui améliore le rendement et la limpidité du jus de mangue dans des conditions douces et économes en énergie.
\end{exoresolu}

\courssub{Fiche méthode 6 : Discuter les enjeux d'une application biotechnologique}

\begin{methode}[Construire une discussion d'enjeux]
Discuter des enjeux exige un propos \textbf{équilibré et structuré}, pas un simple avis. Procède ainsi :
\begin{enumerate}
\item \textbf{Identifier les enjeux économiques} : création d'emplois et d'industries, augmentation des rendements agricoles, réduction des importations (médicaments, enzymes), valorisation des produits locaux, mais aussi coût des équipements et dépendance technologique.
\item \textbf{Identifier les enjeux sanitaires et environnementaux} : production de vaccins et de médicaments, dépollution ; mais risques éventuels (dissémination d'organismes, perte de biodiversité, résidus).
\item \textbf{Identifier les enjeux éthiques} : respect de la vie, brevetabilité du vivant, équité d'accès aux innovations, acceptabilité sociale (ex. : OGM), protection des savoirs traditionnels.
\item \textbf{Peser avantages et limites} : présente les deux faces avec des arguments factuels, sans caricature.
\item \textbf{Conclure par une position argumentée} : par exemple, plaidoyer pour un encadrement réglementaire (biosécurité) permettant de bénéficier des biotechnologies tout en maîtrisant les risques.
\end{enumerate}
\textbf{Réflexe Probatoire} : une bonne discussion cite au moins \textbf{deux avantages} et \textbf{deux limites ou risques}, illustrés d'exemples.
\end{methode}

\begin{exoresolu}[Enjeux de l'introduction de variétés végétales améliorées au Cameroun]
Le Cameroun importe une partie de ses semences améliorées et de ses produits pharmaceutiques. Le développement d'une filière nationale de biotechnologies (sélection assistée, production locale d'enzymes et de médicaments) est envisagé.

Discute des enjeux économiques et éthiques de ce développement.
\tcblower
\textbf{Solution détaillée}

\textbf{Enjeux économiques favorables :}
\begin{itemize}
\item \textbf{Augmentation de la production agricole} : des variétés de maïs, manioc ou cacao à haut rendement et résistantes aux maladies accroissent les revenus des producteurs et la sécurité alimentaire.
\item \textbf{Réduction des importations} : produire localement semences, enzymes et médicaments (ex. : insuline) diminue la dépendance et les dépenses en devises.
\item \textbf{Création d'emplois qualifiés} : laboratoires, unités de transformation agroalimentaire, contrôle qualité.
\end{itemize}

\textbf{Enjeux économiques défavorables :}
\begin{itemize}
\item \textbf{Coût élevé des équipements} (bioréacteurs, laboratoires d'analyse d'ADN), difficilement accessible aux petites structures.
\item \textbf{Risque de dépendance technologique} si les souches, enzymes ou brevets restent détenus par des firmes étrangères.
\end{itemize}

\textbf{Enjeux éthiques :}
\begin{itemize}
\item \textbf{Équité d'accès} : les innovations doivent bénéficier aussi aux petits agriculteurs et aux populations rurales, pas seulement aux grandes exploitations.
\item \textbf{Brevetabilité du vivant} : peut-on s'approprier par brevet une souche microbienne ou un gène ? Cela pose la question de la marchandisation du vivant.
\item \textbf{Protection des savoirs traditionnels} : les pratiques locales de fermentation (vin de palme, lait caillé) et les ressources génétiques locales doivent être valorisées et non confisquées.
\item \textbf{Biosécurité} : l'utilisation d'organismes génétiquement modifiés exige une réglementation stricte pour préserver la biodiversité et la santé.
\end{itemize}

\textbf{Conclusion :} le développement des biotechnologies au Cameroun représente une opportunité économique et sanitaire majeure, à condition d'être encadré par une réglementation de biosécurité, des politiques d'accès équitable et la protection des savoirs locaux : c'est le prix d'un progrès à la fois performant et éthiquement acceptable.
\end{exoresolu}

\begin{attention}[Les 5 erreurs qui coûtent des points au Probatoire]
\begin{enumerate}
\item \textbf{Confondre fermentation et respiration.} La fermentation se déroule \emph{sans dioxygène} (anaérobiose) et produit peu d'ATP ; la respiration consomme du \ce{O2} et produit \ce{CO2} et \ce{H2O}. Écrire que la levure « respire » le glucose en produisant de l'alcool est une erreur de fond.
\item \textbf{Confondre sélection assistée et transgenèse (OGM).} La sélection assistée par marqueurs moléculaires ne modifie \emph{pas} le génome : elle accélère le choix des reproducteurs. La transgenèse introduit un gène étranger. Les confondre dans une copie est sanctionné.
\item \textbf{Décrire un procédé de façon incomplète.} Oublier l'agent biologique, le substrat ou les conditions (anaérobiose, température) revient à décrire une « recette » et non un principe biotechnologique. Pense au quintet : agent + substrat + transformation + conditions + produit.
\item \textbf{Écrire une équation de fermentation fausse ou non équilibrée.} La fermentation alcoolique est \ce{C6H12O6 -> 2C2H5OH + 2CO2} : vérifie les coefficients (2 éthanol, 2 \ce{CO2}). Ne jamais faire apparaître \ce{O2} dans une fermentation.
\item \textbf{Traiter les enjeux de manière unilatérale.} Ne citer que les avantages (ou que les risques) d'une biotechnologie ne répond pas à la consigne « discuter ». Le correcteur attend au moins deux arguments favorables, deux limites, et une conclusion nuancée mentionnant l'encadrement (biosécurité, éthique).
\end{enumerate}
\end{attention}

\newpage

\coursec{Exercices}

\courssub{Je vérifie mes connaissances}

\begin{exercice}[QCM : les définitions de base]
Pour chaque question, une seule réponse est correcte. Recopie la lettre de la bonne réponse.

\textbf{1.} La biotechnologie est :
\begin{itemize}
\item[a)] l'étude des êtres vivants dans leur milieu naturel ;
\item[b)] l'utilisation d'organismes vivants, de cellules ou de molécules biologiques pour produire des biens et des services ;
\item[c)] la fabrication de robots imitant le vivant ;
\item[d)] l'ensemble des techniques de médecine traditionnelle.
\end{itemize}

\textbf{2.} La fermentation alcoolique est :
\begin{itemize}
\item[a)] la dégradation du glucose en éthanol et en dioxyde de carbone par des levures, en absence de dioxygène ;
\item[b)] la combustion complète du glucose en présence de dioxygène ;
\item[c)] la synthèse du glucose par les plantes vertes ;
\item[d)] la transformation du lait en yaourt par des bactéries.
\end{itemize}

\textbf{3.} La biosécurité désigne :
\begin{itemize}
\item[a)] la protection des cultures contre les insectes ravageurs ;
\item[b)] l'ensemble des mesures visant à prévenir les risques liés à la manipulation d'organismes vivants ou de leurs produits pour l'homme et l'environnement ;
\item[c)] la sécurité alimentaire d'une population ;
\item[d)] la conservation des semences dans une banque de gènes.
\end{itemize}

\textbf{4.} La culture in vitro est une biotechnologie qui consiste à :
\begin{itemize}
\item[a)] cultiver des plantes en plein champ sous serre chauffée ;
\item[b)] faire pousser des cellules ou des tissus végétaux sur un milieu nutritif stérile, en conditions contrôlées, pour obtenir des plants entiers ;
\item[c)] croiser deux variétés de plantes par pollinisation naturelle ;
\item[d)] conserver les graines au réfrigérateur.
\end{itemize}
\end{exercice}

\begin{exercice}[Vrai ou faux, en justifiant]
Pour chaque affirmation, indique si elle est \textbf{vraie} ou \textbf{fausse}, puis justifie ta réponse en une ou deux phrases.

\textbf{1.} La fabrication du vin de palme par fermentation du suc de palmier est une biotechnologie moderne.

\textbf{2.} Les enzymes utilisées dans l'industrie agroalimentaire sont des catalyseurs biologiques qui accélèrent les réactions sans être consommées.

\textbf{3.} La sélection assistée par marqueurs moléculaires permet de choisir plus rapidement les variétés végétales porteuses d'un caractère d'intérêt que la sélection traditionnelle.

\textbf{4.} Les biotechnologies ne posent aucune question éthique car elles utilisent des organismes naturels.
\end{exercice}

\begin{exercice}[Définitions à compléter]
Recopie et complète chaque définition avec les mots ou groupes de mots qui conviennent.

\textbf{1.} La \cle{biotechnologie} est l'ensemble des techniques qui utilisent des \dots\dots\dots\dots, des \dots\dots\dots\dots ou des \dots\dots\dots\dots biologiques (enzymes) pour produire des biens et des services utiles à l'homme.

\textbf{2.} On appelle \cle{substrat} la substance \dots\dots\dots\dots par le microorganisme ou l'enzyme au cours d'un procédé biotechnologique.

\textbf{3.} La \cle{sélection assistée} est une biotechnologie moderne qui utilise des \dots\dots\dots\dots moléculaires pour identifier et choisir les individus porteurs d'un caractère d'intérêt.

\textbf{4.} Un \cle{fermenteur} (ou bioréacteur) est une cuve dans laquelle on contrôle les conditions (température, pH, aération) de la \dots\dots\dots\dots réalisée par des microorganismes.
\end{exercice}

\begin{exercice}[Calcul direct : rendement de la fermentation alcoolique]
L'équation globale de la fermentation alcoolique réalisée par la levure \emph{Saccharomyces cerevisiae} s'écrit :
\[ \ce{C6H12O6 -> 2 C2H5OH + 2 CO2} \]
On donne les masses molaires : $M(\ce{C6H12O6}) = \SI{180}{g/mol}$ ; $M(\ce{C2H5OH}) = \SI{46}{g/mol}$.

Une brasserie fait fermenter $\SI{90}{g}$ de glucose.
\begin{enumerate}
\item Calcule la quantité de matière (en mol) de glucose utilisée.
\item En t'aidant de l'équation, détermine la quantité de matière d'éthanol théoriquement produite.
\item Déduis-en la masse d'éthanol théorique obtenue.
\item En réalité, on ne récolte que $\SI{36,8}{g}$ d'éthanol. Calcule le rendement de la fermentation (en \%).
\end{enumerate}
\end{exercice}

\courssub{Je m'entraîne}

\begin{exercice}[Le bobolo : analyser un procédé traditionnel]
Le « bobolo », mets très apprécié au Cameroun, est obtenu à partir de bâtons de manioc. Le procédé traditionnel est le suivant : les bâtons de manioc sont pelés, puis trempés dans l'eau pendant 2 à 4 jours ; pendant ce temps, des microorganismes (bactéries lactiques et levures naturellement présentes) se développent et ramollissent la pulpe en dégradant une partie de l'amidon. La pulpe fermentée est ensuite essorée, moulée en boudins, enveloppée dans des feuilles et cuite à la vapeur.

\begin{enumerate}
\item Identifie, dans ce procédé : l'agent biologique, le substrat, le type de procédé biotechnologique et le produit obtenu.
\item Classe cette biotechnologie (traditionnelle ou moderne) en justifiant ta réponse par deux arguments.
\item Cite un avantage de l'étape de fermentation pour la qualité du produit (goût, conservation ou sécurité alimentaire).
\end{enumerate}
\end{exercice}

\begin{exercice}[Tableau comparatif à compléter]
Recopie et complète le tableau suivant, qui compare une biotechnologie traditionnelle (fabrication du vin de palme par fermentation spontanée) et une biotechnologie moderne (production industrielle d'enzymes par des microorganismes cultivés en bioréacteur), selon les cinq critères imposés.

\begin{center}
\begin{tabularx}{\linewidth}{|l|X|X|}
\hline
\rowcolor{popblueL} \textbf{Critère} & \textbf{Biotechnologie traditionnelle} & \textbf{Biotechnologie moderne} \\
\hline
Maîtrise des microorganismes utilisés & \dots & \dots \\
\hline
Contrôle des conditions du procédé (température, pH) & \dots & \dots \\
\hline
Échelle de production & \dots & \dots \\
\hline
Niveau de connaissances scientifiques requis & \dots & \dots \\
\hline
Régularité de la qualité du produit & \dots & \dots \\
\hline
\end{tabularx}
\end{center}
\end{exercice}

\begin{exercice}[Les amylases dans la bière de maïs]
La fabrication traditionnelle de la bière de maïs (appelée « bil-bil » dans le Nord-Cameroun) comporte une étape de maltage : les grains de maïs sont mis à germer, ce qui provoque la production d'enzymes appelées \cle{amylases}. Lors du brassage, ces amylases agissent sur l'amidon contenu dans les grains ; les levures interviennent ensuite.

\begin{enumerate}
\item Rappelle la nature chimique d'une enzyme.
\item Précise le substrat sur lequel agissent les amylases et le produit de leur action.
\item Explique, en cinq lignes maximum, pourquoi l'action des amylases est indispensable avant l'étape de fermentation alcoolique par les levures.
\end{enumerate}
\end{exercice}

\begin{exercice}[Classer des applications biotechnologiques]
Voici une liste d'applications des biotechnologies :
\begin{itemize}
\item fabrication du yaourt à partir du lait ;
\item production d'insuline par des bactéries modifiées ;
\item épuration des eaux usées de la ville de Douala par des boues activées (bactéries) ;
\item fabrication du « kinkeliba » séché pour la tisane ;
\item production de bioéthanol à partir de la canne à sucre ;
\item sélection de variétés de maïs résistantes à la sécheresse à l'aide de marqueurs moléculaires.
\end{itemize}

\begin{enumerate}
\item Classe ces applications dans les trois grands domaines d'application de la biotechnologie : agroalimentaire, santé, environnement. (Certaines applications peuvent relever de deux domaines : justifie alors ton choix.)
\item Parmi ces applications, identifie deux biotechnologies traditionnelles et deux biotechnologies modernes, en justifiant chaque choix.
\end{enumerate}
\end{exercice}

\courssub{Je me prépare au Probatoire}

\begin{exercice}[Sujet de synthèse : la culture in vitro du plantain]
\textbf{Contexte.} La coopérative agricole de ton village, dans la région du Centre, produit du plantain. Les plants traditionnels (rejets prélevés au pied des plants mères) sont souvent porteurs de maladies (virus, nématodes) et leur production est lente : un pied ne fournit que quelques rejets par an. Un technicien agricole propose à la coopérative d'adopter la \cle{culture in vitro} : on prélève de petits fragments de tissus sains (méristèmes) sur un plant d'élite, on les place sur un milieu nutritif stérile contenant des sucres, des sels minéraux et des hormones de croissance, dans des conditions contrôlées de température et de lumière. Chaque fragment donne de nombreux plantules identiques, qui sont ensuite acclimatés en pépinière avant la mise en champ.

\textbf{Travail à faire.}
\begin{enumerate}
\item Définis la biotechnologie et indique à quel type de biotechnologie (traditionnelle ou moderne) appartient la culture in vitro, en justifiant ta réponse.
\item Décris, en une dizaine de lignes, le principe du procédé de culture in vitro du plantain, en précisant : l'agent biologique utilisé, le milieu de culture et les conditions à contrôler.
\item Présente deux avantages économiques de ce procédé pour la coopérative, en les expliquant.
\item Identifie une question éthique ou environnementale que la coopérative devrait examiner avant de se lancer (par exemple : uniformité génétique des plantations, dépendance vis-à-vis des laboratoires fournisseurs, coût d'accès pour les petits producteurs), et propose une mesure pour y répondre.
\end{enumerate}
\end{exercice}

\begin{exercice}[Étude de document : les bactéries au service de l'assainissement]
\textbf{Document.} La commune de Bafoussam rejette chaque jour des milliers de mètres cubes d'eaux usées chargées de matière organique. Si ces eaux sont déversées directement dans les cours d'eau, la matière organique est dégradée par les microorganismes naturels, qui consomment alors le dioxygène dissous : les poissons meurent et des odeurs nauséabondes apparaissent. Une station d'épuration à \cle{boues activées} reproduit ce phénomène de façon contrôlée : les eaux usées sont brassées et fortement aérées dans de grands bassins contenant une flore bactérienne concentrée (les « boues activées »). Les bactéries dégradent la matière organique en quelques heures. Une partie des boues est recyclée dans le bassin, l'excédent est extrait et peut être composté. L'eau traitée, débarrassée de l'essentiel de sa pollution organique, est rejetée dans le milieu naturel.

\textbf{Questions.}
\begin{enumerate}
\item Dégage du document le problème environnemental que la station d'épuration permet de résoudre.
\item Décris le principe du procédé biotechnologique utilisé, en précisant l'agent biologique, le substrat et le rôle de l'aération.
\item Justifie le classement de cette application dans le domaine « environnement » des biotechnologies.
\item Discute deux enjeux de l'installation d'une telle station pour une commune camerounaise : un enjeu sanitaire ou environnemental (bénéfice attendu) et un enjeu économique ou social (coût, entretien, formation du personnel, valorisation des boues).
\end{enumerate}
\end{exercice}

\begin{exercice}[Problème : produire du bioéthanol à partir de la canne à sucre]
\textbf{Contexte et données.} Pour réduire ses importations de carburant, une unité pilote installée près de Nkoteng envisage de produire du bioéthanol par fermentation alcoolique du jus de canne à sucre par la levure \emph{Saccharomyces cerevisiae}. Le saccharose du jus est d'abord hydrolysé en glucose et en fructose ; on admettra que seul le glucose intervient dans les calculs, selon l'équation :
\[ \ce{C6H12O6 -> 2 C2H5OH + 2 CO2} \]
Données : $M(\ce{C6H12O6}) = \SI{180}{g/mol}$ ; $M(\ce{C2H5OH}) = \SI{46}{g/mol}$ ; masse volumique de l'éthanol : $\rho = \SI{0,79}{kg/L}$. Une cuve contient $\SI{360}{kg}$ de glucose ; le rendement réel de la fermentation est de $90\,\%$.

\textbf{Questions.}
\begin{enumerate}
\item Nomme le type de procédé biotechnologique mis en jeu et précise l'agent biologique et le substrat.
\item Calcule la masse théorique d'éthanol produite par la cuve, puis la masse réelle compte tenu du rendement.
\item Déduis-en le volume d'éthanol (en litres) effectivement obtenu.
\item La fermentation dégage du dioxyde de carbone. Calcule la masse de \ce{CO2} réellement dégagée (on donne $M(\ce{CO2}) = \SI{44}{g/mol}$) et propose une valorisation possible de ce sous-produit.
\item Présente un avantage économique et une limite (environnementale, sociale ou éthique) de la production de bioéthanol à grande échelle au Cameroun, en argumentant chaque point.
\end{enumerate}
\end{exercice}



\newpage

\coursec{Corrigés des exercices}

\courssub{Je vérifie mes connaissances}

\begin{corrige}[QCM : les définitions de base]
\textbf{1. b)} C'est la définition officielle (\cle{OCDE}, \cle{MINESEC}). La \cle{biotechnologie} met en œuvre des \cle{agents biologiques} (microorganismes, cellules, enzymes) pour des applications utiles.
\textbf{2. a)} La \cle{fermentation alcoolique} est une \cle{anaérobiolyse} (absence de \ce{O2}) du glucose en éthanol et \ce{CO2} par des levures (ex: \emph{Saccharomyces cerevisiae}). La proposition d) correspond à la \cle{fermentation lactique}.
\textbf{3. b)} La \cle{biosécurité} (ou sûreté biologique) encadre la manipulation d'\cle{OGM}, de pathogènes, etc., pour protéger la santé humaine et l'environnement (Protocole de Carthagène).
\textbf{4. b)} La \cle{culture in vitro} (micropropagation) utilise la \cle{totipotence} des cellules végétales sur milieu stérile (gélose + sels minéraux + sucres + hormones) en conditions aseptiques.
\end{corrige}

\begin{corrige}[Vrai ou faux, en justifiant]
\textbf{1. Faux.} C'est une \textbf{\cle{biotechnologie traditionnelle}} (fermentation spontanée par la flore microbienne naturelle du suc, sans maîtrise des souches ni contrôle des paramètres physico-chimiques).
\textbf{2. Vrai.} Par définition, une \cle{enzyme} est une protéine catalytique qui abaisse l'énergie d'activation d'une réaction chimique sans être modifiée ni consommée au bilan global.
\textbf{3. Vrai.} La \cle{sélection assistée par marqueurs (SAM)} permet un criblage précoce (au stade plantule) et hors champ (au laboratoire) des allèles d'intérêt, raccourcissant considérablement le cycle de sélection par rapport à la sélection phénotypique classique qui attend l'expression du caractère.
\textbf{4. Faux.} Les \cle{biotechnologies modernes} (OGM, édition de génome, clonage, diagnostic prénatal, brevetage du vivant) soulèvent d'importants débats éthiques (respect de la biodiversité, équité d'accès, brevetage des gènes, risques sanitaires, modification de la lignée germinale humaine).
\end{corrige}

\begin{corrige}[Définitions à compléter]
\textbf{1.} \dots\dots \textbf{organismes vivants}, des \textbf{cellules} ou des \textbf{molécules} biologiques \dots\dots
\textbf{2.} \dots\dots la substance \textbf{transformée} (ou \textbf{consommée}, \textbf{utilisée}) \dots\dots
\textbf{3.} \dots\dots des \textbf{marqueurs} moléculaires \dots\dots
\textbf{4.} \dots\dots de la \textbf{fermentation} (ou \textbf{culture}, \textbf{bioconversion}) \dots\dots
\end{corrige}

\begin{corrige}[Calcul direct : rendement de la fermentation alcoolique]
\begin{enumerate}
\item $n_{\text{glucose}} = \dfrac{m}{M} = \dfrac{90}{180} = \mathbf{\num{0,50}\,mol}$.
\item D'après l'équation \ce{C6H12O6 -> 2 C2H5OH + 2 CO2} : $1$ mol glucose $\rightarrow$ $2$ mol éthanol. Donc $n_{\text{éthanol (th)}} = 2 \times \num{0,50} = \mathbf{\num{1,00}\,mol}$.
\item $m_{\text{éthanol (th)}} = n \times M = \num{1,00} \times 46 = \mathbf{\num{46}\,g}$.
\item Rendement $\eta = \dfrac{m_{\text{réel}}}{m_{\text{théorique}}} \times 100 = \dfrac{\num{36,8}}{46} \times 100 = \mathbf{80\,\%}$.
\end{enumerate}
\end{corrige}

\courssub{Je m'entraîne}

\begin{corrige}[Le bobolo : analyser un procédé traditionnel]
\begin{enumerate}
\item \textbf{Agent biologique :} Bactéries lactiques et levures (flore microbienne naturelle, non sélectionnée). \textbf{Substrat :} Amidon (et glucides solubles) de la pulpe de manioc. \textbf{Type de procédé :} Fermentation lactique et alcoolique mixte (\cle{biotechnologie traditionnelle}). \textbf{Produit obtenu :} Pâte de manioc fermentée (bobolo) cuite à la vapeur.
\item \textbf{Biotechnologie traditionnelle.} Arguments : (1) Utilisation de la flore microbienne naturelle (inoculum non maîtrisé, non purifié) ; (2) Absence de contrôle des paramètres physico-chimiques (température, pH, aération) et du temps de fermentation (empirique, basé sur le savoir-faire).
\item \textbf{Avantages (au choix) :} \textbf{Sécurité alimentaire :} Dégradation des glycosides cyanogènes (linamarine) du manioc amer, réduisant la toxicité (cyanure). \textbf{Conservation :} Production d'acide lactique (abaissement du pH) et d'éthanol qui inhibent les germes de contamination et de putréfaction. \textbf{Qualité organoleptique :} Développement d'arômes caractéristiques et ramollissement de la texture (gélification partielle de l'amidon).
\end{enumerate}
\end{corrige}

\begin{corrige}[Tableau comparatif à compléter]
\begin{center}
\begin{tabularx}{\linewidth}{|l|X|X|}
\hline
\rowcolor{popblueL} \textbf{Critère} & \textbf{Biotechnologie traditionnelle (Vin de palme)} & \textbf{Biotechnologie moderne (Enzymes industrielles)} \\
\hline
Maîtrise des microorganismes utilisés & Faible : flore naturelle, mixte, non sélectionnée, variable & Élevée : souches pures, sélectionnées, souvent modifiées génétiquement (\cle{OGM}), identifiées génomiquement \\
\hline
Contrôle des conditions du procédé (température, pH) & Aucun ou rudimentaire (température ambiante, pas de régulation pH) & Contrôle précis et continu via automate (\cle{bioréacteur}) : thermostatation, pompes pH, débit d'air stérile, agitation \\
\hline
Échelle de production & Artisanale, familiale, faible volume (litres à hectolitres) & Industrielle, grande échelle (fermenteurs de $\SI{10}{m^3}$ à $\SI{500}{m^3}$ et plus) \\
\hline
Niveau de connaissances scientifiques requis & Savoir-faire empirique, transmission orale, observation & Sciences fondamentales (génétique, biochimie, microbiologie, génie des procédés), R\&D \\
\hline
Régularité de la qualité du produit & Variable (dépend de la saison, du collecteur, du contenant) & Constante, standardisée, traçable (cahier des charges, normes ISO) \\
\hline
\end{tabularx}
\end{center}
\end{corrige}

\begin{corrige}[Les amylases dans la bière de maïs]
\begin{enumerate}
\item Une \cle{enzyme} est une \textbf{protéine} (souvent globulaire) à activité catalytique. Elle possède un \cle{site actif} spécifique pour son \cle{substrat}.
\item \textbf{Substrat :} L'\cle{amidon} (polymère de glucose : amylose et amylopectine). \textbf{Produits :} Principalement du \textbf{maltose} (disaccharide) et des \textbf{dextrines} (oligosaccharides), in fine du glucose.
\item Les levures (\emph{Saccharomyces}) ne possèdent pas d'amylases extracellulaires : elles sont incapables d'hydrolyser l'amidon (macromolécule trop grosse pour traverser leur membrane). L'action des amylases (produites pendant la germination du grain = \cle{maltage}) hydrolyse l'amidon en sucres fermentescibles (maltose, glucose) que les levures peuvent alors assimiler par glycolyse pour produire de l'éthanol et du \ce{CO2}. Sans maltage, pas de substrat utilisable par la levure, donc pas de fermentation alcoolique.
\end{enumerate}
\end{corrige}

\begin{corrige}[Classer des applications biotechnologiques]
\begin{enumerate}
\item \textbf{Agroalimentaire :} Yaourt (fermentation lactique), Bioéthanol (fermentation alcoolique - aussi énergie), Sélection maïs (amélioration variétale). \textbf{Santé :} Insuline recombinante (thérapeutique). \textbf{Environnement :} Épuration eaux usées (boues activées - dépollution), Bioéthanol (carburant substitut aux fossiles - atténuation changement climatique). \textit{Note :} Le séchage du kinkeliba est une \textbf{technique de conservation physique} (déshydratation), pas une biotechnologie (pas d'agent biologique transformateur).
\item \textbf{Traditionnelles :} (1) \textbf{Yaourt :} Fermentation lactique spontanée ou sur levain naturel, sans contrôle stérile ni souches pures définies, savoir-faire ancestral. (2) \textbf{Bioéthanol (contexte artisanal/traditionnel) :} Si produit par fermentation de jus de canne ou palme sans contrôle aseptique ni souches sélectionnées. \textbf{Modernes :} (1) \textbf{Insuline :} Bactéries \emph{E. coli} ou levures génétiquement modifiées (\cle{ADN recombinant}), culture en bioréacteur aseptique, purification industrielle. (2) \textbf{Sélection maïs par marqueurs (SAM) :} Utilisation de marqueurs moléculaires (ADN), génotypage au laboratoire, sélection précoce sans attente phénotypique. (3) \textbf{Boues activées (station d'épuration) :} Procédé ingénierisé (aération forcée, recirculation boues, contrôle paramètres), flore sélectionnée/adaptée, échelle industrielle.
\end{enumerate}
\end{corrige}

\courssub{Je me prépare au Probatoire}

\begin{corrige}[Sujet de synthèse : la culture in vitro du plantain]
\begin{enumerate}
\item \textbf{Définition :} La \cle{biotechnologie} est l'ensemble des techniques utilisant des organismes vivants, des cellules ou des molécules biologiques pour produire des biens ou des services. \textbf{Type :} \textbf{\cle{Biotechnologie moderne.}} \textbf{Justification :} Elle nécessite un laboratoire (milieu stérile, conditions aseptiques), une maîtrise scientifique (totipotence, rôle des hormones auxine/cytokinine), un contrôle précis des paramètres (température, photopériode, stérilité) et utilise souvent du matériel végétal sain (méristème exempt de virus) issu de sélection variétale.
\item \textbf{Principe :} On exploite la \cle{totipotence} des cellules végétales. \textbf{Agent biologique :} Explants méristématiques (apex caulinaire, $\approx \SI{0,5}{mm}$) prélevés sur un plant mère d'élite (sain, productif). \textbf{Milieu de culture :} Gélose (agar) + sels minéraux (MS de Murashige \& Skoog) + source de carbone (saccharose $\SI{20}{g/L}$) + vitamines + régulateurs de croissance (cytokinines type BAP pour multiplication des bourgeons, auxines type NAA pour rhizogenèse). \textbf{Conditions contrôlées :} Asepsie stricte (hotte à flux laminaire, autoclave), température $\SI{25 \pm 2}{\celsius}$, photopériode $\SI{16}{h}$ lumière / $\SI{8}{h}$ obscurité, intensité lumineuse $\approx \SI{2000}{lux}$. Étapes : initiation $\rightarrow$ multiplication (subcultures successives) $\rightarrow$ élongation $\rightarrow$ racinement $\rightarrow$ acclimatation (serre humide, substrat stérile) $\rightarrow$ champ.
\item \textbf{Avantage 1 : Assainissement (gain de rendement).} Le prélèvement au méristème élimine virus (BBTV, BSV) et nématodes. Les plants in vitro sont \textbf{sains}, ce qui assure une productivité maximale dès le premier cycle, contrairement aux rejets contaminés qui dégénèrent.
\item \textbf{Avantage 2 : Multiplication massive et rapide (gain de temps/échelle).} Un seul plant mère peut produire des \textbf{milliers de plants uniformes en quelques mois} (vs quelques rejets/an). Cela permet de planter de grandes surfaces rapidement, de répondre à la demande du marché et de renouveler le verger plus souvent.
\item \textbf{Question éthique/environnementale :} \textbf{Uniformité génétique (érosion variétale) et vulnérabilité aux épidémies.} Tous les plants sont des clones du même génotype : si un nouveau pathogène (ex: nouvelle race de \emph{Fusarium} ou virus) apparaît, toute la plantation est sensible simultanément (risque de perte totale). \textbf{Mesure :} Diversifier les sources d'explants (utiliser plusieurs cultivars d'élite différents, introduire de la diversité génétique via graines ou nouveaux méristèmes régulièrement), et ne pas remplacer 100\% du verger par un seul clone in vitro ; maintenir une banque de gènes locale (champ de collection).
\end{enumerate}
\end{corrige}

\begin{corrige}[Étude de document : les bactéries au service de l'assainissement]
\begin{enumerate}
\item \textbf{Problème :} L'\cle{eutrophisation} et l'\cle{anoxie} des cours d'eau récepteurs. Le déversement direct d'eaux usées brutes apporte une charge organique massive (DBO$_5$ élevée). La dégradation naturelle par les bactéries consomme tout l'\ce{O2} dissous $\rightarrow$ asphyxie de la faune aquatique (mortalité piscicole), mise en place de processus anaérobies (fermentation, méthanogenèse) $\rightarrow$ odeurs nauséabondes (\ce{H2S}, \ce{NH3}, mercaptans), dégradation paysagère et risque sanitaire (pathogènes).
\item \textbf{Procédé : \cle{Boues activées} (traitement biologique aérobie).} \textbf{Agent biologique :} Flocons bactériens (boues activées) = consortium de bactéries aérobies (floc-formateurs : \emph{Zoogloea}, \emph{Pseudomonas}, etc.), protozoaires, métazoaires. \textbf{Substrat :} Matière organique carbonée (et azotée/phosphatée) dissoute et colloïdale des eaux usées (mesurée par DBO, DCO). \textbf{Rôle de l'aération :} Apporter le \textbf{dioxygène (\ce{O2})} nécessaire à la respiration aérobie des bactéries (accepteur final d'électrons), maintenir les boues en suspension (contact substrat/biomasse) et éviter la sédimentation/anaérobiose dans le bassin. La réaction globale : \ce{Matiere organique + O2 ->[bacteries] CO2 + H2O + Biomasse (boues) + Energie}.
\item \textbf{Domaine Environnement :} Ce procédé vise explicitement la \textbf{dépollution} (épuration) des eaux usées avant rejet au milieu naturel pour \textbf{protéger les écosystèmes aquatiques} (rivières, nappes) et la santé publique. C'est une application de la \cle{bioremédiation} / bioépuration. La valorisation des boues en compost (amendement organique) s'inscrit aussi dans l'économie circulaire environnementale.
\item \textbf{Enjeu sanitaire/environnemental (Bénéfice) :} \textbf{Protection de la ressource en eau et de la santé des riverains.} Réduction drastique de la charge polluante (DBO, MES, germes pathogènes) $\rightarrow$ cours d'eau vivants, eau potable plus facile à produire en aval, réduction des maladies hydriques (choléra, typhoïde), suppression des nuisances olfactives.
\item \textbf{Enjeu économique/social (Coût/Contrainte) :} \textbf{Investissement lourd et coûts de fonctionnement élevés.} Construction (bassins, soufflantes, bâtiment), consommation électrique massive (aération = 50-60\% de la facture), maintenance technique (pompes, électrovannes, automate), nécessité de \textbf{personnel qualifié} (techniciens traitement eau, biologistes) souvent rare hors grands centres. \textbf{Valorisation des boues :} Nécessite une filière de compostage maîtrisée (hygiénisation, traçabilité métaux lourds) pour être acceptée par l'agriculture ; sinon coût d'enfouissement/incinération. \textbf{Mesure :} Partenariats public-privé (PPP), formation locale (lycées techniques, CFP), biogaz à partir des boues pour réduire la facture électrique.
\end{enumerate}
\end{corrige}

\begin{corrige}[Problème : produire du bioéthanol à partir de la canne à sucre]
\begin{enumerate}
\item \textbf{Procédé :} Fermentation alcoolique (\cle{biotechnologie moderne} : souche sélectionnée, bioréacteur contrôlé, hydrolyse enzymatique/acide du saccharose en amont). \textbf{Agent biologique :} \emph{Saccharomyces cerevisiae} (levure de boulangerie/brasserie, souche industrielle à haut rendement, tolérance éthanol $\approx 15-18\,\%$ vol). \textbf{Substrat :} Glucose (issu de l'hydrolyse du saccharose du vesou de canne).
\item \textbf{Calculs :}
$m_{\text{glucose}} = \SI{360}{kg} = \SI{360\,000}{g}$.
$n_{\text{glucose}} = \dfrac{360\,000}{180} = \SI{2000}{mol}$.
Équation : \ce{C6H12O6 -> 2 C2H5OH + 2 CO2}. $1$ mol glucose $\rightarrow$ $2$ mol éthanol.
$n_{\text{éthanol (th)}} = 2 \times 2000 = \SI{4000}{mol}$.
$m_{\text{éthanol (th)}} = 4000 \times 46 = \SI{184\,000}{g} = \mathbf{\SI{184}{kg}}$.
Rendement $90\,\%$ : $m_{\text{éthanol (réel)}} = 184 \times 0,90 = \mathbf{\SI{165,6}{kg}}$.
\item Volume $V = \dfrac{m}{\rho} = \dfrac{165,6}{0,79} \approx \mathbf{\SI{209,6}{L}}$ (arrondi à $\SI{210}{L}$).
\item \ce{CO2} : $1$ mol glucose $\rightarrow$ $2$ mol \ce{CO2}. $n_{\ce{CO2} (th)} = \SI{4000}{mol}$.
$m_{\ce{CO2} (th)} = 4000 \times 44 = \SI{176\,000}{g} = \SI{176}{kg}$.
Masse réelle (même rendement $90\,\%$ car couplé stœchiométriquement) : $m_{\ce{CO2} (réel)} = 176 \times 0,90 = \mathbf{\SI{158,4}{kg}}$.
\textbf{Valorisation :} Capture et purification pour : industrie des boissons gazeuses (carbonatation), glace carbonique (réfrigération transport vaccins/aliments), serres agricoles (enrichissement en \ce{CO2} pour photosynthèse), synthèse chimique (urée, méthanol).
\item \textbf{Avantage économique :} \textbf{Réduction de la facture pétrolière et création de valeur locale.} Le Cameroun importe 100\% de son carburant raffiné. La filière canne-bioéthanol (SOSUCAM, nouveaux projets) crée des emplois agricoles/industriels non délocalisables, valorise les résidus (bagasse $\rightarrow$ cogénération électricité/vapeur pour l'usine), et améliore la balance commerciale.
\item \textbf{Limite (éthique/sociale/environnementale) :} \textbf{Concurrence alimentaire (Food vs Fuel) et pression foncière.} L'extension des plantations de canne à sucre (monoculture intensive) peut empiéter sur les terres arables vivrières (maïs, manioc, plantain) ou sur les forêts/savanes (déforestation, perte biodiversité), faire monter le prix des denrées de base et déplacer des populations rurales. \textbf{Mesure d'atténuation :} Politique foncière stricte (zonage agricole), priorité aux résidus agricoles (bagasse, paille) et cultures non alimentaires sur terres marginales (jatropha, euphorbe) pour la 2G/3G, certification durable (Bonsucro).
\end{enumerate}
\end{corrige}

\newpage

\coursec{Fiche bilan}

\begin{popfigure}
\begin{center}\tcbox[colback=popblue,colframe=popblue,coltext=white,arc=9pt,boxsep=4pt,left=6pt,right=6pt]{\bfseries\small Biotechnologie}\end{center}
\vspace{-2pt}
\begin{tcbraster}[raster columns=3,raster equal height=rows,raster column skip=6pt,raster row skip=6pt,enhanced,blankest]
\begin{tcolorbox}[enhanced,colback=popgreen,colframe=popgreen,coltext=white,boxrule=0.9pt,arc=3pt,left=4pt,right=4pt,top=3pt,bottom=3pt,boxsep=1pt,halign=left]\bfseries\scriptsize Définition et domaines\end{tcolorbox}
\begin{tcolorbox}[enhanced,colback=poporange,colframe=poporange,coltext=white,boxrule=0.9pt,arc=3pt,left=4pt,right=4pt,top=3pt,bottom=3pt,boxsep=1pt,halign=left]\bfseries\scriptsize Fermentation (traditionnelle)\end{tcolorbox}
\begin{tcolorbox}[enhanced,colback=poppurple,colframe=poppurple,coltext=white,boxrule=0.9pt,arc=3pt,left=4pt,right=4pt,top=3pt,bottom=3pt,boxsep=1pt,halign=left]\bfseries\scriptsize Sélection assistée\end{tcolorbox}
\begin{tcolorbox}[enhanced,colback=poppink,colframe=poppink,coltext=white,boxrule=0.9pt,arc=3pt,left=4pt,right=4pt,top=3pt,bottom=3pt,boxsep=1pt,halign=left]\bfseries\scriptsize Biotechnologies enzymatiques\end{tcolorbox}
\begin{tcolorbox}[enhanced,colback=popgold,colframe=popgold,coltext=white,boxrule=0.9pt,arc=3pt,left=4pt,right=4pt,top=3pt,bottom=3pt,boxsep=1pt,halign=left]\bfseries\scriptsize Enjeux économiques et éthiques\end{tcolorbox}
\end{tcbraster}

\legende{Carte mentale de la leçon : la biotechnologie, des fermentations traditionnelles aux techniques modernes, et ses enjeux.}
\end{popfigure}

\begin{bilan}
\textbf{Définitions clés}
\begin{itemize}
  \item \cle{Biotechnologie} : ensemble des techniques qui utilisent des organismes vivants (micro-organismes, cellules) ou leurs constituants (enzymes) pour produire des biens ou des services.
  \item \cle{Fermentation} : transformation de matières organiques par des micro-organismes (bactéries, levures), en l'absence d'oxygène le plus souvent.
  \item \cle{Sélection assistée} : choix et croisement dirigés d'organismes présentant des caractères d'intérêt, aidés par des outils modernes (marqueurs, analyses).
  \item \cle{Biotechnologie enzymatique} : utilisation d'enzymes purifiées comme catalyseurs dans des procédés industriels.
\end{itemize}

\textbf{Repères à connaître par cœur}
\begin{center}
\begin{tabularx}{\linewidth}{|l|X|X|}
\hline
\rowcolor{popblueL} \textbf{Domaine} & \textbf{Exemple d'application} & \textbf{Type} \\
\hline
Agroalimentaire & bière, vin, yaourt, pain, vin de palme & traditionnelle (fermentation) \\
\hline
Santé & insuline, vaccins, antibiotiques & moderne \\
\hline
Environnement & traitement des eaux usées, dépollution & moderne \\
\hline
Agriculture & variétés améliorées de maïs, manioc & sélection assistée \\
\hline
\end{tabularx}
\end{center}

\textbf{Équations à retenir}
\begin{itemize}
  \item Fermentation alcoolique : \ce{C6H12O6 -> 2C2H5OH + 2CO2}
  \item Respiration cellulaire (à comparer) : \ce{C6H12O6 + 6O2 -> 6CO2 + 6H2O}
\end{itemize}

\textbf{Méthodes express}
\begin{itemize}
  \item Décrire un procédé : matière première $\to$ agent biologique $\to$ conditions $\to$ produit obtenu.
  \item Comparer traditionnel / moderne : critères = outils utilisés, maîtrise du procédé, échelle de production, coût.
  \item Discuter un enjeu : toujours présenter un avantage ET une limite ou un risque (économique, sanitaire, éthique).
\end{itemize}
\end{bilan}

\begin{aretenir}[Checklist avant le Probatoire]
\begin{itemize}
  \item[$\square$] Je sais définir la biotechnologie et citer ses trois grands domaines d'application.
  \item[$\square$] Je sais définir la fermentation et citer les micro-organismes impliqués.
  \item[$\square$] Je connais l'équation de la fermentation alcoolique.
  \item[$\square$] Je sais citer au moins trois produits fermentés courants au Cameroun.
  \item[$\square$] Je sais expliquer le principe de la sélection assistée.
  \item[$\square$] Je sais ce qu'est une biotechnologie enzymatique et donner un exemple.
  \item[$\square$] Je sais comparer une biotechnologie traditionnelle et une moderne selon plusieurs critères.
  \item[$\square$] Je sais décrire les étapes d'un procédé biotechnologique en quatre points.
  \item[$\square$] Je sais citer au moins deux enjeux économiques de la biotechnologie.
  \item[$\square$] Je sais discuter un enjeu éthique (OGM, brevets, sécurité alimentaire) avec des arguments équilibrés.
\end{itemize}
\end{aretenir}

\findecours

\end{document}
